\documentclass[11pt,letterpaper]{article}
\usepackage[T1]{fontenc}
\usepackage[utf8]{inputenc}
\usepackage{lmodern}
\usepackage[margin=1in,headheight=15pt]{geometry}
\usepackage{amsmath,amssymb,amsthm,mathtools,bm}
\usepackage{microtype,booktabs,longtable,array,enumitem}
\usepackage{xcolor,fancyhdr}
\usepackage[colorlinks=true,linkcolor=blue!40!black,citecolor=blue!35!black,urlcolor=blue!40!black]{hyperref}
\usepackage{aliascnt}
\usepackage[nameinlink,noabbrev]{cleveref}
\hypersetup{%
  pdftitle={Euclidean Three-Space over the Surreal Numbers},%
  pdfauthor={Merged research report prepared for Vladimir Reshetnikov},%
  pdfsubject={Coordinate geometry, spherical trigonometry and the rotation group SO(3,No)}}
\pagestyle{fancy}\fancyhf{}
\fancyhead[L]{\small Euclidean three-space over the surreal numbers}
\fancyhead[R]{\small\thepage}
\renewcommand{\headrulewidth}{0.3pt}
\setlength{\parskip}{0.25em}
\setlength{\emergencystretch}{3em}
\setlist{itemsep=0.15em,topsep=0.3em}
\setcounter{tocdepth}{2}
\numberwithin{equation}{section}
\newtheorem{theorem}{Theorem}[section]
\newaliascnt{lemma}{theorem}
\newtheorem{lemma}[lemma]{Lemma}
\aliascntresetthe{lemma}
\crefname{lemma}{lemma}{lemmas}
\newaliascnt{proposition}{theorem}
\newtheorem{proposition}[proposition]{Proposition}
\aliascntresetthe{proposition}
\crefname{proposition}{proposition}{propositions}
\newaliascnt{corollary}{theorem}
\newtheorem{corollary}[corollary]{Corollary}
\aliascntresetthe{corollary}
\crefname{corollary}{corollary}{corollaries}
\theoremstyle{definition}
\newaliascnt{definition}{theorem}
\newtheorem{definition}[definition]{Definition}
\aliascntresetthe{definition}
\crefname{definition}{definition}{definitions}
\newaliascnt{example}{theorem}
\newtheorem{example}[example]{Example}
\aliascntresetthe{example}
\crefname{example}{example}{examples}
\theoremstyle{remark}
\newaliascnt{remark}{theorem}
\newtheorem{remark}[remark]{Remark}
\aliascntresetthe{remark}
\crefname{remark}{remark}{remarks}
\crefname{part}{part}{parts}
\newcommand{\No}{\mathbf{No}}
\newcommand{\R}{\mathbb R}
\newcommand{\C}{\mathbb C}
\newcommand{\Q}{\mathbb Q}
\newcommand{\Z}{\mathbb Z}
\newcommand{\N}{\mathbb N}
\newcommand{\HH}{\mathbb H}
\newcommand{\TT}{\mathbb T}
\newcommand{\OO}{\mathcal O}
\newcommand{\mm}{\mathfrak m}
\newcommand{\Oz}{\mathbf{Oz}}
\newcommand{\eps}{\varepsilon}
\newcommand{\st}{\operatorname{st}}
\newcommand{\supp}{\operatorname{supp}}
\newcommand{\cis}{\operatorname{cis}}
\newcommand{\atan}{\operatorname{atan}}
\newcommand{\atantwo}{\operatorname{atan2}}
\newcommand{\acos}{\operatorname{acos}}
\newcommand{\asin}{\operatorname{asin}}
\newcommand{\hav}{\operatorname{hav}}
\newcommand{\SO}{\operatorname{SO}}
\newcommand{\Orth}{\operatorname{O}}
\newcommand{\Sp}{\operatorname{Sp}}
\newcommand{\Spin}{\operatorname{Spin}}
\newcommand{\SU}{\operatorname{SU}}
\newcommand{\GL}{\operatorname{GL}}
\newcommand{\so}{\mathfrak{so}}
\newcommand{\Ad}{\operatorname{Ad}}
\newcommand{\id}{\mathrm{id}}
\newcommand{\im}{\operatorname{Im}}
\newcommand{\re}{\operatorname{Re}}
\newcommand{\cay}{\operatorname{Cay}}
\newcommand{\Sin}{\operatorname{Sin}}
\newcommand{\Cos}{\operatorname{Cos}}
\newcommand{\Exp}{\operatorname{Exp}}
\newcommand{\Log}{\operatorname{Log}}
\newcommand{\Rot}{\operatorname{Rot}}
\newcommand{\Aut}{\operatorname{Aut}}
\newcommand{\area}{\operatorname{Area}}
\newcommand{\vol}{\operatorname{Vol}}
\newcommand{\dist}{\operatorname{dist}}
\newcommand{\diag}{\operatorname{diag}}
\newcommand{\tr}{\operatorname{tr}}
\newcommand{\sgn}{\operatorname{sgn}}
\newcommand{\Span}{\operatorname{span}}
\newcommand{\norm}[1]{\left\lVert#1\right\rVert}
\newcommand{\abs}[1]{\left|#1\right|}
\newcommand{\ip}[2]{\left\langle#1,#2\right\rangle}
\newcommand{\vv}[1]{\bm{#1}}
\newcommand{\transpose}{^{\mathsf T}}
\newcommand{\ang}{\operatorname{ang}}
\newcommand{\smallO}{\operatorname{O}}
\newcommand{\comm}[2]{\left(#1,#2\right)}
\newcolumntype{L}[1]{>{\raggedright\arraybackslash}p{#1}}
\newcommand{\repo}{\texttt{VladimirReshetnikov/Surreal}}

\begin{document}
\hypersetup{pageanchor=false}
\begin{titlepage}
\centering
\vspace*{0.2cm}
{\large\scshape Merged research report\par}
\vspace{0.5cm}
{\LARGE\bfseries Euclidean Three-Space over the Surreal Numbers\par}
\vspace{0.4cm}
{\large Coordinate geometry, spherical trigonometry,\\[0.15em]
and the rotation group $\SO(3,\No)$\par}
\vspace{0.4cm}
{\large 22 September 2026\par}
\vspace{0.5cm}
\begin{minipage}{0.94\textwidth}
\begin{abstract}
\noindent
We develop Euclidean geometry in $\No^3$ in two parts. Every construction is
carried out in a set-sized real closed subfield of $\No$, or in a set-sized
real-closed Hahn field when analytic evaluation is needed. Distances, areas,
volumes and sphere radii may be infinite or infinitesimal. Angles between
directions always have canonical finite surreal representatives, so spatial
trigonometry needs no oscillatory value at an infinite argument.

Part~I is coordinate geometry. It covers affine incidence, projections, skew
lines, sphere intersections, quadrics and rigid motions. On spheres it proves
the sine and cosine laws, polar duality, and sharp existence and congruence
criteria, including an exceptional continuous SSA family. It also gives
spherical centers, Ceva, a branch-correct solid-angle formula, and a finitely
additive area for geodesic polygons. Parallel transport has quaternionic
holonomy equal to the spherical excess, and a relative planar-limit area
estimate holds for arbitrarily thin infinitesimal triangles.

Part~II studies the class group $G=\SO(3,\No)$. Rotation-matrix entries are
bounded, and standard part gives a split extension $G=K\rtimes\SO(3,\R)$. The
infinitesimal kernel $K$ is uniquely divisible, and each of its elements is a
single commutator of elements of $K$. Every finite subgroup is conjugate to
its standard part by an infinitesimal rotation. We treat strong exponential
and logarithm maps and the valuation-graded cross-product algebra. For the
global rotation exponential we give its fibers and differential; its critical
radii are $2\pi\Oz$. Finally, three topologies are separated.

The rotation algebra shared by the two parts is printed once in Part~I: axes,
Rodrigues, quaternions, and the Cayley chart. The quaternion double cover is
cited from the collection's surquaternion report.
\end{abstract}
\end{minipage}
\vfill
\begin{minipage}{0.94\textwidth}\small
\textbf{Status and provenance.} An AI-assisted research report merged from two
manuscripts of 22 September 2026 (sources~04 and~09 of the batch that placed
them); see \cref{e3:sub:provenance}. It has not been refereed, and no statement
in it has a Lean formalization. Classical identities, the Bays--Peterzil
perfectness precedent and imported analytic foundations are credited, not
claimed. No priority claim or resolution of a named open problem is made. The
accompanying symbolic checks test selected finite identities only.
Repository pin of this merge: \texttt{50cb709}.
\end{minipage}
\end{titlepage}
\hypersetup{pageanchor=true}
\tableofcontents
\clearpage

\section{Scope, sources, and conventions}\label{e3:sec:scope}

An infinite coordinate is not a point at infinity, and an infinitesimal distance
is not zero. Those two distinctions are enough to change many geometric
interpretations, but they do not destroy Euclidean coordinate algebra. The
appropriate starting point is an ordered field equipped with positive square
roots, rather than the completeness of the real line.

The central principle of Part~I is
\[
 \boxed{\text{Coordinates and lengths can have arbitrary scale; geometric
 angles can always be finite.}}
\]
For example, the vectors $(\omega,1,0)$ and $(\omega,0,0)$ make the nonzero
angle $\atan(\omega^{-1})$. The large coordinate does not require evaluation of
$\sin\omega$. Likewise, a great-circle arc on a sphere of radius $R=\omega$
can have an infinite length $R a$ while its central angle $a$ lies in $[0,\pi]$.
The dimensionless quantity entering trigonometric functions is $a$, not $Ra$.

Part~II asks what the group of rotations of this space is. Let
\(G(F)=\SO(3,F)=\{A\in M_3(F):A\transpose A=I,\ \det A=1\}\) for a real closed
field $F$; the group of interest is $G(\No)$, read with the size conventions of
\cref{e3:sec:scalars}. There are three complementary answers:
\begin{align}
G(\No)&\cong \Sp(1,\No)/\{1,-1\},\label{e3:so:eq:overviewspin}\\
G(\No)&\cong K\rtimes\SO(3,\R),\label{e3:so:eq:overviewsplit}\\
K&\cong\bigl(\mm^3,\star\bigr),\qquad
 \vv t\star\vv u=\frac{\vv t+\vv u+\vv t\times\vv u}{1-\vv t\cdot\vv u}.\label{e3:so:eq:overviewcay}
\end{align}
Here $\Sp(1,\No)$ is the unit group of the Hamilton quaternion algebra,
$\mm$ is the class of infinitesimal surreals, and $K$ consists of the rotation
matrices infinitesimally close to the identity. The denominator in
\cref{e3:so:eq:overviewcay} is a unit infinitesimally close to one, so the
formula is valid on all of $\mm^3$. It is a noncommutative group law, not
vector addition. The first isomorphism is the algebraic double cover of the
collection's surquaternion report; the second and third are
\cref{e3:so:thm:split,e3:so:prop:kernelchart}.

The geometry of a rotation still has three scalar degrees of freedom. What
changes is the range of those scalars and the meaning of analysis. Infinitesimal
changes of axis and angle give genuinely different rotations even when their
standard parts agree. A small rotation can move an infinitely distant point by
a finite or infinite distance. Ordinary real paths, compactness arguments, and
power-series convergence cannot simply be imported into the full fine topology.

\subsection{Two manuscripts, one report}\label{e3:sub:provenance}

This report is merged from two manuscripts, both dated 22 September 2026 and
written independently against different commits of the repository \repo.
The source numbers are those of the batch in which they arrived.

{\small
\begin{longtable}{L{0.07\textwidth}L{0.32\textwidth}L{0.09\textwidth}L{0.38\textwidth}}
\toprule
Source & Manuscript & Pin & Contribution here\\
\midrule
\endhead
04 (base) & \emph{Analytic Geometry and Trigonometry in Surreal Three-Space:
 finite angles, arbitrary scales, spherical duality, solid angles, rotations,
 and curvature} (34 pages) & \texttt{d22a5b3} &
 \Crefrange{e3:sec:scalars}{e3:sec:boundaries}, that is, \cref{e3:sec:scalars}
 and all of Part~I; \cref{e3:sub:impl04,e3:sub:lean04}; the first proof of
 \cref{e3:thm:phase}. Files prefixed \texttt{04-three-space-}.\\
09 & \emph{Rotations of Surreal Three-Space: Algebra, Spin, Infinitesimal
 Structure, and Topology} (31 pages, 19 sections, 26 numbered statements) &
 \texttt{dcf8666} &
 All of Part~II
 (\crefrange{e3:so:sec:frames}{e3:so:sec:examples});
 in Part~I the localization proposition and three-part decomposition of
 \cref{e3:sec:scalars}, the hat identities of \cref{e3:sec:vectors}, and the
 algebraic-angle, spin-lift, bounds and Cayley-chart statements of
 \cref{e3:sec:rotations}; \cref{e3:sub:impl09,e3:sub:lean09}. Files prefixed
 \texttt{09-so3-}.\\
\bottomrule
\end{longtable}}

\textbf{Why one report.} The two manuscripts overlap in the finite algebra of
rotations and are otherwise complementary. Source~04 is a geometry of points,
lines, spheres and spherical triangles, for which rotations are one tool;
source~09 is a structure theory of the rotation group. Source~04 is the base: it
is the wider development and fixes the scalar conventions. Source~09 forms
Part~II, after the geometry.

\textbf{Printed once.} The following results are proved in both sources and
appear once here. We say which source's statement and proof are printed.
\begin{enumerate}[label=(\alph*)]
\item The finite phase theorem, \cref{e3:thm:phase}: source~04's statement
 and proof, with source~09's logarithmic proof kept as a marked second route.
\item The classification of distance-preserving maps,
 \cref{e3:thm:isometries}. We print source~04's affine statement, add
 source~09's orientation clause and its remark on field automorphisms, and
 give one proof; the two proofs were the same polarization argument.
\item The axis theorem and Rodrigues formula, \cref{e3:thm:rodrigues}. The
 statement is source~09's algebraic-angle form, which contains source~04's
 finite-angle form \cref{e3:eq:rodrigues}. Both proofs were the same
 determinant argument.
\item Finite axis--angle representatives, \cref{e3:thm:finiteaxis}. This
 combines a paragraph of source~04 with the second half of source~09's
 theorem on finite angles, whose proof is printed.
\item The Hamilton product and the quaternion rotation formula,
 \cref{e3:eq:quaternionmul,e3:eq:quatrotation}. The two sources give these
 identically.
\item The two-to-one spin map $\Sp(1,F)\to\SO(3,F)$ with kernel $\{\pm1\}$,
 \cref{e3:thm:spin}. Both sources reprove it by the argument that the
 surquaternion report already gives as the first route of
 \texttt{squat:thm:rotations}. We cite that theorem and do not repeat the
 proof; only source~09's explicit lift formula is printed.
\item The rational Cayley chart, \cref{e3:so:thm:cayley}. Source~04's vector
 formula \cref{e3:eq:cayley} and source~09's matrix formula are the same map,
 because $\vv t\times(\vv t\times\vv x)=\vv t(\vv t\cdot\vv x)-\norm{\vv t}^2\vv x$.
 The theorem is source~09's; source~04's remark on infinite chart
 coordinates is kept.
\item Scalar foundations. The finite and infinitesimal elements, the standard
 part, the Hahn valuation, strong summability and the commutative
 positive-support lemma appear once, in \cref{e3:sec:scalars}. The
 noncommutative all-lengths version needed for BCH is
 \cref{e3:so:lem:neumann}.
\item The symmetric spectral theorem over a real closed field. It is proved in
 \cref{e3:sub:quadrics} (source~04); source~09's proof of
 \cref{e3:so:lem:positiveroot} used the same eigenvector argument and now
 refers to it.
\end{enumerate}

\textbf{Where the merge had to choose.} Symbols were renamed only where the
two sources, or a source and the rest of the collection, used one symbol for
two things. \Cref{e3:sub:notation} lists each renaming. The largest is that
source~04's scalar workspace $K$ is written $F$ here, because $K$ is source~09's
infinitesimal kernel. The word \emph{fine} also changes meaning: source~04 used
it for the order topology of one set-sized workspace, source~09 for the full
surreal topology. Here \emph{fine} always has the second meaning, in agreement
with the surquaternion and foundations reports. Source~04's usage is written
\emph{intrinsic}. Labels from source~04 carry the prefix \texttt{e3:}, and
labels from source~09 carry \texttt{e3:so:}. Text written for the merge is confined to this section, the framing of
\cref{e3:sec:rotations}, \cref{e3:sec:nonclaims}, the ledgers, and short
cross-reference remarks. It adds no mathematical claim
(\cref{e3:sub:nonclaimsmerge}).

\subsection{Relation to the collection}\label{e3:sub:collection}

\textbf{Trigonometry.} The trigonometry report \cite{RepoTrig}, in
\path{docs/surcomplex/trigonometry}, develops canonical finite phases,
arbitrary-scale plane geometry and the spherical sine and cosine laws. \Cref{e3:thm:phase} is the Hahn-workspace form of its
\texttt{trigonometry:thm:polar}. The choice of representatives is its
\texttt{trigonometry:cor:representatives}, and the planar half-angle chart is its
\texttt{trigonometry:thm:cayley}. The spherical laws of
\cref{e3:thm:spherical-laws} rederive its \texttt{trigonometry:thm:spherical}
from Gram data, adding polar duality, reconstruction and area. The
Ehrlich--Kaplan phase of \cref{e3:so:sec:angles} is its
\texttt{trigonometry:def:globaltrig} and \texttt{trigonometry:thm:globalexp}. A
treatment of $\SO(2,\No)$ is being added to that report at the same time as this
merge. This report cites only labels present at its pin \texttt{50cb709}.

\textbf{Surquaternions.} The report \emph{Surquaternions}
(\path{docs/surquaternions/surquaternions}) \cite{RepoQuaternion} already
contains the following, cited by label with its theorem numbers at pin
\texttt{50cb709}:
\begin{itemize}
\item the rotation double cover, and the isomorphism of $\SO(3,F)$ with
 $\Aut_F(\HH_F)$ (\texttt{squat:sub:spin}; \texttt{squat:thm:rotations},
 Theorem~4.4);
\item the quaternion Cayley and stereographic chart (\texttt{squat:prop:charts},
 Proposition~4.5);
\item the standard-part splitting of the unit group
 (\texttt{squat:sub:unitreduction}, \S4.3);
\item infinitesimal exponential and logarithm, and the filtered group of
 infinitesimal rotations (\texttt{squat:thm:principal-log},
 \texttt{squat:thm:rotation-log}, \texttt{squat:thm:graded}; Theorems~8.1,
 8.5, 8.6);
\item the all-lengths positive-support lemma (\texttt{squat:lem:neumann},
 Lemma~7.1) and set discreteness of the fine topology
 (\texttt{squat:thm:discrete}, Theorem~7.3);
\item the differential of the four-dimensional radial exponential
 (\texttt{squat:thm:exp-derivative}, Theorem~11.2).
\end{itemize}
Part~II works instead with $3\times3$ rotation matrices and matrix Cayley
coordinates. The two descriptions differ by the factor two of the spin
differential (\cref{e3:so:prop:liesimple}). In particular the rank-one critical
spheres of the rotation exponential in \cref{e3:so:thm:expdifferential} are not
the rank-two critical spheres of \texttt{squat:thm:exp-derivative}. At its pin
\texttt{dcf86662b574}, source~09 records these surquaternion results as starting
points, not as its own discoveries. At its pin \texttt{d22a5b35d5b3}, source~04
inspected only the trigonometry report and one Lean file. It does not mention
the surquaternion report, although that report already contained the double cover and
the quaternion Cayley chart at that commit. The citations in
\cref{e3:sec:rotations} correct this omission.

\textbf{Foundations.} The foundations report
(\path{docs/foundations-and-computation/foundations}) \cite{RepoFoundations}
proves set discreteness of the fine topology (\texttt{found:thm:discrete}) and
separates intrinsic from full-class subspace topologies
(\texttt{found:prop:twotopologies}). \Cref{e3:so:thm:finediscrete} specializes the
first to rotation coordinates.

\textbf{Lean.} \texttt{Surreal/Algebra/Geometry.lean} develops
two-dimensional dot products, determinants, Cauchy--Schwarz, the Ptolemy
inequality and the area form of Heron's identity over ordered fields with
nonnegative square roots, for labels of the trigonometry report
\cite{RepoGeometry}. It explicitly distinguishes these field-valued lengths from
the real-valued distances used by conventional metric-space interfaces. The repository contains no three-dimensional Euclidean or
rotation-group Lean development, and no label of this report is mapped to a
Lean declaration.

\textbf{Vector and tensor fields.} The report
\path{docs/surreal/vector-and-tensor-fields}, placed in the same commit as
this one, develops finite tensor algebra, metrics and Levi-Civita connections
over set-sized real closed surreal fields. The explicit sphere connection,
curvature and transport of \cref{e3:sec:differential} are computed directly
here and are not derived from that framework.

\subsection{Three levels of assumptions}\label{e3:sub:levels}

\begin{center}\small
\begin{tabular}{L{0.18\textwidth}L{0.34\textwidth}L{0.38\textwidth}}
\toprule
Level & Scalar assumptions & What the assumptions support\\
\midrule
Algebraic & Ordered field with positive square roots; real closed when roots of
 general polynomials are used & Incidence, norms, cross products, rigid motions,
 Gram matrices, algebraic sine--cosine data, the finite algebra of $\SO(3,F)$\\
Standard part & A real closed field $F\supseteq\R$ & The split extension,
 the infinitesimal kernel, finite subgroups\\
Trigonometric & A real-closed Hahn workspace with its canonical restricted
 analytic structure & Finite angle representatives, inverse trigonometric
 functions, triangle identities, bounded analytic formulas, strong
 exponential and logarithm\\
Differential & Explicit restricted-analytic charts and differentiation in an
 external parameter, or the fine Fr\'echet derivative & Metric tensors, sphere
 curvature, finite-angle geodesics, parallel transport, the differential of
 the rotation exponential\\
\bottomrule
\end{tabular}
\end{center}

The trigonometric and differential levels do not imply an unrestricted
integration theory. Our spherical polygon area is explicitly constructed and
proved finitely additive. The distinction matters even for apparently harmless
real arguments; \cref{e3:sec:boundaries} gives concrete counterexamples to naive
transfers of calculus.

Source~09 states its hypotheses in the same spirit. The proofs in
\cref{e3:sec:rotations,e3:so:sec:frames,e3:so:sec:conjugacy} are finite algebra over
a real closed field. Standard-part results require $F\supseteq\R$. Strong sums
are constructed either in the full surreal normal form or in a named divisible
Hahn workspace. Global infinite angles use the specific Ehrlich--Kaplan
normalization. Semialgebraic homotopy statements are made over set-sized real
closed fields. These are different hypotheses; none is silently substituted for
another.

\subsection{Notation, and the renamings made by the merge}\label{e3:sub:notation}

$F$ denotes a real closed scalar field. In Part~I it is a set-sized workspace
inside $\No$, chosen in \cref{e3:sec:scalars}; in Part~II it is any real
closed field unless $F\supseteq\R$ is stated. Bold letters denote vectors in
$F^3$ in Part~I; Part~II, following source~09, writes vectors in plain type.
Points can be subtracted to give vectors, and an origin is chosen whenever
coordinates are used. The physical radius of a sphere is $R>0$. For a
spherical triangle with unit direction vectors $\vv u,\vv v,\vv w$, the
opposite \emph{angular} side lengths are $a,b,c$, and the interior angles are
$\alpha,\beta,\gamma$. Physical side lengths are $Ra,Rb,Rc$. We consistently use
\begin{equation}\label{e3:eq:xyzconvention}
 x=\vv v\cdot\vv w=\cos a,\qquad
 y=\vv w\cdot\vv u=\cos b,\qquad
 z=\vv u\cdot\vv v=\cos c.
\end{equation}
The determinant $\tau=\det[\vv u\ \vv v\ \vv w]$ is oriented;
$D=\tau^2$ is the Gram determinant; $d=\sqrt D=|\tau|$ is nonnegative.
These symbols will not also denote area or a side length. In
\cref{e3:so:sec:frames} the letters $a,b,c,d$ are instead the coordinates of a
quaternion $a+bi+cj+dk$.

The merge fixed one meaning per symbol as follows.
\begin{center}\small
\begin{tabular}{L{0.2\textwidth}L{0.34\textwidth}L{0.36\textwidth}}
\toprule
Here & Source notation & Reason\\
\midrule
$F$, $F^3$, $\SO(3,F)$ & source~04: $K$, $K^3$, $\SO(3,K)$ & $K$ is the infinitesimal kernel of Part~II\\
$\TT(F)$ & source~04: $U(K)$, the unit circle & $U$ denotes elements of $K$; $\TT$ as in the trigonometry report\\
$\Rot(n;c,s)$ & source~09: $R(n;c,s)$; source~04: $Q_{\vv n,\theta}$ & $R$ is a sphere radius\\
$A$ for rotation matrices in \cref{e3:sec:rotations} & source~04: $Q$ & $Q$ is also a point\\
$\Phi$ for an isometry & source~04: $F$ & $F$ is the field\\
$A_0=\st(A)$, $B,B'$ real rotations & source~09: $R_0$, $R$, $S$ & $R$ is a radius; $S_u$ a reflection\\
$\mathcal E$, $\mathcal J(w)$ & source~09: $E$, $J(w)$ for the global rotation exponential and its differential & $E$ is the spherical excess, $J$ a Jacobian\\
dual number $\kappa$, $\kappa^2=0$ & source~09: $\tau$ & $\tau$ is an oriented determinant\\
$\mathrm{Gn}$, $\mathrm{St}$ & source~04: $G$, $H$ for the gnomonic and stereographic charts & $G(F)=\SO(3,F)$; $H_n$ is a half-turn\\
$\sin,\cos,\cis$ on finite angles & source~09: $\sin_0,\cos_0,\cis_0$ & one finite-angle convention; $\Sin,\Cos,\cis_{\mathrm{EK}}$ remain the global ones\\
$\widehat x$ & source~04: $[x]_\times$ & one name for the cross-product matrix\\
$f=\smallO(g)$ & source~04: $F=\smallO(G)$ & letters in use\\
\emph{intrinsic} topology of a workspace & source~04: \emph{fine} & \emph{fine} means the full surreal topology\\
\bottomrule
\end{tabular}
\end{center}

The local matrix $E=P-I$ in source~09's positive-root lemma is written $Y$.
Two further symbols keep a double use from their sources. The letter $t$ is the
Hahn monomial variable in $t^\gamma$; the vector $\vv t$, or $t\in F^3$ in
Part~II, is a Cayley coordinate. $\Exp_P,\Log_P$ with a point subscript are the
exponential map of a sphere and its inverse, while $\Exp_{\mm},\Log_{\mm}$ are
the strong matrix exponential and logarithm.

\section{Scalar foundations and canonical finite trigonometry}\label{e3:sec:scalars}

\subsection{Set-sized workspaces inside the surreal numbers}

The surreal numbers form a real-closed ordered \emph{class} field, not a set.
We use their Conway normal forms but perform each construction in a set-sized
field. Let $\Gamma\subset\No$ be a set-sized divisible ordered additive
subgroup and put
\begin{equation}\label{e3:eq:hahn}
 F=\R((t^\Gamma)),\qquad
 \sum_{\gamma} a_\gamma t^\gamma\longmapsto
 \sum_{\gamma} a_\gamma\omega^{-\gamma}\in\No.
\end{equation}
Supports in the left-hand expression are well ordered. This identifies $F$
with a real-closed subfield of $\No$. Its canonical restricted-analytic
structure is an elementary extension of $\R_{\mathrm{an}}$; these structures
are compatible under enlargement of $\Gamma$ \cite[\S2]{vdDE}.

Every finite configuration of surreal coordinates belongs to such a
workspace: collect all exponents occurring in its finitely many normal forms
and take their rational span. More generally the same argument works for any
set of inputs. Neither the support set nor this span is a proper class.
All operations used below stay in a workspace or in a specified larger one.
Consequently the theorems apply to arbitrary individual surreal coordinates
without treating $\No$ as a set or forming a proper-class power set.

Source~09 makes the same point for the rotation group. There is no set
containing all surreal numbers, so $\No^3$, $M_3(\No)$ and $G(\No)$ are proper
classes. One may use a class theory such as NBG, or a universe-relative
formulation in which the surreal carrier is large relative to the universe
containing its individual sign sequences. A matrix still has only nine entries,
and every formula for its multiplication, determinant, inverse, or action is
finite. Statements about all elements of this class group are schemes of
ordinary finite algebra. Statements about a finite subgroup, a sequence, or a
set-indexed family retain their ordinary size restriction. We do not form an
unrestricted set of all subsets of $\No$, or invoke ordinary compactness for
such a nonexistent set-sized space. For finite algebra alone a smaller field
suffices.

\begin{proposition}[Localization of finite algebra]\label{e3:so:prop:localization}
Every set-sized family of surreal coordinates lies in a set-sized real closed
subfield $F\subseteq\No$. One can require $\R\subseteq F$. Thus every
set-sized configuration of vectors and matrices can be studied inside an honest
real closed field without changing its finite algebraic identities.
\end{proposition}
\begin{proof}
Adjoin the coordinates to $\R$, obtaining a set-sized ordered subfield.
Take its real closure inside the real closed field $\No$: equivalently, take
the elements of $\No$ algebraic over that subfield. They form a set and a
real closed field. The inclusion preserves the arithmetic and order. It does
not assert that the subfield is closed under arbitrary strong sums or inherits
the same topology as its intrinsic order topology.
\end{proof}

The Hahn workspace \cref{e3:eq:hahn} is larger than this subfield; its purpose
is to contain strong sums and analytic evaluations. An arbitrary birthday
cutoff is not a substitute for either construction. For example, the
finite-birthday surreals are dyadic rationals and are not closed under
division. A workspace must have the closure properties actually needed.

We write
\[
 \OO=\{a\in F: |a|<n\text{ for some }n\in\N\},\qquad
 \mm=\{a\in F: |a|<1/n\text{ for every }n\ge1\}.
\]
These are respectively the finite elements and the infinitesimals. Every
$a\in\OO$ has a unique decomposition
\[
 a=\st(a)+\eps,\qquad \st(a)\in\R,\quad\eps\in\mm.
\]
The standard-part map is an ordered ring homomorphism on $\OO$, with kernel
$\mm$. In a Hahn workspace it is the coefficient of $t^0$ for finite series.
The same definitions apply to any ordered field $F\supseteq\R$, with
subscripts $\OO_F,\mm_F$ where the field matters; $\OO$ and $\mm$ without
subscript refer to the workspace, or to all of $\No$. In that generality the
real cut determined by a bounded $a$ has a real supremum, whose difference from
$a$ is infinitesimal. Uniqueness follows because no nonzero real is
infinitesimal (source~09).

A surreal has a set-supported normal form
$x=\sum_{a\in S} r_a\omega^a$ with $r_a\in\R\setminus\{0\}$ and $S$ reverse
well ordered. Splitting the exponents into positive, zero, and negative parts
gives
\begin{equation}\label{e3:so:eq:threeparts}
 x=x^\uparrow+r+\eps .
\end{equation}
The first part is purely infinite, the second is real, and the third is
infinitesimal. Purely infinite does not mean positive: a purely infinite number
may have either sign. Normal-form and analytic foundations are treated in
\cite{vdDE,EK}; only set-supported sums are used.

For $a\ne0$, let $v(a)$ be the least exponent in its Hahn support, and set
$v(0)=+\infty$. In normal-form terms, $v(x)=-\max\supp_\omega(x)$, the
negative of the leading exponent, so $v(\omega^{-\gamma})=\gamma$; the two
sources use this one convention, which is also the collection's. For a vector
or matrix, $v$ is the minimum of the valuations of its entries. Thus
$v(t^\gamma)=\gamma$, positive valuation means
infinitesimal magnitude, and $v(ab)=v(a)+v(b)$. If $a,b\ne0$, the notation
$a\sim b$ means $a/b=1+\eta$ with $\eta\in\mm$. Later,
$f=\smallO(g)$ means $|f|\le C|g|$ for a positive \emph{real} constant $C$
in the stated range; this is not an assertion about convergence of a sequence
of ordinary partial sums.

\subsection{What an infinite Taylor sum means}

A family of Hahn series is \emph{strongly summable} when the union of its
supports is well ordered and every exponent occurs in only finitely many
members. The sum is then defined coefficientwise. The positive-support lemma
of Hahn-series algebra says that if $S$ is a well-ordered subset of
$\Gamma_{>0}$, its additive monoid is well ordered and every exponent has only
finitely many decompositions as an unordered finite sum from $S$.
This is the summability mechanism used below; one must not replace it by a
claim that $n\delta$ is cofinal in $\Gamma$ for every $\delta>0$.
The latter claim is false in higher-rank ordered groups. See
\cite{Neumann,vdDEErr} for the relevant support theory. The all-lengths
version for ordered finite words, needed for noncommutative matrix series, is
\cref{e3:so:lem:neumann}, which is \texttt{squat:lem:neumann} of the
surquaternion report.

If $f$ is real analytic near a real number $r$ and $h\in\mm$, define
\begin{equation}\label{e3:eq:taylor}
 f^F(r+h)=\sum_{n=0}^{\infty}\frac{f^{(n)}(r)}{n!}h^n.
\end{equation}
The positive-support lemma proves strong summability. It also justifies
multiplication, substitution, and differentiation of these expressions when
used in the corresponding analytic germs. This is consistent with the
restricted-analytic structure in \cref{e3:eq:hahn}.

For a nonzero ordinary real argument, the usual Taylor series about zero
need not be a strongly summable Hahn family: infinitely many nonzero real
terms have exponent zero. One first evaluates the real analytic function at
its real center, and only then evaluates its infinitesimal Taylor expansion.
This convention is essential, not a stylistic choice.

\subsection{Sine, cosine, and inverse tangent}

For $a=r+h\in\OO$, define $\sin a$ and $\cos a$ by
\cref{e3:eq:taylor}. In particular,
\begin{align}
 \sin h&=h-\frac{h^3}{3!}+\frac{h^5}{5!}-\cdots,\label{e3:eq:sininf}\\
 \cos h&=1-\frac{h^2}{2!}+\frac{h^4}{4!}-\cdots.\label{e3:eq:cosinf}
\end{align}
The formal addition identities, together with the ordinary real identities,
give
\begin{equation}\label{e3:eq:trigaddition}
 \sin(a+b)=\sin a\cos b+\cos a\sin b,\qquad
 \cos(a+b)=\cos a\cos b-\sin a\sin b,
\end{equation}
and $\cos^2a+\sin^2a=1$. In particular $\sin h=h(1+\smallO(h^2))$.
The real sign and monotonicity properties on any fixed real bounded interval
hold with surreal inputs as well. They follow either from the local
expansions and endpoint checks, or from restricted-analytic elementarity.

Define $\atan a$ by real-centered Taylor evaluation for finite $a$, and for
infinite $a$ by
\begin{equation}\label{e3:eq:ataninfinite}
 \atan a=\sgn(a)\frac\pi2-\atan(1/a).
\end{equation}
It is the unique angle in $(-\pi/2,\pi/2)$ whose tangent is $a$.
Thus inverse tangent accepts arbitrary surreal inputs and always returns a
finite surreal angle. Tangent itself is used only where its cosine denominator
is nonzero.

\begin{theorem}[Every direction has a finite phase]\label{e3:thm:phase}
Let
$\TT(F)=\{(c,s)\in F^2:c^2+s^2=1\}$, with multiplication
$(c,s)(c',s')=(cc'-ss',sc'+cs')$. The map
\[
 \theta\longmapsto(\cos\theta,\sin\theta)
\]
induces an isomorphism $\OO/(2\pi\Z)\simeq \TT(F)$.
Consequently every unit-circle direction has a unique representative in
$[0,2\pi)$, and also a unique representative in $(-\pi,\pi]$.
\end{theorem}
\begin{proof}
Both coordinates of a unit-circle point are finite. Its standard part is an
ordinary real unit-circle point. Rotate this standard part to $(1,0)$ using
an ordinary real phase. The resulting point $(C,S)$ has $C=1+\mm$, $S\in\mm$,
and $C>0$. Put $h=\atan(S/C)\in\mm$. Since the unit point with positive first
coordinate and slope $S/C$ is unique, $(C,S)=(\cos h,\sin h)$. This proves
surjectivity. Addition follows from \cref{e3:eq:trigaddition}.

For the kernel, the real standard part of a phase must be $2\pi k$ for an
ordinary integer $k$. Subtract it. An infinitesimal $h$ with $\sin h=0$
must vanish because $\sin h=h(1+\smallO(h^2))$ and the second factor is a
unit. Thus the kernel is exactly $2\pi\Z$, not an unspecified group of
infinite periods. Finally a finite element can be reduced into either stated
interval by an ordinary integer multiple of $2\pi$.
\end{proof}

\begin{proof}[Second proof (source~09)]
Write $\cis\theta=\cos\theta+i\sin\theta$ in $F(i)$. For $\theta=r+\eps$ with
$r\in\R$ and $\eps\in\mm$, real-centered evaluation gives
$\cis\theta=e^{ir}\sum_{k\ge0}(i\eps)^k/k!$; this is \cref{e3:so:eq:finitephase}.
For a unit surcomplex $z=c+is$, both coordinates are finite. Its standard
part $z_0$ is an ordinary unit complex number. Choose a real angle $r$ with
$e^{ir}=z_0$. Then $z/z_0=1+\delta$ for infinitesimal $\delta$, so its
strong logarithm is defined. Conjugation and inversion give
\[
 \overline{\log(z/z_0)}=\log(\overline{z/z_0})
 =\log((z/z_0)^{-1})=-\log(z/z_0).
\]
It is therefore $i\eps$ with real infinitesimal $\eps$, and $z=\cis(r+\eps)$.
This proves surjectivity. If $\cis(r+\eps)=1$, standard part gives
$r\in2\pi\Z$; then logarithmic injectivity gives $\eps=0$.
\end{proof}

The second proof is the argument of \texttt{trigonometry:thm:polar}, which proves
the same isomorphism over all of $\No$ \cite{RepoTrig}. The choice of
representatives, including the endpoint correction for infinitesimal
excursions, is \texttt{trigonometry:cor:representatives}.

For $(X,Y)\ne(0,0)$, define $\atantwo(Y,X)\in(-\pi,\pi]$ as the phase
of $(X,Y)/\sqrt{X^2+Y^2}$. With $r=\sqrt{X^2+Y^2}$, an explicit formula is
\begin{equation}\label{e3:eq:atan2}
 \atantwo(Y,X)=2\atan\frac{Y}{r+X}
 \quad\text{when }r+X>0,
\end{equation}
with value $\pi$ on the negative $X$-axis. This definition works when either
coordinate is infinite or infinitesimal.

For $-1\le x\le1$, let
\begin{equation}\label{e3:eq:acos}
 \acos x=\atantwo\bigl(\sqrt{1-x^2},x\bigr)\in[0,\pi],\qquad
 \asin x=\frac\pi2-\acos x.
\end{equation}
Here the phase at $(1,0)$ is zero. These inverse functions include their
endpoints. For $-1<x\le1$, another exact formula is
$\acos x=2\atan\sqrt{(1-x)/(1+x)}$.

\begin{remark}[Infinite input is a separate question]\label{e3:rem:infinite-input}
The present construction does not assert that no distinguished global sine
or cosine exists. Ehrlich and Kaplan construct such functions using the
initial integer part of $\No$ \cite[\S11]{EK}; their construction supplies
additional structure and agrees with the finite theory used here. What is
unnecessary for geometry is choosing a phase for every infinite input.
The naive zero-centered sine series at $\omega$ is not a strongly summable
Hahn series, regardless of whether a separate global definition is chosen.
Part~II uses that global phase explicitly, with its period class, in
\cref{e3:so:sec:angles}.
\end{remark}

\part{Coordinate geometry of \texorpdfstring{$\No^3$}{No\^{}3}}\label{e3:part:geometry}

Part~I is source~04. Rotations enter as one tool among several, and the shared
rotation algebra of both sources is in \cref{e3:sec:rotations}.

\section{Vectors, distances, orientation, and volume}\label{e3:sec:vectors}

For $\vv x=(x_1,x_2,x_3)$ and $\vv y=(y_1,y_2,y_3)$ set
\[
 \vv x\cdot\vv y=\sum_{j=1}^3x_jy_j,\qquad
 \norm{\vv x}=\sqrt{\vv x\cdot\vv x},\qquad
 d(P,Q)=\norm{Q-P}.
\]
The cross product is given by the usual determinant formula. We use the
orientation for which $\vv e_1\times\vv e_2=\vv e_3$.

\begin{theorem}[Euclidean vector identities over the scalar field]\label{e3:thm:vector}
For all vectors in $F^3$,
\begin{align}
 \norm{\vv x\times\vv y}^{2}
 &=\norm{\vv x}^{2}\norm{\vv y}^{2}-(\vv x\cdot\vv y)^2,\label{e3:eq:lagrange}\\
 \vv x\times(\vv y\times\vv z)
 &=\vv y(\vv x\cdot\vv z)-\vv z(\vv x\cdot\vv y),\label{e3:eq:triple}\\
 (\vv x\times\vv y)\cdot(\vv z\times\vv w)
 &=(\vv x\cdot\vv z)(\vv y\cdot\vv w)
 -(\vv x\cdot\vv w)(\vv y\cdot\vv z).\label{e3:eq:crossdot}
\end{align}
Moreover $|\vv x\cdot\vv y|\le\norm{\vv x}\norm{\vv y}$ and
$\norm{\vv x+\vv y}\le\norm{\vv x}+\norm{\vv y}$.
For nonzero $\vv x,\vv y$, equality in the latter holds exactly when
$\vv y=\lambda\vv x$ for some $\lambda>0$.
\end{theorem}
\begin{proof}
The three displayed identities follow by expansion in coordinates over any
commutative ring. Positivity of the left side of \cref{e3:eq:lagrange} and
nonnegative square roots give Cauchy--Schwarz. Squaring the proposed triangle
inequality reduces it to Cauchy--Schwarz. Equality requires
$\vv x\cdot\vv y=\norm{\vv x}\norm{\vv y}$, hence
$\vv x\times\vv y=0$. The two vectors are then proportional, and the positive
dot product makes the proportionality factor positive. The converse is direct.
Source~09 proves Cauchy--Schwarz by a second finite route: for $\vv y\ne0$,
minimizing the quadratic $\norm{\vv x-t\vv y}^2$ algebraically at
$t=(\vv x\cdot\vv y)/(\vv y\cdot\vv y)$ gives the inequality.
\end{proof}

Thus $d$ is a positive-definite, symmetric, triangle-inequality distance with
values in $F_{\ge0}$. It is a \emph{field-valued} distance, not an ordinary
real-valued metric. The coordinate formulas do not implicitly make $F^3$ a
Banach space over $\R$.

The cross product is represented by the matrix
\begin{equation}\label{e3:so:eq:hat}
 \widehat{x}=
 \begin{pmatrix}0&-x_3&x_2\\x_3&0&-x_1\\-x_2&x_1&0\end{pmatrix},
 \qquad \widehat{x}\,y=x\times y .
\end{equation}
Source~09 records the following polynomial identities, valid over every field
of characteristic different from two:
\begin{align}
 \widehat{x}^{\,2}&=xx\transpose-\norm{x}^2I,\qquad
 \widehat{x}^{\,3}=-\norm{x}^2\widehat{x},\label{e3:so:eq:hatpowers}\\
 [\widehat{x},\widehat{y}]&=\widehat{x\times y}.\label{e3:so:eq:hatbracket}
\end{align}
Applied to a vector, the first identity is the case of equal first two
arguments of the vector triple product \cref{e3:eq:triple}. No limiting process
is involved.

For nonzero vectors define their unoriented angle by
\begin{equation}\label{e3:eq:vectorangle}
 \ang(\vv x,\vv y)=
 \atantwo\bigl(\norm{\vv x\times\vv y},\vv x\cdot\vv y\bigr)
 \in[0,\pi].
\end{equation}
Lagrange's identity shows that its cosine is the normalized dot product and
its sine is the normalized cross-product norm. If the vectors lie in a plane
with specified unit normal $\vv n$, their oriented angle is instead
$\atantwo(\vv n\cdot(\vv x\times\vv y),\vv x\cdot\vv y)$.

\subsection{Areas and tetrahedra}

The area of a flat triangle $PQR$ and the oriented volume of a tetrahedron
$PQRS$ are
\begin{equation}\label{e3:eq:flatvol}
 \area(PQR)=\frac12\norm{(Q-P)\times(R-P)},\qquad
 V_{\mathrm{or}}(PQRS)=\frac16\det[Q-P\ R-P\ S-P].
\end{equation}
The unoriented volume is $|V_{\mathrm{or}}|$. These are algebraic definitions:
finite dissection and linear change-of-variables identities follow from
multilinearity. They preserve orientation and multiplicity information at
infinitesimal scales.

If $M$ is the matrix of the three edge vectors from $P$, then
\[
 \det(M\transpose M)=\det(M)^2=36\vol(PQRS)^2.
\]
In terms of the six pairwise distances $d_{ij}$, this becomes
\begin{equation}\label{e3:eq:cayleymenger}
 288\vol(PQRS)^2=
 \det\begin{pmatrix}
 0&1&1&1&1\\
 1&0&d_{01}^2&d_{02}^2&d_{03}^2\\
 1&d_{01}^2&0&d_{12}^2&d_{13}^2\\
 1&d_{02}^2&d_{12}^2&0&d_{23}^2\\
 1&d_{03}^2&d_{13}^2&d_{23}^2&0
 \end{pmatrix}.
\end{equation}
Indeed $2M_i\cdot M_j=d_{0i}^2+d_{0j}^2-d_{ij}^2$;
subtracting the corresponding rows and columns in the displayed determinant
reduces it to eight times $\det(M\transpose M)$.

\subsection{Finite-dimensional minimization does not need completeness}

Let $V\subset F^3$ have an orthonormal basis $\vv e_1,\ldots,\vv e_k$.
Gram--Schmidt constructs one whenever a finite independent spanning family
is given; it uses only finitely many nonzero divisions and positive square
roots, so it constructs orthonormal frames in $F^3$ exactly as over $\R$, with
no completeness assumption. The orthogonal projection
\[
 \operatorname{proj}_V(\vv x)=\sum_{j=1}^k(\vv x\cdot\vv e_j)\vv e_j
\]
satisfies, for $\vv y\in V$,
\[
 \norm{\vv x-\vv y}^2=
 \norm{\vv x-\operatorname{proj}_V\vv x}^2+
 \norm{\operatorname{proj}_V\vv x-\vv y}^2.
\]
It is therefore the unique nearest point. This proves a useful minimization
theorem by a finite identity, without a compactness argument or an infimum
construction. The same method gives least squares for full-column-rank
matrices: the minimizer of $\norm{A\vv x-\vv b}^2$ is
$(A\transpose A)^{-1}A\transpose\vv b$.

\section{Affine analytic geometry}\label{e3:sec:affine}

\subsection{Lines, planes, and incidence}

A line is $L=P+F\vv u$ with $\vv u\ne0$. A plane is
$\Pi=\{X:\vv n\cdot X=h\}$ with $\vv n\ne0$.
A line meets this plane uniquely when $\vv n\cdot\vv u\ne0$, at parameter
\begin{equation}\label{e3:eq:lineplane}
 t=\frac{h-\vv n\cdot P}{\vv n\cdot\vv u}.
\end{equation}
A zero denominator means that the line lies in the plane if the numerator
also vanishes, and is disjoint from it otherwise. An infinitesimal nonzero
denominator is not a reason to declare parallelism: it can produce an
infinitely distant intersection.

Two nonparallel planes have line direction $\vv n_1\times\vv n_2$.
Three planes meet uniquely exactly when their normal determinant is nonzero.
All remaining cases are ordinary finite-rank alternatives. For two lines
$P+F\vv u$ and $Q+F\vv v$ with $\vv u\times\vv v\ne0$, they intersect
if and only if
\[
 (Q-P)\cdot(\vv u\times\vv v)=0.
\]
Otherwise they are skew.

The perpendicular foot and distance from $P$ to $\Pi$ are
\begin{equation}\label{e3:eq:pointplane}
 P_\Pi=P-\frac{\vv n\cdot P-h}{\norm{\vv n}^2}\vv n,
 \qquad d(P,\Pi)=\frac{|\vv n\cdot P-h|}{\norm{\vv n}}.
\end{equation}
For a line $Q+F\vv u$, they are
\[
 P_L=Q+\frac{(P-Q)\cdot\vv u}{\norm{\vv u}^2}\vv u,
 \qquad d(P,L)=\frac{\norm{(P-Q)\times\vv u}}{\norm{\vv u}}.
\]
Reflection in $\Pi$ sends $P$ to $2P_\Pi-P$. These formulas are valid at every
surreal scale.

\begin{theorem}[Distance and closest points of nonparallel lines]\label{e3:thm:skew}
Let $\vv r=Q-P$ and put
\[
 A=\vv u^2,\quad B=\vv u\cdot\vv v,\quad C=\vv v^2,
 \quad D_L=AC-B^2>0,
\]
where $\vv u^2$ means $\vv u\cdot\vv u$. The unique closest parameters on
$P+t\vv u$ and $Q+s\vv v$ are
\begin{equation}\label{e3:eq:skewparams}
 t=\frac{C(\vv r\cdot\vv u)-B(\vv r\cdot\vv v)}{D_L},\qquad
 s=\frac{B(\vv r\cdot\vv u)-A(\vv r\cdot\vv v)}{D_L}.
\end{equation}
Their distance is
\begin{equation}\label{e3:eq:skewdist}
 d(L_1,L_2)=\frac{|\vv r\cdot(\vv u\times\vv v)|}
 {\norm{\vv u\times\vv v}}.
\end{equation}
\end{theorem}
\begin{proof}
The difference $\vv r+s\vv v-t\vv u$ must be perpendicular to both direction
vectors. Its two dot products give the linear system
$At-Bs=\vv r\cdot\vv u$, $Bt-Cs=\vv r\cdot\vv v$, whose solution is
\cref{e3:eq:skewparams}. The remaining difference is parallel to
$\vv u\times\vv v$, giving \cref{e3:eq:skewdist}. Every other pair of points
adds a vector in $\Span(\vv u,\vv v)$ to this perpendicular difference.
Pythagoras proves minimality and uniqueness of the two parameters.
\end{proof}

For parallel lines use the point-to-line distance; the closest pair is then
not unique. Angles between unoriented lines or planes are conventionally
chosen acute, by taking an absolute value in the normalized dot product.
Oriented normal vectors instead retain the supplementary-angle distinction.
The acute line--plane angle $\theta$ for line direction $\vv u$ and
plane normal $\vv n$ satisfies
$\sin\theta=|\vv u\cdot\vv n|/(\norm{\vv u}\norm{\vv n})$.
The line--normal angle is its complement.

\subsection{Spheres, circles, and discriminants}

The sphere with center $C$ and radius $R>0$ is
$\{X:\norm{X-C}^2=R^2\}$. For the line $P+t\vv u$, let $t_0$ be the
parameter of its perpendicular foot from $C$, and let $h=d(C,L)$.
Completing the square gives
\begin{equation}\label{e3:eq:linesphere}
 \norm{\vv u}^2(t-t_0)^2=R^2-h^2.
\end{equation}
There are two, one, or no intersections according as $R^2-h^2$ is positive,
zero, or negative. The two parameters, when distinct, are
$t_0\pm\sqrt{R^2-h^2}/\norm{\vv u}$.
A plane at distance $h$ from the center cuts a circle of radius
$\sqrt{R^2-h^2}$ when $h<R$, a point when $h=R$, and nothing when $h>R$.

For two spheres with distinct centers $C_1,C_2$, let
$d=\norm{C_2-C_1}>0$ and $\vv e=(C_2-C_1)/d$. Their intersection lies in
the radical plane perpendicular to $\vv e$, at signed distance
\[
 q=\frac{R_1^2-R_2^2+d^2}{2d}
\]
from $C_1$. Its candidate circle has center $C_1+q\vv e$ and squared radius
$R_1^2-q^2$. The sign of this expression gives the complete classification.
Coincident centers must be handled separately. These formulas never compare
an infinitesimal with zero by rounding it away.

Four affinely independent points determine a unique sphere. Translate the
first point to zero and write the other three as columns of $M$. Its center
$\vv c$ solves
\begin{equation}\label{e3:eq:tetra-sphere}
 M\transpose\vv c=\frac12
 (\norm{M_1}^2,\norm{M_2}^2,\norm{M_3}^2)\transpose.
\end{equation}
The radius is $\norm{\vv c}$. Small nonzero $\det M$ can make this center
infinite even when all four input points are finite.

\subsection{Quadrics and principal axes}\label{e3:sub:quadrics}

A quadric has an equation
\[
 \vv x\transpose A\vv x+2\vv b\transpose\vv x+c=0,
 \qquad A=A\transpose.
\]
Over a real-closed field, a symmetric matrix has an orthonormal eigenbasis.
Here is a proof that does not use compactness of a real unit sphere. A root
$\lambda$ of its characteristic polynomial lies in $F[i]$, and there is a
nonzero vector $z$ with $Az=\lambda z$. Since
$\bar z\transpose A z=\lambda\bar z\transpose z$ is real and
$\bar z\transpose z=\sum |z_j|^2>0$, $\lambda\in F$. Its eigenspace has a
nonzero real vector. Normalize it, restrict to its invariant perpendicular
complement, and continue by induction. The argument works verbatim for
symmetric $d\times d$ matrices over any real closed field; Part~II uses it in
\cref{e3:so:lem:positiveroot}.

After an orthogonal change of coordinates the quadric becomes
$\sum\lambda_j y_j^2+2\sum b'_j y_j+c=0$.
For $\lambda_j\ne0$, completing the square removes the corresponding linear
term. If all $\lambda_j$ are nonzero this gives the familiar central forms,
classified by coefficient signs and the remaining constant. Ellipsoids,
hyperboloids, cones, empty loci, and point loci are distinguished exactly as
over $\R$, with arbitrary surreal semiaxes. If a null direction retains a
linear term, solving for that coordinate gives a paraboloid-type form.
If all linear terms in the nullspace vanish, those directions are cylindrical.
Thus the classification is algebraic, while the scales of the axes carry
additional information.

\subsection{Affine points are not projective ideal points}

Homogeneous coordinates embed $F^3$ into $\mathbb P^3(F)$ by
$\vv x\mapsto[x_1:x_2:x_3:1]$. An affine point such as $(\omega,0,0)$ is
not $[1:0:0:0]$: the last homogeneous coordinate distinguishes them exactly.
Projective completion is a separate incidence construction. It does not
identify a very large finite-coordinate point with an ideal direction.

\section{Rigid motions, rotations, quaternions, and screws}\label{e3:sec:rotations}

This section prints once the rotation algebra that both sources develop
(\cref{e3:sub:provenance}). Except for \cref{e3:thm:finiteaxis}, everything here
is finite algebra over a real closed field $F$. Source~04 stated it over its
Hahn workspace, which is such a field; source~09 stated it over an arbitrary
real closed field. We write $G(F)=\SO(3,F)$ and use the cross-product matrix
$\widehat x$ of \cref{e3:so:eq:hat}.

\begin{theorem}[Classification of field-distance isometries]\label{e3:thm:isometries}\label{e3:so:thm:isometries}
Every map $\Phi:F^3\to F^3$ preserving $d(P,Q)=\norm{P-Q}$ has the form
\[
 \Phi(\vv x)=A\vv x+\vv b,\qquad A\transpose A=I.
\]
In particular, it is automatically onto. No continuity or affine assumption
is needed. If $\Phi$ fixes the origin, it is $F$-linear with orthogonal matrix
$A$, and it preserves the orientation of the standard basis exactly when
$\det A=1$. Conversely, every such matrix preserves the norm and fixes the
origin.
\end{theorem}
\begin{proof}
Subtract $\Phi(0)$ and put $\vv q_j=\Phi(\vv e_j)-\Phi(0)$. Distances between
$0,\vv e_1,\vv e_2,\vv e_3$ and polarization,
\[
 \bigl(\Phi(\vv x)-\Phi(0)\bigr)\cdot\bigl(\Phi(\vv y)-\Phi(0)\bigr)
 =\tfrac12\bigl(\norm{\vv x}^2+\norm{\vv y}^2-\norm{\vv x-\vv y}^2\bigr)
 =\vv x\cdot\vv y,
\]
show $\vv q_i\cdot\vv q_j=\delta_{ij}$. For arbitrary $\vv x$, the same identity
with $\vv y=\vv e_j$ gives $(\Phi(\vv x)-\Phi(0))\cdot\vv q_j=x_j$. The $\vv q_j$
form an orthonormal basis, so $\Phi(\vv x)-\Phi(0)=\sum x_j\vv q_j$. Its matrix
is orthogonal and therefore invertible. Its determinant is $1$ or $-1$, and
preserving the orientation of the given basis is precisely the first
alternative.
\end{proof}

The theorem excludes a potential ambiguity. A field automorphism applied to
coordinates may carry squared distances to their images under that automorphism.
That is a semilinear transformation of a different geometry, not an isometry
preserving every surreal distance as an actual value. The rotations of $F^3$
are therefore exactly the elements of $G(F)$.

\begin{proposition}[Bounds and equivariance]\label{e3:so:prop:bounds}
For $A\in G(F)$, every entry satisfies $\abs{A_{ij}}\le1$, and
\[
 (Ax)\times(Ay)=A(x\times y),\qquad
 A\widehat{x}A^{-1}=\widehat{Ax}.
\]
Thus no rotation matrix has an infinite entry, even when its vectors have
infinite coordinates.
\end{proposition}
\begin{proof}
Each column has sum of squares one. The cross-product identity follows by
pairing both sides with $Az$: both scalar triple products are multiplied by
$\det A=1$. The hat identity follows immediately.
\end{proof}

\subsection{Axis--angle representation}

\begin{theorem}[Axis theorem and Rodrigues formula]\label{e3:thm:rodrigues}\label{e3:so:thm:axis}
Every $A\in G(F)$ fixes a unit vector $n$. There exist $c,s\in F$ with
$c^2+s^2=1$ such that
\begin{equation}\label{e3:so:eq:rodrigues}
 A=\Rot(n;c,s)=cI+(1-c)nn\transpose+s\widehat n.
\end{equation}
If $A\ne I$, its fixed space is exactly $Fn$. Replacing $(n,s)$ by $(-n,-s)$
leaves the rotation unchanged. Conversely $\Rot(n;c,s)\in G(F)$ for every unit
$n$ and every $(c,s)$ on the unit circle. When $(c,s)=(\cos\theta,\sin\theta)$
for a finite angle $\theta$, which is always possible over a Hahn workspace by
\cref{e3:thm:finiteaxis}, the formula reads
\begin{equation}\label{e3:eq:rodrigues}
 \Rot(\vv n;\cos\theta,\sin\theta)\,\vv x=
 (\cos\theta)\vv x+(\sin\theta)(\vv n\times\vv x)
 +(1-\cos\theta)(\vv n\cdot\vv x)\vv n.
\end{equation}
\end{theorem}
\begin{proof}
Since $A\transpose=A^{-1}$ and $\det A=1$,
\[
 \det(A-I)=\det(A\transpose-I)=\det(A^{-1}(I-A))
 =\det(I-A)=-\det(A-I).
\]
Characteristic zero implies $\det(A-I)=0$. Normalize a nonzero fixed vector
to obtain $n$. Its perpendicular plane is invariant under $A$. In an
oriented orthonormal basis $(u,n\times u,n)$, the restriction to that plane
is an orthogonal $2\times2$ matrix of determinant one, necessarily
$\left(\begin{smallmatrix}c&-s\\s&c\end{smallmatrix}\right)$.
This gives \cref{e3:so:eq:rodrigues}. Conversely, on $n^\perp$ the operator
$J x=n\times x$ satisfies $J^2=-I$, so $cI+sJ$ is an orientation-preserving
orthogonal map of that plane; on the axis the formula is the identity. Hence
$\Rot(n;c,s)\in G(F)$.

The determinant of the planar matrix minus the identity is
$(c-1)^2+s^2=2(1-c)$. It vanishes only when $c=1,s=0$, which gives the
identity rotation. Otherwise the perpendicular plane contains no fixed vector.
(Source~04's variant: if two independent vectors were fixed, their cross
product would also be fixed, forcing $A=I$.) Expanding $nn\transpose$ and
$\widehat n$ on a vector gives \cref{e3:eq:rodrigues}.
\end{proof}

The pair $(c,s)$ is an \emph{algebraic angle}. It needs neither a sine function
on all of $F$ nor a real number measuring arc length. It is often the safest
primary representation of a rotation over an arbitrary real closed field.
In the basis used in the proof,
\begin{equation}\label{e3:so:eq:trace}
 \tr A=1+2c,\qquad -1\le\tr A\le3,
 \qquad \chi_A(X)=(X-1)(X^2-2cX+1).
\end{equation}
Over $F(i)$, the eigenvalues are $1,c+is,c-is$. A half-turn is exactly a
rotation with $c=-1,s=0$, so
\begin{equation}\label{e3:so:eq:halfturn}
 H_n=2nn\transpose-I,\qquad H_n^2=I,\qquad \tr H_n=-1.
\end{equation}
Its axis is an unoriented line. These are all nonidentity involutions.

For a non-half-turn, the skew part is
$A-A\transpose=2s\widehat n$. At a half-turn it vanishes, although the rotation is
not the identity. Algorithms that infer the axis only from the skew part must
handle this case separately.

\begin{theorem}[Every rotation has a finite axis--angle description]\label{e3:thm:finiteaxis}
Let $F$ be a real-closed Hahn workspace as in \cref{e3:sec:scalars}; by
localization the statement applies to $G(\No)$. Every $A\in G(F)$ has a
representation
\[
 A=\Rot(n;\cos\theta,\sin\theta),\qquad 0\le\theta\le\pi .
\]
For $A\ne I$, the number $\theta$ is unique. The axis is oriented uniquely when
$0<\theta<\pi$, by the requirement $\sin\theta>0$, and is unoriented at
$\theta=\pi$. Thus every rotation, even one involving infinitely separated
spatial points, has a finite axis--angle description.
\end{theorem}
\begin{proof}
By \cref{e3:thm:rodrigues} and the finite phase theorem \cref{e3:thm:phase},
$A=\Rot(n;\cos\theta,\sin\theta)$ for some finite $\theta$. Reduce $\theta$ by
an ordinary integral multiple of $2\pi$ to $(-\pi,\pi]$. To handle an
infinitesimal excursion past an endpoint, shift once more by $2\pi$; the
representative must lie in the actual surreal interval, not merely have
standard part there. Reverse the axis if the angle is negative, obtaining
$[0,\pi]$. On $(0,\pi)$ the finite sine is positive: away from the endpoints its
standard part is positive, and at an infinitesimal endpoint displacement the
leading Taylor term decides the sign. Equality of two circle values is exactly
equality modulo $2\pi\Z$. With the axis-reversal convention, this gives
uniqueness of $\theta$ and the stated axis ambiguity.
\end{proof}

For unit vectors $\vv u\ne-\vv v$, put $c=\vv u\cdot\vv v$ and
$\vv k=\vv u\times\vv v$. The shortest-angle rotation taking $\vv u$ to
$\vv v$ is
\begin{equation}\label{e3:eq:minimalrotation}
 A(\vv u,\vv v)=I+\widehat{\vv k}+
                   \frac{\widehat{\vv k}^{\,2}}{1+c}.
\end{equation}
The formula also gives the identity when $\vv u=\vv v$.
Using $\vv k\times\vv u=\vv v-c\vv u$ and
$\vv k\times(\vv k\times\vv u)=-(1-c^2)\vv u$ verifies the mapping.
At an antipodal pair the formula is undefined for a genuine reason: there
are many half-turns taking one vector to the other.

\subsection{Quaternion representation}\label{e3:sub:quaternions}

Write a quaternion as $(a,\vv b)\in\HH_F=F\oplus F^3$ with multiplication
\begin{equation}\label{e3:eq:quaternionmul}
 (a,\vv b)(c,\vv d)=
 (ac-\vv b\cdot\vv d,\ a\vv d+c\vv b+\vv b\times\vv d).
\end{equation}
The basis $(1,i,j,k)$ satisfies $i^2=j^2=k^2=ijk=-1$. This is an associative
algebra; associativity follows either from Hamilton's relations and the faithful
matrix representation \cref{e3:so:eq:su2}, or from the triple-product identity.
Quaternion conjugation and norm square are
\[
 \overline{(a,\vv b)}=(a,-\vv b),\qquad N(a,\vv b)=a^2+\norm{\vv b}^2.
\]
Conjugation reverses products, $q\bar q=N(q)$, and $N(pq)=N(p)N(q)$. Since a
nonzero sum of squares is positive, $q^{-1}=\bar q/N(q)$ for every nonzero $q$.
The center is exactly $F$. These are ordered-field facts, not real-analytic
facts \cite{Voight,RepoQuaternion}.

Put $\Sp(1,F)=\{q\in\HH_F:N(q)=1\}$. Unit quaternions act on pure vectors by
conjugation, $p(q)(\vv x)=q(0,\vv x)q^{-1}$. Direct expansion gives
\begin{equation}\label{e3:eq:quatrotation}
 p(a,\vv b)\,\vv x=(a^2-\norm{\vv b}^2)\vv x
       +2(\vv b\cdot\vv x)\vv b+2a(\vv b\times\vv x),
\end{equation}
which is \texttt{squat:eq:rodrigues} of the surquaternion report. It preserves
the imaginary subspace and the norm. Its determinant is $+1$ without an appeal
to connectedness. For $\vv b\ne0$, put $n=\vv b/\norm{\vv b}$ and compare with
\cref{e3:so:eq:rodrigues} using $c=a^2-\norm{\vv b}^2$ and $s=2a\norm{\vv b}$. If
$\vv b=0$, the map is the identity.

\begin{theorem}[The spin exact sequence; \texttt{squat:thm:rotations}]\label{e3:thm:spin}\label{e3:so:thm:spin}
Conjugation gives a surjective group homomorphism with precisely two elements
in every fiber:
\begin{equation}\label{e3:so:eq:spinexact}
 1\longrightarrow\{1,-1\}\longrightarrow\Sp(1,F)
 \xrightarrow{\ p\ }G(F)\longrightarrow1.
\end{equation}
In particular, $\Spin(3,F)\cong\Sp(1,F)$.
\end{theorem}

This is Theorem~4.4 (\texttt{squat:thm:rotations}) of the surquaternion
report \cite{RepoQuaternion}, which also identifies $G(F)$ with the
automorphism group of $\HH_F$. Both sources reprove it by the argument of that
theorem's first surjectivity route: a quaternion commuting with every pure
vector is central, and a fixed axis gives a lift. The proof is not repeated
here. Source~09's explicit lift of $\Rot(n;c,s)$ is used below. If $c\ne-1$,
\begin{equation}\label{e3:so:eq:spinlift}
 a=\sqrt{(1+c)/2}>0,\qquad u=\frac{s}{2a}\,n
\end{equation}
give $a^2+\norm{u}^2=1$, $a^2-\norm{u}^2=c$ and $2au=sn$, so $(a,u)$ is a lift. If
$c=-1$, the unit imaginary quaternion $(0,n)$ lifts $H_n$. For a finite angle,
source~04's lift is $(\cos(\theta/2),\sin(\theta/2)\vv n)$, which gives
\cref{e3:eq:rodrigues}. The standard Clifford-algebra spin group of the positive
definite ternary form has this Hamilton realization;
\cref{e3:so:eq:spinexact} supplies an entirely explicit model without Clifford
terminology.

The phrase ``double cover'' here denotes this two-to-one algebraic group
homomorphism, not an unproved topological covering assertion in the fine
topology. The category in which it is a two-sheeted covering with the familiar
fundamental groups is identified in \cref{e3:so:thm:sahomotopy}.

\subsection{The rational Cayley chart}\label{e3:sub:cayley}

\begin{theorem}[The matrix Cayley chart]\label{e3:so:thm:cayley}
For $t\in F^3$, define
\begin{equation}\label{e3:so:eq:cayley}
 \cay(t)=(I+\widehat t)(I-\widehat t)^{-1}
 =\frac{(1-\norm{t}^2)I+2tt\transpose+2\widehat t}{1+\norm{t}^2}.
\end{equation}
This is a bijection from $F^3$ onto the rotations which are not half-turns.
Its inverse is
\begin{align}
 \widehat t&=(A-I)(A+I)^{-1},\label{e3:so:eq:cayleyinv}\\
 t&=\frac{(A_{32}-A_{23},\ A_{13}-A_{31},\ A_{21}-A_{12})\transpose}
                  {1+\tr A}.\label{e3:so:eq:cayleyinvcoords}
\end{align}
Moreover, whenever $1-t\cdot u\ne0$,
\begin{equation}\label{e3:so:eq:cayleyproduct}
 \cay(t)\cay(u)=\cay\left(\frac{t+u+t\times u}{1-t\cdot u}\right).
\end{equation}
If $1-t\cdot u=0$, the product is a half-turn.
\end{theorem}
\begin{proof}
The determinant of $I-\widehat t$ is $1+\norm{t}^2>0$.
Using \cref{e3:so:eq:hatpowers} gives the displayed rational expression.
Equivalently it is the rotation of the unit quaternion
\[
 q(t)=\frac{(1,t)}{\sqrt{1+\norm{t}^2}}.
\]
Its scalar coordinate is nonzero, so the rotation is not a half-turn.
Conversely, the lift \cref{e3:so:eq:spinlift} of a non-half-turn has
nonzero scalar part, and dividing its imaginary part by that scalar gives
$t$. Solving the matrix equation gives \cref{e3:so:eq:cayleyinv}; taking skew
coordinates gives \cref{e3:so:eq:cayleyinvcoords}. In particular,
\begin{equation}\label{e3:so:eq:caytrace}
 \tr\cay(t)=\frac{3-\norm{t}^2}{1+\norm{t}^2},\quad
 1+\tr\cay(t)=\frac4{1+\norm{t}^2},\quad
 \det(\cay(t)+I)=\frac8{1+\norm{t}^2}.
\end{equation}
Finally, before normalization,
$(1,t)(1,u)=(1-t\cdot u,t+u+t\times u)$.
A nonzero scalar part gives the quotient in \cref{e3:so:eq:cayleyproduct}. If that
scalar vanishes, the product is a nonzero imaginary quaternion and its rotation
is a half-turn.
\end{proof}

In vector form, and with $\vv t=\tan(\theta/2)\vv n$ for a rotation of finite
angle $\theta$ away from half-turns, source~04 writes the same chart as
\begin{equation}\label{e3:eq:cayley}
 \cay(\vv t)\,\vv x=\vv x+\frac{2}{1+\norm{\vv t}^2}
   \bigl(\vv t\times\vv x+\vv t\times(\vv t\times\vv x)\bigr).
\end{equation}
This rational chart can have infinite coordinates near a half-turn,
although every entry of the rotation matrix remains between $-1$ and $1$
(\cref{e3:so:prop:bounds}). It is the three-dimensional analogue of the planar
half-angle chart \texttt{trigonometry:thm:cayley}.

This is a chart on $G(F)$, not the quaternion stereographic chart
$(1-x)^{-1}(1+x)$ on $\Sp(1,F)$ of \texttt{squat:prop:charts}. The two are
related, but their parameters have different half-angle factors. Confusing them
changes the composition law and the leading infinitesimal rotation by a factor
of two.

The half-turn locus is naturally $\mathbb P^2(F)$, by $[n]\mapsto2nn\transpose-I$
after normalization. The chart and its translates show that the group has
algebraic dimension three. After extending to $F(i)$, its standard algebraic
form is $\operatorname{PGL}_2$; thus the algebraic group is connected and
absolutely simple of adjoint type. These are algebraic assertions. They do not
imply that its $F$-point group is abstractly simple or that its fine topology
is connected. For $F\supsetneq\R$, and for $\No$,
\cref{e3:so:sec:standardpart,e3:so:sec:topology} show that neither holds.

\subsection{Screw motions}

Let $\Phi(\vv x)=A\vv x+\vv b$ with $A\in\SO(3,F)$ nonidentity. Take its
unit axis $\vv n$, put $h=\vv b\cdot\vv n$, and
$\vv b_\perp=\vv b-h\vv n$. On $\vv n^\perp$, $I-A$ is invertible, so
there is a unique $\vv a\perp\vv n$ satisfying
$(I-A)\vv a=\vv b_\perp$. Therefore
\[
 \Phi(\vv a+t\vv n)=\vv a+(t+h)\vv n.
\]
The motion is rotation about the line $\vv a+F\vv n$, followed by translation
$h$ along it. With $0<\theta\le\pi$,
\begin{equation}\label{e3:eq:screwaxis}
 \vv a=\frac12\left(\vv b_\perp+
       \cot(\theta/2)\,\vv n\times\vv b_\perp\right).
\end{equation}
This follows by inverting $(1-\cos\theta)I-\sin\theta J$ on the perpendicular
plane, with $J\vv x=\vv n\times\vv x$. An infinitesimal rotation with appreciable
transverse translation can have an infinitely distant screw axis. The identity
rotation is treated separately as a pure translation.

\section{Coordinates and ordinary triangle trigonometry in space}\label{e3:sec:coordinates}

\subsection{Cylindrical and spherical coordinates}

For $\vv x=(x,y,z)$ put $r=\sqrt{x^2+y^2}$. When $r>0$, cylindrical
coordinates are
\[
 x=r\cos\phi,\qquad y=r\sin\phi,\qquad z=z,
 \qquad \phi=\atantwo(y,x)\pmod{2\pi}.
\]
At $r=0$, azimuth is undetermined. There is no restriction on the scale of
$r$ or $z$.

For $\rho=\norm{\vv x}>0$, define colatitude
$\theta=\acos(z/\rho)\in[0,\pi]$. Away from the polar axis choose the same
azimuth $\phi$. Then
\begin{equation}\label{e3:eq:sphericalcoords}
 (x,y,z)=(\rho\sin\theta\cos\phi,
          \rho\sin\theta\sin\phi,\rho\cos\theta).
\end{equation}
Every nonzero vector has this representation. At the poles the azimuth is
free; at the origin both angles are free. The radial coordinate may be
infinite, but both angular coordinates are finite. Choosing an azimuth branch
introduces a coordinate seam, not a singularity of the sphere.

Differentiate the coordinate map with respect to its external variables.
Its three coordinate derivatives are mutually perpendicular and have squared
norms $1,\rho^2,\rho^2\sin^2\theta$. Therefore
\begin{equation}\label{e3:eq:sphericalmetric}
 ds^2=d\rho^2+\rho^2d\theta^2+
                    \rho^2\sin^2\theta\,d\phi^2,
 \qquad J=\rho^2\sin\theta.
\end{equation}
These are finite identities in restricted-analytic coordinate derivatives.
They do not assert an integral transformation theorem for every function on
$F^3$. On a sphere of radius $R$, the local area density is
$R^2\sin\theta$ and the tangent metric is obtained by setting $d\rho=0$.

\subsection{Flat triangles retain their usual formulas}

Every noncollinear triple of points lies in a unique affine plane. An
orthonormal basis of that plane reduces its geometry to $F^2$ without
changing lengths. Let its opposite side lengths be $A,B,C>0$, its angles
$\alpha,\beta,\gamma\in(0,\pi)$, and its area $\Delta>0$. Then
\begin{align}
 \alpha+\beta+\gamma&=\pi,\label{e3:eq:planeangles}\\
 A^2&=B^2+C^2-2BC\cos\alpha,\label{e3:eq:planecos}\\
 \frac{A}{\sin\alpha}
 &=\frac{B}{\sin\beta}=\frac{C}{\sin\gamma}
   =\frac{ABC}{2\Delta}=2R_c,\label{e3:eq:planesine}\\
 \Delta^2&=s(s-A)(s-B)(s-C),\qquad s=(A+B+C)/2.\label{e3:eq:heron}
\end{align}
These are the usual classical identities \cite[\S4.42(ii)]{DLMF}, now with
field-valued lengths.

For completeness, put the first two vertices at $(0,0)$ and $(C,0)$, and the
third in the upper half-plane. If the ray from the first vertex to the third
has phase $\alpha$ and the ray from the second to the third has phase
$\delta$, orientation gives $0<\alpha<\delta<\pi$. The remaining angles
are $\pi-\delta$ and $\delta-\alpha$, proving \cref{e3:eq:planeangles}.
Expanding the squared difference proves \cref{e3:eq:planecos}. Computing the
same area from each pair of sides proves the first equalities in
\cref{e3:eq:planesine}. The circumcenter is the solution of two perpendicular
bisector equations, and the chord relation gives $A=2R_c\sin\alpha$.
Finally squaring $2\Delta=BC\sin\alpha$ and substituting the cosine law gives
\[
 16\Delta^2=(A+B+C)(-A+B+C)(A-B+C)(A+B-C),
\]
which proves Heron's formula.

The triangle inequalities are strict exactly in the noncollinear case.
None of these statements requires the side lengths to be finite. For example,
a triangle can have sides of scale $\omega$ and an altitude of scale
$\omega^{-\omega}$; its nonzero area is the product scale, not automatically
zero or infinite.

\section{The surreal sphere and its intrinsic angular distance}\label{e3:sec:sphere}

Let
\[
 S_R^2(C)=\{C+R\vv u:\vv u\cdot\vv u=1\}.
\]
A great circle is its intersection with a plane through $C$. Translating and
rescaling reduces all angular calculations to the unit sphere $S^2(F)$.

\subsection{Minor arcs without a global sine function}

For unit vectors $\vv u\ne-\vv v$, define the minor arc by the algebraic set
\begin{equation}\label{e3:eq:algebraicarc}
 \left\{
 \frac{(1-t)\vv u+t\vv v}{\norm{(1-t)\vv u+t\vv v}}:0\le t\le1
 \right\}.
\end{equation}
The denominator is nonzero: vanishing would make the endpoints antipodal
unless they coincide. If $a=\ang(\vv u,\vv v)\in(0,\pi)$ and
\[
 \vv e=\frac{\vv v-\cos a\,\vv u}{\sin a},
\]
the same arc has the trigonometric parametrization
\begin{equation}\label{e3:eq:arcparam}
 \vv g(s)=\cos s\,\vv u+\sin s\,\vv e,
 \qquad 0\le s\le a.
\end{equation}
To see equality, work in the oriented plane basis $(\vv u,\vv e)$.
Normalized positive combinations of $(1,0)$ and $(\cos a,\sin a)$ have
exactly the phases between $0$ and $a$, as is checked by the signs of their
two oriented determinants with the endpoints. This also proves uniqueness
of the minor arc. Antipodal points instead have a family of semicircles,
indexed by unit directions perpendicular to either endpoint.

Define the angular distance on the unit sphere and the physical distance on
$S_R^2(C)$ by
\begin{equation}\label{e3:eq:angulardistance}
 d_S(\vv u,\vv v)=\acos(\vv u\cdot\vv v),\qquad
 d_R(C+R\vv u,C+R\vv v)=R\,d_S(\vv u,\vv v).
\end{equation}
The length assigned to a minor arc is $Ra$. Subarcs add because their finite
phase differences add. One complete great circle has length $2\pi R$.

\begin{theorem}[Angular triangle inequality]\label{e3:thm:angularmetric}
$d_S$ is a positive-definite $F$-valued distance on $S^2(F)$. For
$a=d_S(\vv u,\vv v)$,
\begin{equation}\label{e3:eq:chordangle}
 \norm{\vv u-\vv v}=2\sin(a/2),\qquad
 \frac{2}{\pi}a\le\norm{\vv u-\vv v}\le a.
\end{equation}
A minor great-circle arc minimizes the sum of arc lengths among all finite
chains of minor great-circle arcs with the same endpoints.
\end{theorem}
\begin{proof}
Only the triangle inequality needs explanation. Set
$a=d_S(\vv u,\vv v)$ and $b=d_S(\vv v,\vv w)$. If $a+b\ge\pi$ the
claim follows from $d_S\le\pi$. Otherwise write
\[
 \vv u=\cos a\,\vv v+\sin a\,\vv e,\qquad
 \vv w=\cos b\,\vv v+\sin b\,\vv f
\]
with unit perpendicular vectors, treating zero angles separately.
Cauchy--Schwarz gives $\vv e\cdot\vv f\ge-1$, hence
$\vv u\cdot\vv w\ge\cos(a+b)$. Cosine is decreasing on $[0,\pi]$, so
$d_S(\vv u,\vv w)\le a+b$. The chord identity follows by squaring.
The inequalities follow from $2t/\pi\le\sin t\le t$ for
$0\le t\le\pi/2$, valid in the restricted-analytic field. Iterating the
triangle inequality proves the finite-chain minimizing assertion.
\end{proof}

This is an intrinsic distance construction with an explicit minimizing class.
No supremum over arbitrary partitions, no real-valued length metric, and no
unrestricted class of nondefinable curves has been silently introduced.

\subsection{Gnomonic and stereographic charts}\label{e3:subsec:charts}

The gnomonic chart on the northern open hemisphere is
\begin{equation}\label{e3:eq:gnomonic}
 \mathrm{Gn}(p,q)=\frac{(p,q,1)}{\sqrt{1+p^2+q^2}},\qquad
 (p,q)=(u_1/u_3,u_2/u_3).
\end{equation}
Great circles become straight lines because $\vv n\cdot\vv u=0$ becomes
$n_1p+n_2q+n_3=0$. Positive-cone spherical triangles become ordinary affine
convex triangles. The chart is valid for arbitrary $p,q\in F$, including
infinite values near the horizon. With $\vv p=(p,q)$ its metric on a sphere
of radius $R$ is
\begin{equation}\label{e3:eq:gnomonicmetric}
 ds^2=R^2\frac{(1+\norm{\vv p}^2)\norm{d\vv p}^2
                 -(\vv p\cdot d\vv p)^2}
                 {(1+\norm{\vv p}^2)^2},
 \qquad J_{\mathrm{Gn}}=\frac{R^2}{(1+\norm{\vv p}^2)^{3/2}}.
\end{equation}
Both follow by direct differentiation of \cref{e3:eq:gnomonic}.

Stereographic projection from the north pole has inverse
\begin{equation}\label{e3:eq:stereo}
 \mathrm{St}(p,q)=\frac{(2p,2q,p^2+q^2-1)}{1+p^2+q^2}.
\end{equation}
It covers the sphere minus that pole, with metric
\[
 ds^2=\frac{4R^2(dp^2+dq^2)}{(1+p^2+q^2)^2}.
\]
The scalar multiple of $dp^2+dq^2$ proves conformality, including at
infinitesimal and infinite chart coordinates. A spherical circle cut out by
$\vv n\cdot\vv u=h$ becomes
\[
 (n_3-h)(p^2+q^2)+2n_1p+2n_2q-(n_3+h)=0,
\]
a circle or line. Infinite chart coordinates are not the omitted pole;
they describe genuine points infinitesimally near it.

\section{Spherical triangles and the Gram-matrix calculus}\label{e3:sec:spherical-laws}

\begin{definition}[Proper spherical triangle]
Three linearly independent unit vectors $\vv u,\vv v,\vv w$ determine the
closed spherical triangle consisting of the normalized nonzero vectors in
\[
 \{\lambda\vv u+\mu\vv v+\nu\vv w:\lambda,\mu,\nu\ge0\}.
\]
Its edges are the minor great-circle arcs. Its interior angles are the angles
between the inward-pointing edge tangents at each vertex.
\end{definition}

Such a triangle lies in an open hemisphere. Indeed, if
$P=[\vv u\ \vv v\ \vv w]$, then
$\vv h=P^{-\mathsf T}(1,1,1)\transpose$ has positive dot product with every
nonzero vector in the positive cone. This algebraic observation fixes which
region and which angles are intended, even when its real standard part is
degenerate. All sides and angles of a proper triangle lie strictly between
$0$ and $\pi$.

Use \cref{e3:eq:xyzconvention} and put
\begin{equation}\label{e3:eq:grammatrix}
 \mathcal G=P\transpose P=
 \begin{pmatrix}1&z&y\\z&1&x\\y&x&1\end{pmatrix},
 \qquad D=\det\mathcal G=1+2xyz-x^2-y^2-z^2=\tau^2>0.
\end{equation}
Write $s_a=\sqrt{1-x^2}=\sin a$ and cyclically $s_b,s_c$.
The unit tangents at $\vv u$ pointing toward $\vv v$ and $\vv w$ are
\begin{equation}\label{e3:eq:tangents}
 T_{uv}=\frac{\vv v-z\vv u}{s_c},\qquad
 T_{uw}=\frac{\vv w-y\vv u}{s_b}.
\end{equation}
Consequently
\begin{equation}\label{e3:eq:angledata}
 \cos\alpha=\frac{x-yz}{s_bs_c},\qquad
 \sin\alpha=\frac{d}{s_bs_c},\qquad
 \alpha=\atantwo(d,x-yz),
\end{equation}
with cyclic formulas for $\beta,\gamma$. The common positive denominator
in the first two formulas is irrelevant to \texttt{atan2}.

\begin{theorem}[Spherical sine and cosine laws]\label{e3:thm:spherical-laws}
Every proper spherical triangle satisfies
\begin{align}
 \cos a&=\cos b\cos c+\sin b\sin c\cos\alpha,\label{e3:eq:sphericalcos}\\
 \frac{\sin\alpha}{\sin a}
 &=\frac{\sin\beta}{\sin b}
  =\frac{\sin\gamma}{\sin c}
  =\frac{d}{\sin a\sin b\sin c}.\label{e3:eq:sphericalsine}
\end{align}
Cyclic permutations give the other cosine laws. A useful additional identity
is
\begin{equation}\label{e3:eq:sphericalfourpart}
 \sin a\cos\beta=
 \cos b\sin c-\sin b\cos c\cos\alpha.
\end{equation}
\end{theorem}
\begin{proof}
The dot product of the two tangents in \cref{e3:eq:tangents} is
$(x-yz)/(s_bs_c)$. Their cross product is parallel to $\vv u$ and has norm
$|\det[\vv u\ \vv v\ \vv w]|/(s_bs_c)$. This proves
\cref{e3:eq:angledata}, and hence both main laws. For
\cref{e3:eq:sphericalfourpart}, substitute the cosine law into the right-hand
side and reduce it to $(y-xz)/s_c$, which is the left-hand side by its cyclic
tangent formula.
\end{proof}

The sine and cosine laws are classical \cite[\S4.42(iii)]{DLMF}. Their
proof here uses only finite vector identities plus the already constructed
finite angle representatives. The same calculation over a field with square
roots still gives the algebraic sine--cosine data even without naming angles.

\subsection{Polar triangles and the cosine law for angles}\label{e3:subsec:polar}

Relabel the vertices so that $\tau>0$ and define inward unit side normals
\begin{equation}\label{e3:eq:polarnormals}
 \vv n_a=\frac{\vv v\times\vv w}{s_a},\quad
 \vv n_b=\frac{\vv w\times\vv u}{s_b},\quad
 \vv n_c=\frac{\vv u\times\vv v}{s_c}.
\end{equation}
For example $\vv n_a\cdot\vv u=d/s_a>0$, and its dot products with the
other two vertices vanish. These normals are independent and positively
oriented: their determinant is $d^2/(s_as_bs_c)>0$.

\begin{theorem}[Polar duality]\label{e3:thm:polar}
The triangle with vertices $\vv n_a,\vv n_b,\vv n_c$ has sides and angles
\begin{equation}\label{e3:eq:polarrelations}
 a^*=\pi-\alpha,\quad b^*=\pi-\beta,\quad c^*=\pi-\gamma,
 \qquad
 \alpha^*=\pi-a,\quad\beta^*=\pi-b,\quad\gamma^*=\pi-c.
\end{equation}
Its polar is the original triangle. In particular,
\begin{equation}\label{e3:eq:anglecosine}
 \cos\alpha=-\cos\beta\cos\gamma+
                 \sin\beta\sin\gamma\cos a,
\end{equation}
and cyclically.
\end{theorem}
\begin{proof}
By \cref{e3:eq:crossdot},
$\vv n_b\cdot\vv n_c=(yz-x)/(s_bs_c)=-\cos\alpha$.
The side is in $(0,\pi)$, so it equals $\pi-\alpha$.
Also
$(\vv w\times\vv u)\times(\vv u\times\vv v)=\tau\vv u$;
after normalization the corresponding inward polar normal is $\vv u$.
This proves involutivity and the angle relations. Apply the side cosine law
to the polar triangle and substitute \cref{e3:eq:polarrelations}.
\end{proof}

\subsection{Half-angle, right-triangle, and haversine formulas}

Let $s=(a+b+c)/2$. The identity
\begin{equation}\label{e3:eq:gramfactor}
 D=4\sin s\sin(s-a)\sin(s-b)\sin(s-c)
\end{equation}
follows by factoring
\[
 D=\bigl(\cos(b-c)-\cos a\bigr)
       \bigl(\cos a-\cos(b+c)\bigr).
\]
For a proper triangle $s<\pi$ and $s-a,s-b,s-c>0$, as proved in the next
section. The half-angle formulas therefore have unambiguous positive roots:
\begin{align}
 \sin(\alpha/2)&=
 \sqrt{\frac{\sin(s-b)\sin(s-c)}{\sin b\sin c}},\label{e3:eq:halfsin}\\
 \cos(\alpha/2)&=
 \sqrt{\frac{\sin s\sin(s-a)}{\sin b\sin c}}.\label{e3:eq:halfcos}
\end{align}
They follow by substituting the cosine law into
$\sin^2(\alpha/2)=(1-\cos\alpha)/2$ and
$\cos^2(\alpha/2)=(1+\cos\alpha)/2$.
Their ratio gives the half-angle tangent formula. Applying the same formulas
to the polar triangle gives half-side formulas in terms of the angles.

For a right triangle with $\gamma=\pi/2$, the Napier relations can be written
in the particularly safe form
\begin{align}
 \cos c&=\cos a\cos b,&
 \sin a&=\sin c\sin\alpha,&
 \sin b&=\sin c\sin\beta,\label{e3:eq:napier1}\\
 \cos\alpha&=\cos a\sin\beta,&
 \cos\beta&=\cos b\sin\alpha,&
 \cos c&=\cot\alpha\cot\beta.\label{e3:eq:napier2}
\end{align}
They follow from the side and angle cosine laws and the sine law. The last
expression is defined since $\sin\alpha,\sin\beta>0$. Formulas involving
$\tan a$, $\tan b$, or $\tan c$ require additional nonzero-cosine checks.
For instance a proper right spherical triangle can itself have a side of
length $\pi/2$.

Define $\hav t=\sin^2(t/2)$. Rearranging the cosine law gives
\begin{equation}\label{e3:eq:haversine}
 \hav a=\hav(b-c)+\sin b\sin c\,\hav\alpha.
\end{equation}
This formulation separates small positive quantities rather than obtaining
them by subtracting two almost equal cosines.

\section{Existence, reconstruction, and ambiguous data}\label{e3:sec:reconstruction}

\begin{theorem}[Complete side-data criterion and reconstruction]\label{e3:thm:SSS}
For $a,b,c\in(0,\pi)$ the following are equivalent:
\begin{enumerate}
\item They are the angular sides of a proper spherical triangle.
\item The matrix $\mathcal G$ in \cref{e3:eq:grammatrix} is positive definite.
\item $D>0$.
\item $a<b+c$, $b<c+a$, $c<a+b$, and $a+b+c<2\pi$.
\end{enumerate}
When these conditions hold, a positively oriented realization is
\begin{equation}\label{e3:eq:SSSconstruction}
 \vv u=(1,0,0),\quad
 \vv v=(z,s_c,0),\quad
 \vv w=\left(y,\frac{x-yz}{s_c},\frac{\sqrt D}{s_c}\right).
\end{equation}
It is unique up to an orthogonal transformation; the positively oriented
realizations form one $\SO(3,F)$ orbit.
\end{theorem}
\begin{proof}
A genuine triple has $\mathcal G=P\transpose P$ with $P$ invertible, hence
positive definiteness. Conversely $1>0$ and $1-z^2>0$ are its first two
leading principal minors, so $D>0$ is exactly the remaining condition in
Sylvester's criterion. Alternatively \cref{e3:eq:SSSconstruction} directly
constructs a realization: the identity
$D=(1-y^2)(1-z^2)-(x-yz)^2$ makes the last vector unit and its dot products
are $y,x$. Its determinant is $\sqrt D>0$.

To prove the inequality criterion, use \cref{e3:eq:gramfactor}. Here
$0<s<3\pi/2$ and each $s-a,s-b,s-c$ belongs to $(-\pi/2,\pi)$.
At most one of these three numbers is nonpositive. If $s\ge\pi$, all three
are positive because $a,b,c<\pi$, whereas $\sin s\le0$, giving $D\le0$.
If $s<\pi$, positivity of $D$ is equivalent to all three differences being
positive. This is exactly the stated set of inequalities.

Finally, if $P$ and $P'$ are two realizations, then
$Q=P'P^{-1}$ satisfies $Q\transpose Q=I$ because their Gram matrices agree.
The sign of its determinant is the ratio of their orientation signs.
\end{proof}

The perimeter restriction is essential. The three inequalities alone are
not sufficient: sides all equal to $3\pi/4$ satisfy the ordinary triangle
inequalities but have perimeter exceeding $2\pi$.

\subsection{SAS, AAA, and ASA}

Any two sides $b,c\in(0,\pi)$ and included angle
$\alpha\in(0,\pi)$ determine exactly one proper triangle up to isometry.
An explicit realization is
\[
 \vv u=(1,0,0),\quad
 \vv v=(\cos c,\sin c,0),\quad
 \vv w=(\cos b,\sin b\cos\alpha,\sin b\sin\alpha).
\]
Its determinant is $\sin b\sin c\sin\alpha>0$. Compute $a$ from the dot
product and the remaining angles from \cref{e3:eq:angledata}. This avoids any
ambiguous inverse sine.

By polar duality, angles $\alpha,\beta,\gamma\in(0,\pi)$ occur exactly when
\begin{equation}\label{e3:eq:AAAcriterion}
 \alpha+\beta+\gamma>\pi,\qquad
 \beta+\gamma-\alpha<\pi,
\end{equation}
with the last inequality also imposed cyclically. They then determine the
triangle up to isometry. In particular, \emph{three angles fix the shape and
angular size} on a fixed sphere; there is no nontrivial similarity scaling of
angular sides. The condition that the angle sum exceed $\pi$ alone is not
sufficient. Two angles and their included side (ASA) reduce to SAS for the
polar triangle and likewise determine a unique proper triangle for inputs
in $(0,\pi)$.

\subsection{The spherical SSA ambiguity has an exceptional family}\label{e3:subsec:SSA}

Suppose that $a,b,\alpha\in(0,\pi)$ are given, with $a$ opposite
$\alpha$. Set
\[
 P=\cos b,\quad Q=\sin b\cos\alpha,\quad X=\cos a.
\]
The unknown side $c$ must solve
\begin{equation}\label{e3:eq:SSAlinecircle}
 P C+Q S=X,\qquad C^2+S^2=1,\qquad S>0,
 \quad (C,S)=(\cos c,\sin c).
\end{equation}
Every solution gives a proper triangle via the SAS construction, whose
opposite side is the prescribed $a$.

\begin{proposition}[Exact SSA alternatives]\label{e3:prop:SSA}
If $L=P^2+Q^2>0$, SSA data have zero, one, or two solutions up to isometry.
The candidate unit-circle pairs are
\begin{equation}\label{e3:eq:SSAcandidates}
 (C,S)=\frac{X}{L}(P,Q)
 \ \pm\ \frac{\sqrt{L-X^2}}{L}(-Q,P),
\end{equation}
provided $L-X^2\ge0$; only candidates with $S>0$ are retained, and duplicate
candidates are counted once.

If $P=Q=0$, there are no solutions unless $a=\pi/2$. If $a=\pi/2$,
then $b=\alpha=\pi/2$ as well, and \emph{every} $c\in(0,\pi)$ gives a
solution, with $\beta=\pi/2$ and $\gamma=c$.
\end{proposition}
\begin{proof}
For $L>0$, intersect the line with the unit circle by decomposing a point into
components parallel and perpendicular to $(P,Q)$. This gives
\cref{e3:eq:SSAcandidates}; a line has at most two intersections. The upper-half
condition is exactly $c\in(0,\pi)$.
If $P=Q=0$, the original equation is $0=X$. Since $b\in(0,\pi)$ and
$\sin b>0$, the vanishing of $P,Q$ means $b=\alpha=\pi/2$.
When also $a=\pi/2$, take
$\vv u=(1,0,0)$, $\vv v=(\cos c,\sin c,0)$, $\vv w=(0,0,1)$.
Their determinant is $\sin c>0$. The angle formulas give the claimed
remaining angles.
\end{proof}

Thus a blanket assertion that spherical SSA always has at most two solutions
misses a genuine continuous exceptional family. The polar construction gives
the corresponding AAS ambiguity. The effect is already present over $\R$;
surreal coordinates add arbitrarily small perturbations of the exceptional
data, for which branch-sensitive computation is especially important.

\subsection{Latitude, longitude, and bearings}

For latitude $\varphi\in[-\pi/2,\pi/2]$ and longitude $\lambda$ modulo
$2\pi$, the unit direction is
\[
 (\cos\varphi\cos\lambda,\cos\varphi\sin\lambda,\sin\varphi).
\]
Two such points have separation $a$ given by
\begin{align}
 \cos a&=\sin\varphi_1\sin\varphi_2+
          \cos\varphi_1\cos\varphi_2\cos(\lambda_2-\lambda_1),\label{e3:eq:latcos}\\
 \hav a&=\hav(\varphi_2-\varphi_1)+
          \cos\varphi_1\cos\varphi_2\hav(\lambda_2-\lambda_1).\label{e3:eq:lathav}
\end{align}
At a nonpolar starting point and for a distinct nonantipodal endpoint, the
initial bearing measured from north toward east is
\begin{equation}\label{e3:eq:bearing}
 \atantwo\left(
 \cos\varphi_2\sin\Delta\lambda,
 \cos\varphi_1\sin\varphi_2-
 \sin\varphi_1\cos\varphi_2\cos\Delta\lambda\right).
\end{equation}
Indeed these are the east and north components of the projected tangent to
the target. Undefined bearings at poles, coincident points, or antipodes are
geometric ambiguities, not numerical failures. The physical great-circle
distance remains $Ra$ for any positive surreal radius.

\section{Spherical centers, medians, and Ceva}\label{e3:sec:centers}

Assume $\tau=d>0$ and retain the matrix $P$ and inward normals of
\cref{e3:eq:polarnormals}.

\subsection{Circumcenter and incenter}

Put $\vv h=P^{-\mathsf T}(1,1,1)\transpose$. Then
\begin{equation}\label{e3:eq:circumcenter}
 \vv o=\frac{\vv h}{\norm{\vv h}},\qquad
 \rho_c=\acos\frac1{\norm{\vv h}}
\end{equation}
is the center and angular radius of the circumcircle with radius less than
$\pi/2$. All three dot products with $\vv o$ equal
$1/\norm{\vv h}>0$. Cauchy--Schwarz makes this number at most one, and
distinct vertices make it strictly less than one. Every equidistant unit
center solves $P\transpose\vv o=k(1,1,1)\transpose$, so the two possibilities
are exactly $\vv o$ and $-\vv o$, with radii $\rho_c$ and $\pi-\rho_c$.
The chosen circumcenter need not lie inside the triangle.

The incenter has the particularly simple weighted-vector expression
\begin{equation}\label{e3:eq:incenter}
 \vv j=\frac{s_a\vv u+s_b\vv v+s_c\vv w}
               {\norm{s_a\vv u+s_b\vv v+s_c\vv w}},\qquad
 r_i=\asin\frac{d}{\norm{s_a\vv u+s_b\vv v+s_c\vv w}}.
\end{equation}
To prove this, take dot products of its numerator with the three inward
normals. Each equals $d$. Hence $\vv j$ has equal positive signed distances
from the three side great circles, and it lies inside the positive cone.
Those signed angular distances are the inverse sines in
\cref{e3:eq:incenter}. Uniqueness follows by solving the nonsingular linear
system $N\transpose\vv j=k(1,1,1)\transpose$ for the inward normal matrix
$N$. Equivalently, the inside angle bisectors below concur at this point.
The physical circumradius and inradius along the sphere are $R\rho_c$ and
$Rr_i$, not $\rho_c$ and $r_i$ themselves.

\subsection{Medians and angle bisectors}

The midpoint of a minor edge from $\vv v$ to $\vv w$ is
$(\vv v+\vv w)/\norm{\vv v+\vv w}$. The three vertex-to-opposite-midpoint
great circles concur at
\[
 \frac{\vv u+\vv v+\vv w}{\norm{\vv u+\vv v+\vv w}}.
\]
This follows by writing the common numerator as a vertex plus the sum of
the other two. It is a median center, not an assertion that it is the
surface-area centroid.

At $\vv u$, an internal angle bisector has tangent direction
$T_{uv}+T_{uw}$. Its opposite-side intersection is proportional to
$s_b\vv v+s_c\vv w$. If that point is $D_a$ on the edge $vw$, then
\begin{equation}\label{e3:eq:spherebisector}
 \frac{\sin d_S(\vv v,D_a)}{\sin d_S(D_a,\vv w)}
       =\frac{\sin c}{\sin b}.
\end{equation}
This is a ratio of sines of angular subarcs, not in general a ratio of their
angular lengths.

\begin{theorem}[Spherical Ceva]\label{e3:thm:ceva}
For interior points $D_a,D_b,D_c$ on the minor edges $vw,wu,uv$, the three
cevians concur inside the triangle if and only if
\begin{equation}\label{e3:eq:sphericalceva}
 \frac{\sin d_S(\vv v,D_a)}{\sin d_S(D_a,\vv w)}
 \frac{\sin d_S(\vv w,D_b)}{\sin d_S(D_b,\vv u)}
 \frac{\sin d_S(\vv u,D_c)}{\sin d_S(D_c,\vv v)}=1.
\end{equation}
\end{theorem}
\begin{proof}
Write $D_a$ as the normalization of $\lambda\vv v+\mu\vv w$ with
$\lambda,\mu>0$. Cross products with the endpoints show that its first sine
ratio is $\mu/\lambda$. Do the same on the other two edges.
A positive linear combination $A\vv u+B\vv v+C\vv w$ belongs to the
first cevian exactly when $C/B=\mu/\lambda$. The cyclic conditions have a
positive solution precisely when their product is one. Normalizing that
solution gives the interior concurrence point. This is ordinary barycentric
Ceva expressed in spherical sine ratios.
\end{proof}

\section{Spherical excess, solid angles, and polygon area}\label{e3:sec:area}

Define
\begin{equation}\label{e3:eq:Mdef}
 M=1+x+y+z,\qquad B_0=\sqrt{M^2+D}.
\end{equation}
An important positive identity is
\begin{equation}\label{e3:eq:Midentity}
 M^2+D=2(1+x)(1+y)(1+z)>0.
\end{equation}
The sign of $M$ must not be discarded. It determines whether the triangle's
area is smaller or larger than a quarter of the full sphere.

\begin{theorem}[Branch-correct spherical excess]\label{e3:thm:excess}
For a proper spherical triangle,
\begin{equation}\label{e3:eq:excess}
 E:=\alpha+\beta+\gamma-\pi
   =2\atantwo(d,M),\qquad 0<E<2\pi.
\end{equation}
In particular, $E<\pi$, $E=\pi$, or $E>\pi$ according as $M>0$, $M=0$, or
$M<0$. The associated half-angle data are
\begin{equation}\label{e3:eq:excessphase}
 \cos(E/2)=\frac{M}{B_0},\qquad
 \sin(E/2)=\frac{d}{B_0}.
\end{equation}
\end{theorem}
\begin{proof}
The polar triangle has perimeter less than $2\pi$ by \cref{e3:thm:SSS}, so
$3\pi-(\alpha+\beta+\gamma)<2\pi$. Each interior angle is less than
$\pi$. Hence $0<\alpha+\beta+\gamma-\pi<2\pi$.

For an algebraic evaluation, work in $F[i]$, and put
$A_1=x-yz$, $A_2=y-zx$, $A_3=z-xy$ and
$T=(1-x)(1-y)(1-z)$. Direct multiplication, using
$d^2=1+2xyz-x^2-y^2-z^2$, gives
\begin{equation}\label{e3:eq:excesspoly}
 2(A_1+id)(A_2+id)(A_3+id)+T(M+id)^2=0.
\end{equation}
This can be checked by separately comparing real and imaginary parts;
they reduce respectively to
\[
 2\bigl(A_1A_2A_3-D(A_1+A_2+A_3)\bigr)+T(M^2-D)=0
\]
and $A_1A_2+A_2A_3+A_3A_1-D+TM=0$.
By \cref{e3:eq:angledata}, the product of the three angle phases is the product
of the first three factors in \cref{e3:eq:excesspoly}, divided by
$(1-x^2)(1-y^2)(1-z^2)$. Since that denominator equals
$T(M^2+D)/2$, we obtain
\[
 \cis(\alpha+\beta+\gamma)=-\frac{(M+id)^2}{M^2+D}.
\]
It follows that $E$ and $2\atantwo(d,M)$ have the same phase. Both lie in
$(0,2\pi)$, where phase is injective, so they are equal. The sign alternatives
and \cref{e3:eq:excessphase} follow immediately.
\end{proof}

When $M>0$, \cref{e3:eq:excess} can be written $2\atan(d/M)$. When $M<0$,
the same uncorrected one-argument arctangent would give a negative number
instead of the correct value in $(\pi,2\pi)$. This is a geometric branch error,
not merely a floating-point issue.

\subsection{L'Huilier's formula}

For $s=(a+b+c)/2$, all four half-arguments below lie in $(0,\pi/2)$.
We have the exact formula
\begin{equation}\label{e3:eq:lhuilier}
 \tan(E/4)=\sqrt{
 \tan(s/2)\tan((s-a)/2)\tan((s-b)/2)\tan((s-c)/2)}.
\end{equation}
Here the root is positive, even for infinitesimal excess or for excess
infinitesimally close to $2\pi$.

To prove it, use $\tan(E/4)=d/(B_0+M)$ from
\cref{e3:eq:excessphase}. The product-to-sum identity
\[
 B_0+M=8\cos(s/2)\cos((s-a)/2)
               \cos((s-b)/2)\cos((s-c)/2)
\]
follows by expansion, using
$B_0=4\cos(a/2)\cos(b/2)\cos(c/2)>0$ from
\cref{e3:eq:Midentity}. Replace the four full-angle sines in
\cref{e3:eq:gramfactor} by twice the corresponding sine--cosine products.
After division by $(B_0+M)^2$, one obtains the square of
\cref{e3:eq:lhuilier}. Positivity supplies the indicated root.

\subsection{Area as a precise finite-polygon functional}

\begin{definition}[Spherical polygon area]\label{e3:def:area}
For a proper triangle on $S_R^2(C)$ define its area to be $R^2E$.
For a convex geodesic polygon contained in an open hemisphere, triangulate
it by minor great-circle diagonals and add these triangle areas.
\end{definition}

\begin{theorem}[Consistency, positivity, and finite additivity]\label{e3:thm:polygonarea}
The definition is independent of triangulation. For a convex polygon with
$n$ boundary vertices and interior angles $\alpha_1,\ldots,\alpha_n$,
\begin{equation}\label{e3:eq:polygonexcess}
 \area(P)=R^2\left(\sum_{j=1}^n\alpha_j-(n-2)\pi\right).
\end{equation}
It is positive for a nonempty two-dimensional polygon, invariant under
orthogonal transformations, and additive under finite geodesic polygonal
dissections with disjoint interiors, with edges and vertices assigned area zero. It extends by such dissections to finite
unions of spherical geodesic polygons.
\end{theorem}
\begin{proof}
Use a gnomonic chart to make the polygon an ordinary convex polygon.
In a triangulation with boundary vertices only, summing all the triangle
angles gives precisely the sum of boundary angles, and there are $n-2$
triangles. This proves \cref{e3:eq:polygonexcess} and independence of the
triangulation.

A refinement can introduce interior vertices or subdivide boundary edges.
At an interior vertex the angles sum to $2\pi$; at an inserted straight
boundary vertex the angle is $\pi$. If a planar triangulation has $I$
interior vertices, $B$ boundary vertices and $T_\triangle$ triangles, the
finite Euler count is $T_\triangle=2I+B-2$. Thus the extra angles cancel
exactly the extra triangle subtractions. A common finite refinement of two
dissections proves additivity and consistency. Positivity follows by adding
the positive triangle excesses. Orthogonal invariance follows from invariance
of dot products and interior angles. For regions not contained in one open
hemisphere, intersect with the eight closed coordinate octants, subdividing
edges where necessary. Each octant lies in the open hemisphere whose normal
has its three coordinate signs. Thus all the resulting pieces admit the
preceding construction, and common refinement proves the same consistency.
No topological compactness is needed for this explicit finite subdivision.
\end{proof}

An octant triangle has three right angles, hence excess $\pi/2$.
Eight octants therefore give
\begin{equation}\label{e3:eq:spherearea}
 \area(S_R^2)=4\pi R^2,\qquad
 \area(\text{hemisphere})=2\pi R^2.
\end{equation}
A lune of angle $\lambda\in(0,\pi)$ splits along the equator into two
triangles with angles $\lambda,\pi/2,\pi/2$. Its area is $2R^2\lambda$.
These values reproduce the usual geometric normalization of area.

For a circular cap of angular radius $\theta$, the usual elementary
chartwise primitive of the area density gives the canonical expression
\[
 A_{\mathrm{cap}}(\theta)=2\pi R^2(1-\cos\theta),\qquad0\le\theta\le\pi.
\]
This is the extension of a specific classical analytic formula, obtained by
evaluating the primitive $-\cos\theta$ of $\sin\theta$. It is not an
assertion that arbitrary subsets have a surreal measure or that ordinary
Darboux sums define this value. Our finite-additivity theorem concerns the
geodesic polygonal class just constructed.

\subsection{Solid angle of three arbitrary rays}

Let $\vv A,\vv B,\vv C\in F^3$ be linearly independent vectors from an
observation point, with lengths $A,B,C>0$. Their spherical triangle has
solid angle
\begin{equation}\label{e3:eq:solidangle}
 \Omega=2\atantwo\left(
 |\det[\vv A\ \vv B\ \vv C]|,
 ABC+C(\vv A\cdot\vv B)+A(\vv B\cdot\vv C)+B(\vv C\cdot\vv A)
 \right).
\end{equation}
Normalize each ray to a unit vector in \cref{e3:eq:excess}; multiplying both
arguments by the positive number $ABC$ proves the formula. This is the
branch-complete form of the classical plane-triangle solid-angle formula
associated with van Oosterom and Strackee \cite{Oosterom}. The formula is
independent of positive rescaling of any individual ray. Replacing the
absolute determinant by its signed value gives the oriented angle, with
sign set by the ray order.

Solid angle is dimensionless and finite, whereas the corresponding area
$R^2\Omega$ can have any surreal scale. A determinant that is infinitesimal
relative to ray lengths does not by itself imply small solid angle: when the
second argument is negative, the angle can instead be close to $2\pi$.

\section{Differential geometry and parallel transport}\label{e3:sec:differential}

All derivatives in this section are with respect to explicit chart or curve
parameters. Scalar parameters such as $R$ are held fixed. These derivatives
are not an intrinsic derivation of the scalar field $\No$. The calculations
are identities between restricted-analytic germs and rational operations;
no existence theorem for arbitrary differential equations is being invoked.

\subsection{Connection, geodesics, and curvature}

For a sphere centered at zero, its tangent plane at $P$ is
$T_PS_R^2=\{X:P\cdot X=0\}$. Orthogonally project ambient derivatives to
obtain the connection
\begin{equation}\label{e3:eq:connection}
 \nabla_XY=D_XY+\frac{X\cdot Y}{R^2}P.
\end{equation}
Indeed differentiating $P\cdot Y=0$ gives
$P\cdot D_XY=-X\cdot Y$, so the displayed correction makes the result
tangent. The product rule shows metric compatibility; symmetry of the
ambient derivative gives zero torsion.

For a unit tangent vector $V$ at $P$, a unit-speed great-circle geodesic is
\begin{equation}\label{e3:eq:geodesic}
 \gamma(s)=\cos(s/R)P+R\sin(s/R)V,
 \qquad s/R\in\OO.
\end{equation}
Direct differentiation gives $\norm{\gamma'}=1$ and
$\gamma''=-\gamma/R^2$, hence $\nabla_{\gamma'}\gamma'=0$.
All minor geodesic segments have $0\le s/R\le\pi$, so they are included.
No oscillatory value at an infinitely large $s/R$ is required.

With curvature convention
$\mathcal R(X,Y)Z=\nabla_X\nabla_YZ-\nabla_Y\nabla_XZ-\nabla_{[X,Y]}Z$,
substitution of \cref{e3:eq:connection} and cancellation of commuting ambient
derivatives gives
\begin{equation}\label{e3:eq:curvature}
 \mathcal R(X,Y)Z=\frac1{R^2}
       \bigl((Y\cdot Z)X-(X\cdot Z)Y\bigr).
\end{equation}
Thus the Gaussian curvature is exactly $1/R^2$. In particular a sphere of
infinite radius has positive infinitesimal curvature, not zero curvature.
An infinitesimal sphere has infinite curvature. Both statements are literal
field identities.

The exponential map at $P$, for $X\in T_PS_R^2$ with $r=\norm X$ and
$r/R$ finite, is
\[
 \operatorname{Exp}_P(X)=\cos(r/R)P+R\sin(r/R)\frac{X}{r},
\]
with value $P$ at $r=0$. For $0\le r<\pi R$, it is bijective onto the sphere
minus the antipode. Its inverse for $Q\ne-P$ is
\begin{equation}\label{e3:eq:logmap}
 \operatorname{Log}_P(Q)=\frac{a}{\sin a}(Q-\cos a\,P),
 \qquad a=\acos(P\cdot Q/R^2),
\end{equation}
with the removable value zero at $Q=P$. Taking norms and using the unique
minor arc verifies both identities.

\subsection{Explicit parallel transport}

For unit directions $\vv u\ne-\vv v$ and $X\cdot\vv u=0$, define
\begin{equation}\label{e3:eq:transport}
 \mathcal P_{uv}(X)=X-
       \frac{X\cdot\vv v}{1+\vv u\cdot\vv v}(\vv u+\vv v).
\end{equation}
This is the parallel transport along the minor great-circle arc. To verify
it, choose the tangent to the arc and its constant binormal as a basis of the
initial tangent plane. The tangent rotates with the great circle and has
zero covariant derivative; the binormal is ambient-constant and also has
zero covariant derivative. The endpoint values of these two parallel fields
are exactly \cref{e3:eq:transport}. Alternatively, restricting
\cref{e3:eq:minimalrotation} to $\vv u^\perp$ reduces algebraically to the
same formula. It follows that transport is an isometry between tangent planes.
This specifies the canonical transport by explicit analytic continuation
along the parametrized arc. It does not assert uniqueness among arbitrary
intrinsically differentiable fields: the locally constant counterexample in
\cref{e3:sec:boundaries} explains that distinction. Within the definable
restricted-analytic class, uniqueness follows because the two coefficients
in the parallel frame have zero derivative and hence are constant.
For a radius-$R$ sphere the unit directions are used but tangent vectors may
retain their physical lengths.

\subsection{A quaternion proof of spherical holonomy}

Use the positive orientation $\tau=d>0$. The unit quaternion of the minimal
rotation from $\vv u$ to $\vv v$ is
\begin{equation}\label{e3:eq:edgequaternion}
 q_{uv}=\frac{(1+z,\vv u\times\vv v)}{\sqrt{2(1+z)}}.
\end{equation}
Define $q_{vw}$ and $q_{wu}$ cyclically. Quaternion multiplication gives the
especially informative identity
\begin{equation}\label{e3:eq:holonomyquat}
 q_{wu}q_{vw}q_{uv}=
       \frac{(M,d\vv u)}{\sqrt{M^2+D}}.
\end{equation}
For a direct verification, multiply the three numerators using
\cref{e3:eq:quaternionmul}. The scalar component is $2M$ and the vector
component is $2d\vv u$; the denominator is
$\sqrt{8(1+x)(1+y)(1+z)}=2\sqrt{M^2+D}$.
The simplifications use only unit norms and the vector triple-product
identities.

\begin{theorem}[Holonomy equals spherical excess]\label{e3:thm:holonomy}
Parallel transport around the positively oriented boundary
$\vv u\to\vv v\to\vv w\to\vv u$ is rotation of the initial tangent
plane through angle $E$ modulo $2\pi$. Equivalently,
\[
 \text{holonomy phase}=\cis E,
 \qquad E=\frac{\area(\triangle)}{R^2}.
\]
\end{theorem}
\begin{proof}
Each edge transport is the restriction of the corresponding quaternion
rotation. By \cref{e3:eq:holonomyquat,e3:eq:excessphase}, their product is the
quaternion $(\cos(E/2),\sin(E/2)\vv u)$, which acts on
$\vv u^\perp$ by rotation through $E$. The area equality is
\cref{e3:def:area}.
\end{proof}

This is a finite-triangle version of the curvature--area holonomy relation,
proved without an unconstructed global integral. The observed tangent-plane
rotation records $E$ modulo $2\pi$, while the area branch has already been
fixed by $0<E<2\pi$. Reversing the boundary orientation changes the holonomy
sign but not the unoriented area.

\subsection{Curvature correction in normal coordinates}\label{e3:subsec:normal-limit}

Let $p,q$ be vectors in the tangent plane at $P$, with
$r=\norm p$, $s=\norm q$, and $(r+s)/R$ infinitesimal. The squared intrinsic
distance between their exponential images satisfies
\begin{equation}\label{e3:eq:normal-limit}
 \begin{split}
 d_R(\operatorname{Exp}_P p,\operatorname{Exp}_P q)^2
  ={}&\norm{p-q}^2
  -\frac{r^2s^2-(p\cdot q)^2}{3R^2}\\
  &+\smallO\left(\frac{(r+s)^6}{R^4}\right).
 \end{split}
\end{equation}
The error estimate is absolute and uniform over the relative directions;
it does not claim a small \emph{relative} error when $p-q$ is exceptionally
small.

To derive it, for $r,s\ne0$ use the spherical cosine law with
$\cos\alpha=(p\cdot q)/(rs)$. Expanding to degree four in $r/R,s/R$ gives
\[
 \cos(d_R/R)=1-\frac{\norm{p-q}^2}{2R^2}
 +\frac{r^4+s^4+6r^2s^2-4(p\cdot q)(r^2+s^2)}{24R^4}
 +\smallO\bigl((r+s)^6/R^6\bigr).
\]
The function $\acos(1-h)^2=2h+h^2/3+\smallO(h^3)$ has a real analytic
extension through $h=0$. Substitution yields \cref{e3:eq:normal-limit}. Expressions with
$r=0$ or $s=0$ are obtained by their removable analytic values. Uniform
real-germ remainder estimates transfer after rescaling by $R$.
The coefficient $-1/(3R^2)$ explicitly records positive curvature.

\section{Infinitesimal scales and geometric conditioning}\label{e3:sec:scales}

The extra information carried by surreal coordinates is not a new algebraic
identity for an ordinary triangle. It is the ability to retain exact rank,
orientation, angle, and size when several different infinitesimal scales
occur in the same configuration.

\subsection{A valuation dictionary}

Write $v$ for the Hahn valuation, with $v(0)=+\infty$. For a nonzero
vector $X\in F^3$,
\begin{equation}\label{e3:eq:norm-valuation}
 v(\norm X)=\min_j v(X_j).
\end{equation}
Indeed, in $\sum_jX_j^2$ the summands of smallest valuation have positive
leading coefficients, so their leading terms cannot cancel. Taking a
positive square root halves the resulting valuation.

For unit vectors $u,v$, put $a=d_S(u,v)$. The chord bounds and their
antipodal versions give
\begin{align}
 v(\norm{u-v})&=v(a) &&(a>0),\label{e3:eq:chord-val}\\
 v(\norm{u+v})&=v(\pi-a) &&(a<\pi).
\end{align}
For $0<a<\pi$ one also has
\begin{equation}\label{e3:eq:sin-val}
 v(\sin a)=v(a)+v(\pi-a).
\end{equation}
To check the last assertion, when both endpoints are separated from $a$
by positive real amounts all three valuations are zero. When $a$ is
infinitesimal, $\sin a/a=1+\smallO(a^2)$ and $\pi-a$ is a unit; the
case $\pi-a$ infinitesimal is identical after reflection. These exhaust
the possibilities by standard parts.

Consequently a proper spherical triangle satisfies the exact scale laws
\begin{equation}\label{e3:eq:triangle-val}
 \begin{split}
 v(d)&=v(\sin b)+v(\sin c)+v(\sin\alpha)\\
 &=v(\sin c)+v(\sin a)+v(\sin\beta)\\
 &=v(\sin a)+v(\sin b)+v(\sin\gamma).
 \end{split}
\end{equation}
A small determinant can therefore come from a small side, a nearly
antipodal side, a small interior angle, or an angle close to $\pi$.
Neither its standard part nor its valuation alone identifies which
geometric mechanism is responsible.

\subsection{A shape-independent local area estimate}\label{e3:subsec:local-area}

The following estimate does not assume that a shrinking triangle has
angles bounded away from zero. This matters over the surreals: the area
of a triangle can be smaller than every ordinary power of its diameter.

\begin{theorem}[Uniform gnomonic area comparison]\label{e3:thm:local-area}
Let $p_0,p_1,p_2\in F^2$ be noncollinear, let
$\rho=\max_i\norm{p_i}$, and suppose $0<\rho\le1/4$. On the unit
sphere take the vertices $u_i=(p_i,1)/\sqrt{1+\norm{p_i}^2}$.
Set
\[
 H=\det[p_1-p_0\ \ p_2-p_0],\qquad A_{\rm flat}=\frac{|H|}{2}.
\]
The spherical excess $E$ of their minor-arc triangle satisfies
\begin{equation}\label{e3:eq:local-area-bound}
 \left|\frac{E}{A_{\rm flat}}-1\right|\le4\rho^2.
\end{equation}
In particular, for infinitesimal $\rho$,
\[
 E=\frac{|H|}{2}(1+\smallO(\rho^2)),\qquad
 v(E)=v(H).
\]
On the sphere of radius $R>0$ the spherical area is $R^2E$, so
$v(\area)=2v(R)+v(H)$, without any restriction on the size of $R$.
\end{theorem}

\begin{proof}
Put $A_i=\sqrt{1+\norm{p_i}^2}$ and
\[
 B=A_0A_1A_2+(1+p_0\cdot p_1)A_2
 +(1+p_1\cdot p_2)A_0+(1+p_2\cdot p_0)A_1.
\]
The determinant and denominator in the solid-angle formula are
$d=|H|/(A_0A_1A_2)$ and $M=B/(A_0A_1A_2)$, respectively. Thus
\begin{equation}\label{e3:eq:gnomonic-exact-area}
 E=2\atan(|H|/B).
\end{equation}
Here $B>0$; the triangle lies near the north pole, not near the
hemispherical branch cut.

The inequalities $1\le A_i\le1+\rho^2/2$ and
$|p_i\cdot p_j|\le\rho^2$ imply
\[
 |B-4|\le8\rho^2,\qquad B\ge7/2,\qquad |H|\le4\rho^2.
\]
For example, the deviation of the product $A_0A_1A_2$ from $1$ is at
most $2\rho^2$, and each of the other three terms deviates from $1$
by at most $2\rho^2$. For the lower bound one may use the sharper
$B\ge4-3\rho^2$.

Set $\zeta=|H|/B$, so $0<\zeta\le2\rho^2$. The ordinary analytic
inequality $1-\zeta^2/3\le\atan\zeta/\zeta\le1$ holds here by
restricted-analytic transfer. Therefore
\[
 \begin{split}
 \left|\frac{E}{A_{\rm flat}}-1\right|
 &=\left|\frac4B\frac{\atan\zeta}{\zeta}-1\right|\\
 &\le\frac{16}{7}\rho^2+\frac{32}{21}\rho^4
 \le4\rho^2.
 \end{split}
\]
Crucially, $H$ was divided out only after the exact factorization of
$E$; no lower bound on $|H|/\rho^2$ was used. The valuation claims
follow because the final multiplicative factor has standard part $1$.
\end{proof}

\begin{remark}
For $p_0=0$, $p_1=(\eps,0)$, and
$p_2=(\eps,\omega^{-\omega})$ with $\eps=\omega^{-1}$, the
ratio $|H|/\rho^2$ is smaller than every ordinary positive power of
$\eps$. The theorem still gives a relative, rather than merely
absolute, estimate for the spherical area. A generic error bound
$E=A_{\rm flat}+\smallO(\rho^4)$ would not give that information.
\end{remark}

\subsection{Stable vector data and unstable derived data}

\begin{proposition}[Three conditioning estimates]\label{e3:prop:conditioning}
For unit vectors $u,v,u',v'$,
\begin{equation}\label{e3:eq:angle-Lipschitz}
 |d_S(u,v)-d_S(u',v')|
 \le\frac\pi2\bigl(\norm{u-u'}+\norm{v-v'}\bigr).
\end{equation}
For nonzero $X,Y$,
\begin{equation}\label{e3:eq:normalize-bound}
 \left\|\frac X{\norm X}-\frac Y{\norm Y}\right\|
 \le\frac{2\norm{X-Y}}{\norm X}.
\end{equation}
Finally, let $a=d_S(u,v)\in(0,\pi)$ and
$\eta=\norm{u-u'}+\norm{v-v'}<\sin a$. Then $u'\times v'\ne0$ and
\begin{equation}\label{e3:eq:normal-bound}
 \left\|\frac{u\times v}{\norm{u\times v}}
       -\frac{u'\times v'}{\norm{u'\times v'}}\right\|
 \le\frac{2\eta}{\sin a}.
\end{equation}
\end{proposition}

\begin{proof}
The angular triangle inequality bounds the left side of
\cref{e3:eq:angle-Lipschitz} by
$d_S(u,u')+d_S(v,v')$. Apply the lower chord bound to each summand.
For \cref{e3:eq:normalize-bound}, write the difference as
\[
 \frac{X-Y}{\norm X}
 +Y\left(\frac1{\norm X}-\frac1{\norm Y}\right)
\]
and use $|\norm X-\norm Y|\le\norm{X-Y}$. For the last estimate,
\[
 \norm{u\times v-u'\times v'}
 \le\norm{(u-u')\times v}+\norm{u'\times(v-v')}
 \le\eta.
\]
The strict inequality $\eta<\norm{u\times v}=\sin a$ prevents the
new cross product from vanishing; normalization gives the claimed bound.
\end{proof}

Thus angular distance is uniformly stable under consistent perturbations
of the unit vectors. The normal to their spanning plane is not uniformly
stable near coincident or antipodal directions: the amplification factor
$1/\sin a$ is then infinite. There is no contradiction between these
statements; they concern different outputs.

If one perturbs only the scalar cosine as though it were independent
of all other data, inverse cosine is also ill-conditioned near an endpoint.
Away from the endpoints its derivative is
\[
 \frac{d}{dc}\acos c=-\frac1{\sqrt{1-c^2}},
\]
whereas for $h>0$ infinitesimal,
\begin{equation}\label{e3:eq:acos-endpoint}
 \acos(1-h)=\sqrt{2h}\left(1+\frac h{12}+\smallO(h^2)\right).
\end{equation}
The square-root law follows by inverting
$1-\cos\theta=\theta^2/2-\theta^4/24+\cdots$ after extracting
$\sqrt{2h}$. It is a genuinely different scale from an ordinary
linear perturbation estimate.

\section{Worked configurations at several surreal scales}\label{e3:sec:examples}

\subsection{Infinite lengths with infinitesimal directions}

The vectors $X=(\omega,1,0)$ and $Y=(\omega,0,0)$ have infinite
lengths and the finite infinitesimal angle
\[
 \ang(X,Y)=\atan(\omega^{-1})
 =\omega^{-1}-\frac13\omega^{-3}+\frac15\omega^{-5}-\cdots.
\]
This displayed series is strongly summable because its parameter is
infinitesimal. No evaluation of $\sin\omega$ occurs.

For the flat triangle
\[
 P=0,\qquad Q=(\omega,0,0),\qquad
 S=(\omega,\omega^{-\omega},0),
\]
the area is exactly $\tfrac12\omega^{1-\omega}$, while the angle at
$P$ is $\atan(\omega^{-\omega-1})$. The base is infinite, the
altitude is smaller than every ordinary negative power of $\omega$,
and the area is nevertheless a well-defined positive surreal number.
A real standard-part picture alone does not encode this configuration.

\subsection{Nearly parallel lines meet their closest points infinitely far away}

Let $\eps>0$ be infinitesimal and let $h=\eps^2$. Consider
\[
 L_1=\{(t,0,0):t\in F\},\qquad
 L_2=\{(s,1+\eps s,h):s\in F\}.
\]
Their squared distance between parameter points is
\[
 (t-s)^2+(1+\eps s)^2+h^2.
\]
Its minimum is attained at $t=s=-1/\eps$ and is $h^2$. Hence the
lines have distance $h$, although the standard parts of their
\emph{finite-coordinate} points lie on the real lines $y=0,z=0$
and $y=1,z=0$, whose distance is $1$.

The discrepancy is not a failure of the distance formula. The minimizing
points have infinite first coordinates and disappear from a finite-coordinate
standard-part picture. Taking standard parts and minimizing over an entire
affine line are different operations. With $h=0$, the same lines intersect
at a genuine affine point infinitely far from the origin.

\subsection{Tangency splits at a square-root scale}

Intersect the unit sphere with the line $(t,1-\eps,0)$, where
$0<\eps\in\mm$. The exact solutions are
\[
 t=\pm\sqrt{2\eps-\eps^2}
 =\pm\sqrt{2\eps}\left(1-\frac\eps4+\smallO(\eps^2)\right).
\]
A displacement of valuation $v(\eps)$ splits a double tangency into
points separated at valuation $v(\eps)/2$. Replacing $1-\eps$ by
$1+\eps$ produces no intersection, while $\eps=0$ gives a single
point. All three cases have the same ordinary standard-part tangency.

\subsection{An almost planar tetrahedron with an infinite circumradius}

Take vertices $0,e_1,e_2,(2,0,\eps)$ with $0<\eps\in\mm$.
The tetrahedron has volume $\eps/6$. Its circumcenter is
\[
 C=\left(\frac12,\frac12,\frac1\eps+\frac\eps2\right),
\]
since $2C\cdot e_i=1$ for $i=1,2$ and
$2C\cdot(2,0,\eps)=4+\eps^2$. The circumradius is
\[
 \norm C=\sqrt{\frac12+\left(\frac1\eps+\frac\eps2\right)^2},
\]
which is infinite. Nonzero volume, however small, is enough for uniqueness
of the center; it is not enough for a finite conditioning constant.

\subsection{A right spherical triangle and the flat limit of an infinite sphere}

For positive infinitesimals $\eps,\eta$, set
\[
 u=e_3,\qquad
 v=\frac{(\eps,0,1)}{A},\qquad
 w=\frac{(0,\eta,1)}{B},\qquad
 A=\sqrt{1+\eps^2},\quad B=\sqrt{1+\eta^2}.
\]
The angle at $u$ is exactly $\pi/2$, the adjacent side angles are
$c=\atan\eps$ and $b=\atan\eta$, and
\[
 x=\frac1{AB},\qquad y=\frac1B,\qquad z=\frac1A,
 \qquad d=\frac{\eps\eta}{AB}.
\]
The exact excess is therefore
\begin{equation}\label{e3:eq:right-small-excess}
 E=2\atan\frac{\eps\eta}{(A+1)(B+1)}
 =\frac{\eps\eta}{2}-\frac{\eps^3\eta+\eps\eta^3}{8}
 +\smallO\bigl(\eps\eta(\eps^2+\eta^2)^2\bigr).
\end{equation}
This is a relative expansion with an explicit common area factor.
In the symmetric case $\eps=\eta=h$ it becomes
\begin{equation}\label{e3:eq:symmetric-area}
 E=\frac{h^2}{2}-\frac{h^4}{4}+\frac{7h^6}{48}+\smallO(h^8).
\end{equation}

Now choose $R=\omega$ and $h=\omega^{-1}$, and scale all vertices
by $R$. The physical area is
\[
 R^2E=\frac12-\frac1{4\omega^2}
       +\frac7{48\omega^4}+\smallO(\omega^{-6}),
\]
and the physical adjacent side lengths are
$\omega\atan(\omega^{-1})=1-1/(3\omega^2)+\cdots$.
The sphere has positive infinitesimal curvature $\omega^{-2}$.
A bounded patch is asymptotically planar, but its curvature and area
corrections remain exact surreal information rather than being discarded.

\subsection{A tiny determinant with almost hemispherical area}

Let $0<\eps\in\mm$ and consider
\[
 u=(1,0,0),\qquad v=(-1/2,\sqrt3/2,0),\qquad
 w=\frac{(-1/2,-\sqrt3/2,\eps)}{\sqrt{1+\eps^2}}.
\]
Writing $A=\sqrt{1+\eps^2}$, one obtains
\[
 d=\frac{\sqrt3\,\eps}{2A},\qquad M=\frac12-\frac1A<0.
\]
Thus $D=d^2$ is infinitesimal, but the correct branch gives
\begin{equation}\label{e3:eq:hemisphere-example}
 E=2\pi-2\atan\frac{d}{-M}
   =2\pi-2\sqrt3\,\eps+\smallO(\eps^3).
\end{equation}
The triangle nearly fills a hemisphere rather than collapsing to zero
area. Its polar triangle is small. Replacing $\atantwo(d,M)$ by the
principal $\atan(d/M)$ would give the wrong branch and even the wrong
sign for the area. This example and \cref{e3:eq:right-small-excess}
show why $d\to0$ is not a complete description of spherical degeneration.

\section{What the development does not silently assume}\label{e3:sec:boundaries}

\subsection{The intrinsic order topology is not ordinary real topology}

Suppose $F$ is non-Archimedean. Its infinitesimal ideal $\mm$ is both
open and closed in the order topology: each infinitesimal has a
sufficiently small infinitesimal neighborhood inside $\mm$, while for
$x\notin\mm$ a neighborhood of radius $|x|/2$ stays outside $\mm$.
Consequently the function
\[
 f(x)=\begin{cases}0,&x\in\mm,\\1,&x\notin\mm\end{cases}
\]
is locally constant and has derivative zero everywhere in the intrinsic
ordered-field sense (derivatives with radii from $F$), but $f(0)\ne f(1)$. Local differentiability alone
therefore does not give the ordinary global mean value theorem or
uniqueness of an antiderivative.

There is no conflict with the use of restricted-analytic transfer.
The externally specified set of all real infinitesimals is not a
new definable predicate in the restricted-analytic language used
here. Our formulas come from compatible real analytic germs or explicit
algebraic definitions, not from arbitrary functions that happen to be
locally constant on intrinsic neighborhoods. The same function is locally
constant for the full fine topology of $\No$ as well, since $\mm$ is also open
there; \cref{e3:so:sec:topology} compares the two topologies.

\subsection{Ordinary finite partitions do not automatically define an integral}

Even the function $f(x)=x$ on $[0,1]\subset F$ exposes a difficulty.
For any partition with an ordinary finite number $N$ of subintervals,
put $\Delta_i=x_i-x_{i-1}>0$. Its upper and lower Darboux sums satisfy
\[
 U-L=\sum_{i=1}^N\Delta_i^2\ge\frac1N,
\]
by Cauchy--Schwarz and $\sum_i\Delta_i=1$. Given any positive
infinitesimal $\eps$, this difference is always greater than $\eps$,
whatever ordinary finite $N$ is chosen. Hence the usual Darboux
criterion requiring gaps smaller than \emph{every positive element
of $F$} fails for these partitions. A mesh criterion without a
requirement that the relevant partitions exist would instead become
vacuous below infinitesimal mesh sizes.

Nevertheless the polynomial antiderivative $x^2/2$ is perfectly
well-defined, as are our polygon-area formula and the selected analytic
cap primitive. A theory of generalized integration must say which
functional, which functions, and which summation or partition notion it
uses. This report does not infer a countably additive area measure
from finite-polygon additivity. It also does not claim completeness,
ordinary topological compactness, arbitrary smooth-ODE existence, or
a supremum-based length for every curve in surreal space.

\subsection{Finite dimension does not remove size or support conditions}

The ambient coordinate count is three, but the class $\No^3$ remains
a proper class. The set-field construction in \cref{e3:sec:scalars}
allows every finite configuration in the report to be treated in a
single honest field and every needed analytic evaluation to be interpreted
there. Passing to a larger workspace preserves the answers. Neither
this compatibility nor the finite dimension permits unrestricted
summation of arbitrary families of surreals.

Likewise, ordinary-size angles do not imply ordinary-size arc lengths.
For an infinite $R$ and a positive real $a$, the minor arc has infinite
length $Ra$. Conversely, a bounded physical arc on an infinite sphere
has infinitesimal angular length and is governed by the local expansions
above. Choosing a global oscillatory phase at infinite scalar arguments
would add another analytic convention; it would not change these
coordinate-defined directions or finite geometric angles.

\part{The rotation group \texorpdfstring{$\SO(3,\No)$}{SO(3,No)}}\label{e3:part:so}

Part~II is source~09, apart from the material printed once in
\cref{e3:sec:scalars,e3:sec:vectors,e3:sec:rotations}. Its axis theorem, Rodrigues
formula, quaternion action, spin map and Cayley chart are \cref{e3:thm:rodrigues},
\cref{e3:eq:quatrotation}, \cref{e3:thm:spin} and \cref{e3:so:thm:cayley}.
Vectors are written in plain type, and $G(F)=\SO(3,F)$ for a real closed field $F$.

Source~09 organizes its material specifically around the three-dimensional
rotation group. Centralizers at half-turns, a globally valid matrix Cayley chart
with its exceptional locus, infinitesimal conjugacy of finite subgroups, and the
differential of the \emph{rotation} exponential are developed together. The
ordinary quaternion algebra and spin background is standard \cite{Voight}. The
surreal field and restricted-analytic background is due to Conway, Gonshor, van
den Dries--Ehrlich and Ehrlich--Kaplan; the analytic constructions here rely
specifically on \cite{vdDE,EK}. The perfectness of the infinitesimal $\SO(3)$
group has a published precedent: Bays--Peterzil explicitly record that every
element is a commutator \cite[\S3.1]{BP}. \Cref{e3:so:thm:perfect} gives an
elementary, exact quaternion construction rather than presenting that fact as
new.

\section{Frames, Euler coordinates, and quaternion matrices}\label{e3:so:sec:frames}

This section collects source~09's further finite algebra of $G(F)$.

\subsection{Products of reflections and transitivity}

Reflection in a plane perpendicular to a unit vector \(u\) is
\(S_u=I-2uu^T\), of determinant \(-1\). In an oriented perpendicular plane,
two such reflections produce the usual doubled planar angle; the assertion is
a finite matrix multiplication. Every rotation is a product of two plane
reflections: take the two planes through its axis, choosing their normal lines
by a half-angle square root on the algebraic circle. The half-turn and identity
are included. Alternatively this follows from the explicit quaternion lifts
of \cref{e3:sub:quaternions}.

Given two unit vectors, complete each to an oriented orthonormal frame. The
matrix sending the first frame to the second belongs to \(G(F)\). Thus the
orbits of \(G(F)\) on \(F^3\) are the origin and the spheres of each positive
\(F\)-radius. The stabilizer of a unit vector \(n\) is
\begin{equation}\label{e3:so:eq:torus}
 T_n=\{\Rot(n;c,s):c^2+s^2=1\}\cong\SO(2,F),
 \qquad G(F)/T_n\cong S^2(F).
\end{equation}
This homogeneous-space statement is algebraic and semialgebraic; it does not
assert that these are ordinary real manifolds.

\subsection{Euler coordinates and their degeneracy}

Write \(\Rot_z(c,s)=\Rot(e_3;c,s)\) and \(\Rot_y(c,s)=\Rot(e_2;c,s)\).
For every \(A\in G(F)\) there is a decomposition
\begin{equation}\label{e3:so:eq:euler}
 A=\Rot_z(c_\alpha,s_\alpha)\,
   \Rot_y(c_\beta,s_\beta)\,\Rot_z(c_\gamma,s_\gamma),
 \qquad s_\beta\ge0.
\end{equation}
To construct it, put \(n=Ae_3\), \(c_\beta=n_3\), and
\(s_\beta=\sqrt{n_1^2+n_2^2}\). When \(s_\beta\ne0\), choose
\(c_\alpha=n_1/s_\beta\), \(s_\alpha=n_2/s_\beta\). Then the first two factors
send \(e_3\) to \(n\), and their inverse times \(A\) fixes \(e_3\), giving the
last factor. If \(s_\beta=0\), choose any suitable first factor and solve in the
remaining stabilizer. The nonuniqueness there is the algebraic form of gimbal
lock; infinitesimal \(s_\beta\) also makes these coordinates poorly conditioned.
For actual surreals all three circle pairs admit finite angles, as shown in
\cref{e3:thm:finiteaxis}.

\subsection{Quaternion coordinates}

For \(q=a+bi+cj+dk\) of norm one the rotation matrix is
\begin{equation}\label{e3:so:eq:qmatrix}
\begin{pmatrix}
1-2(c^2+d^2)&2(bc-ad)&2(bd+ac)\\
2(bc+ad)&1-2(b^2+d^2)&2(cd-ab)\\
2(bd-ac)&2(cd+ab)&1-2(b^2+c^2)
\end{pmatrix}.
\end{equation}
Composition of rotations is quaternion multiplication in the same order as
matrix multiplication, with column-vector convention.

\subsection{Complex matrices, projective coordinates, and a caution}

The faithful algebra representation
\begin{equation}\label{e3:so:eq:su2}
 a+bi+cj+dk\ \longmapsto\
 \begin{pmatrix}a+bi&c+di\\-c+di&a-bi\end{pmatrix}
 \quad\text{in }M_2(F(i))
\end{equation}
has determinant \(N(q)\) and sends quaternion conjugation to conjugate
transpose. Norm-one quaternions are exactly the corresponding \(\SU(2,F(i))\)
matrices, where conjugation fixes \(F\) and sends \(i\) to \(-i\).

Every nonzero quaternion can be normalized by its positive norm. Therefore
\begin{equation}\label{e3:so:eq:projective}
 G(F)\cong\HH_F^\times/F^\times
       \cong\mathbb P^3(F)
\end{equation}
as a group with the quaternion-induced law, and as a semialgebraic space.
The second statement concerns \(F\)-points: the norm has no zero on
\(F^4\setminus\{0\}\). It is \emph{not} an isomorphism of algebraic varieties
\(G\cong\mathbb P^3\). The algebraic projective model is the open subset
\(N\ne0\); the norm quadric has points after algebraic closure. This distinction
prevents an incorrect inference that the affine algebraic group is projective.

The quaternion quotient also explains why three dimensions are special. It is
not necessary to choose an axis continuously to compose rotations. A unit
quaternion carries a uniformly bounded four-coordinate representation, with
only the unavoidable sign identification.

\subsection{A quaternion extraction algorithm without a small divisor}

Cayley coordinates become large near a half-turn. A uniformly bounded alternative
works directly from \(A\). Set
\begin{align*}
 d_0&=1+\tr A,&
 d_1&=1+2A_{11}-\tr A,\\
 d_2&=1+2A_{22}-\tr A,&
 d_3&=1+2A_{33}-\tr A.
\end{align*}
By \cref{e3:so:eq:qmatrix}, these are \(4a^2,4b^2,4c^2,4d^2\), respectively, for
either unit lift. They are nonnegative and sum to four. Choose a largest one,
so \(d_j\ge1\), and choose the sign of the lift so that its \(j\)-th coordinate
is \(\sqrt{d_j}/2\ge1/2\). The other coordinates are recovered from
\begin{align*}
 4ab&=A_{32}-A_{23},&4ac&=A_{13}-A_{31},&4ad&=A_{21}-A_{12},\\
 4bc&=A_{12}+A_{21},&4bd&=A_{13}+A_{31},&4cd&=A_{23}+A_{32}.
\end{align*}
All divisions are by a quantity of magnitude at least two. This is an exact
ordered-field algorithm with four branches and positive square roots. Ties can
be resolved by the first index. It does not depend on a floating-point tolerance,
and it does not become singular when the rotation is infinitesimally close to a
half-turn.

An axis itself cannot be extracted with comparable conditioning at the identity:
there is no preferred axis there. The trace loses first-order information near
identity, since \(1-c\) is quadratic in a small angle; the skew coordinates
retain the linear part.

\section{Conjugacy, centralizers, and maximal axis groups}\label{e3:so:sec:conjugacy}

\begin{theorem}[Conjugacy and the half-turn exception]\label{e3:so:thm:centralizers}
For \(A,B\in G(F)\):
\begin{enumerate}[label=\textup{(\roman*)}]
\item They are conjugate in \(G(F)\) if and only if \(\tr A=\tr B\).
\item If \(A\ne I\) is not a half-turn and has axis \(Fn\), then
      \(C_{G(F)}(A)=T_n\).
\item For a half-turn \(H_n\),
\[
 C_{G(F)}(H_n)=\{B\in G(F):Bn=\pm n\}
             \cong\Orth(2,F)\cong T_n\rtimes C_2.
\]
\item The normalizer of \(T_n\) is the same line stabilizer. Its quotient by
      \(T_n\) is \(C_2\), acting by inversion on \(T_n\).
\end{enumerate}
\end{theorem}
\begin{proof}
Conjugation preserves trace and sends \(\Rot(n;c,s)\) to \(\Rot(Bn;c,s)\).
Equal traces give equal \(c\). The possible \(s\)'s differ only in sign;
changing \(n\) to \(-n\) absorbs that sign. Transitivity on unit vectors proves
(i), including the half-turn case.

An element commuting with a nonidentity rotation preserves its one-dimensional
fixed space, so sends \(n\) to \(n\) or \(-n\). In the first case it is in
\(T_n\), which is abelian. In the second case it conjugates the planar rotation
to its inverse. A nonidentity circle element equals its inverse exactly at the
half-turn. This proves (ii) and (iii). In a frame ending in \(n\), the line
stabilizer has matrices
\(\operatorname{diag}(B,\det B)\) with \(B\in\Orth(2,F)\).
A half-turn about a perpendicular axis realizes its nontrivial component.

Finally, the common fixed space of all of \(T_n\) is \(Fn\), so a normalizer
must preserve that line. Every line-preserving rotation either fixes or reverses
its orientation, and accordingly centralizes or inverts \(T_n\). This proves (iv).
\end{proof}

In particular, \(G(F)\) has trivial center. The centralizers of non-half-turns
also show that every nontrivial connected algebraic axis torus is conjugate to
\(T_{e_3}\). ``Connected'' in that sentence refers to the algebraic torus, not
to its fine topological point set.

\begin{corollary}[Commuting rotations]\label{e3:so:cor:commuting}
Two nonidentity rotations commute only in the following cases: they have the
same axis, or they are half-turns about perpendicular axes. Both possibilities
do occur.
\end{corollary}
\begin{proof}
If either is not a half-turn, use \cref{e3:so:thm:centralizers}(ii). For two
half-turns, \(H_n\) preserves the line \(Fm\) precisely when \(m\) is parallel
to \(n\) or perpendicular to it: its two eigenspaces are \(Fn\) and \(n^\perp\).
\end{proof}

Thus the three coordinate half-turns and the identity form a Klein four group.
The slogan that commuting spatial rotations always share an axis is false,
over both the reals and the surreals.

\section{The standard-part extension}\label{e3:so:sec:standardpart}

From now on suppose \(F\supseteq\R\). Every matrix entry of a rotation belongs
to \(\OO_F\), so standard part can be applied entrywise to every rotation,
not only to a restricted subgroup.

\begin{theorem}[Real rotation plus an infinitesimal rotation]\label{e3:so:thm:split}
There is a split exact sequence
\begin{equation}\label{e3:so:eq:split}
 1\longrightarrow K_F\longrightarrow G(F)
 \xrightarrow{\ \st\ }\SO(3,\R)\longrightarrow1,
 \qquad K_F=G(F)\cap(I+M_3(\mm_F)).
\end{equation}
Every \(A\in G(F)\) has the unique decomposition
\begin{equation}\label{e3:so:eq:UR}
 A=UA_0,\qquad A_0=\st(A),\quad U=AA_0^{-1}\in K_F.
\end{equation}
Consequently \(G(F)\cong K_F\rtimes\SO(3,\R)\), with the usual conjugation action.
\end{theorem}
\begin{proof}
The equations \(A^TA=I\), \(\det A=1\), and matrix multiplication use finitely
many sums and products. Applying the ring homomorphism \(\st\) preserves them.
The inclusion of constant real matrices is a section. The kernel is exactly
as displayed, and multiplying by \(\st(A)^{-1}\) gives the unique decomposition.
The transported product on pairs, for \(U,V\in K_F\) and
\(B,B'\in\SO(3,\R)\), is
\[
 (U,B)(V,B')=\bigl(U(BVB^{-1}),BB'\bigr).
\]
\end{proof}

\begin{proposition}[Exact coordinates on the kernel]\label{e3:so:prop:kernelchart}
The Cayley map restricts to a bijection \(\mm_F^3\to K_F\), with multiplication
\(t\star u=(t+u+t\times u)/(1-t\cdot u)\). The unit quaternions infinitesimally
close to \(1\) map isomorphically onto \(K_F\).
\end{proposition}
\begin{proof}
For \(t\in\mm_F^3\), \cref{e3:so:eq:cayley} gives \(\cay(t)-I\in M_3(\mm_F)\).
Conversely, for \(A\in K_F\), \(1+\tr A\) has standard part four, so
\cref{e3:so:eq:cayleyinvcoords} is defined and infinitesimal. The denominator in the
product law has standard part one. The distinguished spin lift is
\((1,t)/\sqrt{1+\norm{t}^2}\), which has standard part \(1\). Among the two
lifts only this one has that standard part, proving the isomorphism.
\end{proof}

The quaternion side of \cref{e3:so:thm:split,e3:so:prop:kernelchart} is the
split sequence \(1\to U_1\to S^3(\No)\to S^3(\R)\to1\) of the surquaternion
report (\texttt{squat:sub:unitreduction}), whose kernel \(U_1\) is the group of
unit quaternions infinitesimally close to \(1\)
(\texttt{squat:thm:rotation-log}). The spin map carries it isomorphically onto
\(K\).

When \(F\ne\R\), the kernel is nontrivial: any nonzero
\(\eps\in\mm_F\) gives \(\cay(\eps e_1)\ne I\). It is proper because a
nonidentity constant real rotation has nontrivial standard part. Thus
\(\SO(3,\No)\) is \emph{not abstractly simple}. Its simple algebraic group and
simple Lie algebra are compatible with this fact.

The split extension is not a direct product. For a real rotation \(B\),
\begin{equation}\label{e3:so:eq:conjcay}
 B\cay(t)B^{-1}=\cay(Bt),
\end{equation}
and this action is visibly nontrivial. Two rotations with the same standard
part can therefore interact at finer scales rather than differing by a central
error.

The quotient has a useful metric interpretation. For \(x\in\OO_F^3\),
\(\st(Ax)=\st(A)\st(x)\). Infinitesimal rotations are invisible on the real
shadow of every finite vector. This conclusion need not hold for infinite
vectors; \cref{e3:so:sec:examples} quantifies the failure.

\section{A perfect, uniquely divisible infinitesimal group}\label{e3:so:sec:perfect}

The word infinitesimal does not mean abelian. Nor does a Baker--Campbell--Hausdorff
formula imply that the entire kernel is nilpotent. The following results make
both points precise. Throughout this section, \(F\supsetneq\R\) is real closed;
the same finite proofs apply to the full surreal class.

\begin{lemma}[No infinitesimal torsion]\label{e3:so:lem:torsion}
If \(U\in K_F\) and \(U^n=I\) for an ordinary positive integer \(n\), then
\(U=I\).
\end{lemma}
\begin{proof}
Write \(U=\Rot(n_0;c,s)\) and regard \(z=c+is\) as an element of \(F(i)\).
The trace gives \(c\in1+\mm_F\). Since \(s^2=1-c^2\), also \(s\in\mm_F\).
The equation \(U^n=I\) gives \(z^n=1\). The polynomial \(X^n-1\) already
splits into distinct linear factors over the ordinary complex subfield.
Hence \(z\) is an ordinary complex root of unity. Its standard part is one,
so \(z=1\), and \(U=I\).
\end{proof}

\begin{theorem}[Unique divisibility]\label{e3:so:thm:divisible}
For every ordinary \(n\ge1\), the power map \(U\mapsto U^n\) is a bijection
of \(K_F\). No nontrivial finite group is a homomorphic image of \(K_F\).
\end{theorem}
\begin{proof}
For \(U=\Rot(n_0;c,s)\), choose an \(n\)-th root \(w\) of \(z=c+is\) in the
algebraically closed field \(F(i)\). Its squared norm satisfies
\(N(w)^n=N(z)=1\), and \(N(w)>0\), so \(N(w)=1\). Its coordinates are bounded
and \(\st(w)\) is an ordinary complex \(n\)-th root of unity. Dividing \(w\)
by this root gives a circle element \(w_0\) infinitesimally close to one with
\(w_0^n=z\). Rotation around \(n_0\) by this algebraic angle gives a root in
\(K_F\).

For uniqueness, the identity case follows from \cref{e3:so:lem:torsion}. Otherwise,
if \(V^n=U\), then \(V\) commutes with \(U\). A nonidentity element of \(K_F\)
is not a half-turn, so \(V\in T_{n_0}\) by \cref{e3:so:thm:centralizers}. Two roots
in that abelian group have quotient of finite order in \(K_F\), and are equal.
Finally, an image of a divisible group is divisible. In a finite image of order
\(m\), every \(m\)-th power is the identity, so surjectivity of that power map
forces the image to be trivial.
\end{proof}

Unique divisibility is a statement about solving one power equation at a time.
It does not say that \(K_F\) is a vector space over \(\Q\), because its group
operation is not commutative.

\begin{theorem}[Every infinitesimal rotation is one commutator]\label{e3:so:thm:perfect}
For every \(U\in K_F\) there are \(V,W\in K_F\) with
\[
 U=\comm{V}{W}:=VWV^{-1}W^{-1}.
\]
Thus \(K_F=[K_F,K_F]\). In particular, it is nonabelian and neither solvable
nor nilpotent, and it has no nontrivial solvable quotient.
\end{theorem}
\begin{proof}
The identity is immediate. Take the unique lift of \(U\ne I\) near one,
\(q=(c,w)\in\Sp(1,F)\). Then \(0<c<1\), \(c-1\in\mm_F\), and \(w\ne0\).
Put
\[
 s=\left(\frac{1-c}{2}\right)^{1/4},\qquad a=\sqrt{1-s^2},
 \qquad p=(a,se_1),\quad r=(a,se_2).
\]
These are unit quaternions near one. Direct quaternion multiplication gives
\begin{equation}\label{e3:so:eq:commutatorscalar}
 \re\comm{p}{r}=1-2s^4=c.
\end{equation}
For completeness, a version before the equal-parameter specialization is
\[
 \re\bigl((a,se_1)(b,te_2)(a,-se_1)(b,-te_2)\bigr)
 =(a^2+s^2)(b^2+t^2)-2s^2t^2.
\]
Since \(\comm{p}{r}\) and \(q\) have norm one and the same scalar part, their
imaginary vectors have the same nonzero norm. A unit quaternion \(h\) conjugates
one vector to the other, by the transitivity and spin surjectivity already
proved. Therefore
\[
 q=\comm{hph^{-1}}{hrh^{-1}}.
\]
Conjugation by a unit quaternion preserves infinitesimal imaginary coordinates:
its matrix entries are bounded. Thus both factors are still near one. Applying
\(p\) proves the claim in \(K_F\). A nontrivial perfect group is neither
solvable nor nilpotent; every homomorphic image is perfect, and a solvable perfect
group is trivial.
\end{proof}

The conclusion has a literature precedent \cite[proof of Claim 3.3]{BP}; the
proof above exposes the exact fourth-root scale used in a commutator
construction. The two factors generally live at a \emph{larger infinitesimal
magnitude} than their commutator. There is no contradiction with the abelian
leading layers established in \cref{e3:so:sec:filtration}.

The whole group \(G(F)\) likewise has commutator width one. Choose an algebraic
square root \(B\) of a rotation \(A\) around the same axis, and let \(H\) be a
half-turn about a perpendicular axis. Then \(HBH^{-1}=B^{-1}\), so
\(\comm{B}{H}=B^2=A\). This latter proof alone does not prove the infinitesimal
claim, since the perpendicular half-turn is not infinitesimal.

\section{Finite subgroups have no new surreal forms}\label{e3:so:sec:finitegroups}

\begin{lemma}[Positive matrix square roots over a real closed field]
\label{e3:so:lem:positiveroot}
Every symmetric positive definite matrix over \(F\) has a unique symmetric
positive definite square root. If \(P-I\in M_d(\mm_F)\), its positive square
root and inverse square root are also infinitesimally close to \(I\).
\end{lemma}
\begin{proof}
The symmetric spectral theorem holds over any real closed field, by the
argument of \cref{e3:sub:quadrics}: an eigenvector \(z\ne0\) over \(F(i)\) has
eigenvalue \((z^*Pz)/(z^*z)\in F\), and induction on perpendicular complements
gives an orthonormal eigenbasis. (Source~09 gave the same argument in full.)
For positive definite \(P\), its eigenvalues
\(\lambda_j\) are positive, so take \(\sqrt{\lambda_j}\) in that basis.
Any symmetric positive square root commutes with \(P\), preserves its
eigenspaces, and has only the eigenvalue \(\sqrt{\lambda_j}\) on each such
space; this proves uniqueness.

If every entry of \(Y=P-I\) is infinitesimal, a unit eigenvector gives
\(\abs{\lambda_j-1}=\abs{x^TYx}\le d\max\abs{Y_{ij}}\), by
\(\sum_i\abs{x_i}\le\sqrt d\). Thus \(\lambda_j-1\),
\(\sqrt{\lambda_j}-1\), and \(\lambda_j^{-1/2}-1\) are infinitesimal.
Conjugation by the bounded orthonormal eigenbasis preserves this property.
\end{proof}

\begin{theorem}[Infinitesimal conjugacy of finite subgroups]\label{e3:so:thm:finiterigidity}
Let \(H\le G(F)\) be finite. Standard part is injective on \(H\), and there is
\(C\in K_F\) such that
\begin{equation}\label{e3:so:eq:finiteconjugacy}
 C^{-1}hC=\st(h)\quad\text{for every }h\in H.
\end{equation}
In particular, every finite surreal rotation group is conjugate to an ordinary
real rotation group by a rotation with standard part \(I\).
\end{theorem}
\begin{proof}
The kernel of \(\st|_H\) is finite and lies in the torsion-free group \(K_F\),
so it is trivial. Write \(\rho(h)=\st(h)\) and average the two representations:
\begin{equation}\label{e3:so:eq:average}
 B=\frac1{\abs H}\sum_{h\in H}h\rho(h)^{-1}.
\end{equation}
Each summand is infinitesimally close to \(I\); hence so is \(B\), and its
determinant has standard part one. In particular it is invertible and has
positive determinant. Reindexing the finite sum shows
\(gB=B\rho(g)\) for every \(g\in H\).

Set \(P=B^TB\). This is symmetric positive definite and infinitesimally close
to \(I\). Orthogonality of both representations gives
\(\rho(g)^TP\rho(g)=P\), so \(P\) commutes with \(\rho(g)\).
By uniqueness of the positive square root, \(P^{-1/2}\) commutes with all those
matrices as well. Put
\[
 C=BP^{-1/2}.
\]
Then \(C^TC=I\), and both determinants in the product are positive, so
\(\det C=1\). The preceding lemma gives \(C\in K_F\).
Finally \(gC=C\rho(g)\), which is exactly \cref{e3:so:eq:finiteconjugacy}.
\end{proof}

This proof is finite averaging followed by algebraic polar decomposition. It
uses neither an integral over the full group nor a completeness theorem.
The same argument works in \(\Orth(d,F)\) for any fixed finite dimension, with
the appropriate determinant statement.

\subsection{Classification and what it does not claim}

The familiar finite-subgroup classification of \(\SO(3,\R)\) therefore applies:
cyclic groups \(C_n\), rotational dihedral groups \(D_n\) of order \(2n\), and
the tetrahedral, octahedral, and icosahedral groups \(A_4,S_4,A_5\), together
with the trivial group. The word dihedral here denotes a group of
\emph{three-dimensional rotations}: the elements outside its principal cyclic
subgroup are perpendicular half-turns, not spatial reflections.

One can see the numerical restriction by counting poles on the ordinary unit
sphere. If \(H\ne1\), each nonidentity element fixes exactly two points. At a
pole the stabilizer is cyclic. If the pole orbits have stabilizer orders
\(m_1,\ldots,m_k\ge2\), counting incident pairs gives
\begin{equation}\label{e3:so:eq:polecount}
 \sum_{j=1}^k(1-1/m_j)=2-2/\abs H.
\end{equation}
It follows that \(k=2\) or \(3\). For \(k=2\), both stabilizers have order
\(\abs H\), giving a cyclic group. For \(k=3\), order the \(m_j\)'s; the only
possibilities are
\[
 (2,2,n),\quad(2,3,3),\quad(2,3,4),\quad(2,3,5),
\]
with group orders \(2n,12,24,60\). The usual spherical realizations identify
them with the dihedral and three polyhedral groups. This last identification
is the classical finite rotation-group classification; the new field introduces
no additional possibility because of \cref{e3:so:thm:finiterigidity}.

Individual finite-order surreal rotations need not have real axes. A conjugate
\(C H C^{-1}\) of a real finite group by a nonreal \(C\in K_F\) generally has
nonreal entries. The theorem says that these are conjugate realizations, not
that every finite subgroup is literally contained in the constant matrices.
It also says nothing about arbitrary infinite or set-generated subgroups.

\section{The Lie algebra and angular velocity}\label{e3:so:sec:lie}

The algebraic tangent space at the identity is computed with the formal dual
number \(\kappa\), satisfying \(\kappa^2=0\). It is not a nonzero surreal
infinitesimal: a field has no such nilpotents. Substituting \(I+\kappa X\) in the
orthogonality equation gives
\[
 (I+\kappa X)^T(I+\kappa X)=I+\kappa(X^T+X).
\]
Thus
\begin{equation}\label{e3:so:eq:liealgebra}
 \so(3,F)=\{X\in M_3(F):X^T=-X\}
         =\{\widehat u:u\in F^3\}.
\end{equation}
Its bracket and adjoint action are exactly
\begin{equation}\label{e3:so:eq:liecross}
 [\widehat u,\widehat v]=\widehat{u\times v},\qquad
 \Ad_A(\widehat u)=\widehat{Au}.
\end{equation}
This identifies the Lie algebra with the cross-product algebra of three-space.

\begin{proposition}[Simplicity and the spin factor]\label{e3:so:prop:liesimple}
The Lie algebra \(\so(3,F)\) is simple. The differential of the spin map is
\(u\mapsto2\widehat u\) on imaginary quaternions.
\end{proposition}
\begin{proof}
Let a nonzero ideal contain \(u\ne0\) in cross-product coordinates. Then it
contains \(u\times v\) for every \(v\). The image of this linear map is exactly
\(u^\perp\), by \cref{e3:so:eq:hatpowers}. Together with \(u\), these vectors span
\(F^3\), so the ideal is the entire algebra. The algebra is nonabelian.
For the spin differential, conjugate \(x\) by \(1+\kappa u\): the linear term is
\(ux-xu=2u\times x\). The imaginary-quaternion Lie bracket is itself
\(2u\times v\), consistent with the factor two.
\end{proof}

For a differentiable \(F\)-parameter family of rotations \(A(t)\), whenever the
chosen derivative has the product rule, differentiation of \(A^TA=I\) gives
\[
 \Omega_{\mathrm{space}}=A'A^{-1}\in\so(3,F),\qquad
 \Omega_{\mathrm{body}}=A^{-1}A'\in\so(3,F),
 \qquad \Omega_{\mathrm{space}}=A\Omega_{\mathrm{body}}A^{-1}.
\]
These are the two angular-velocity conventions. Their vector representatives
are related by \(\omega_{\mathrm{space}}=A\omega_{\mathrm{body}}\).
The algebraic identity remains true for an entrywise field derivation, but the
existence or uniqueness of solutions of the resulting differential equation is
an additional question. In particular, the Berarducci--Mantova derivation of
surreal coefficients is not the variable derivative used in the fine
differential of \cref{e3:so:sec:globalexp}.

A further distinction is essential: \(\so(3,\No)\) is a three-dimensional
\(\No\)-vector space, whereas \(\so(3,\mm)\) is not closed under multiplication
by arbitrary infinite scalars. The latter is the domain of the infinitesimal
logarithm, not a replacement for the full algebraic tangent space.

\section{Strong exponential, logarithm, and BCH}\label{e3:so:sec:hahn}

\subsection{A support theorem, not a convergence assumption}

Let \(\Gamma\) be a set-sized divisible ordered abelian group, and work in
\(F_\Gamma=\R((t^\Gamma))\), the Hahn field with well-ordered supports. It is
real closed \cite{vdDE,RepoQuaternion}. A common surreal workspace uses \(t^\gamma=\omega^{-\gamma}\)
for an ordered subgroup of surreal exponents. The rational span of the supports
of finitely or set-many given normal forms supplies a set-sized divisible
exponent group. This localization is larger than the finite-algebra subfield
of \cref{e3:so:prop:localization}; its purpose is to contain the strong sums.

Strong summability is the notion of \cref{e3:sec:scalars}: a well-ordered union
of supports, each exponent in only finitely many members, and coefficientwise
sums. It is independent of ordinary sequence convergence. The following
all-lengths form of the positive-support lemma is \texttt{squat:lem:neumann}
(Lemma~7.1 of the surquaternion report), with the same Higman argument; the
repository's \texttt{Surreal/HahnSeries/Neumann.lean} proves the
support statements for other reports' labels. Source~09's statement and proof
are kept because they add fixed-size matrix arguments.

\begin{lemma}[All-length support control]\label{e3:so:lem:neumann}
Let \(S\subseteq\Gamma_{>0}\) be well ordered. The additive monoid \(S^*\)
generated by \(S\) is well ordered, and every \(\gamma\in\Gamma\) is the sum of
only finitely many finite ordered words in \(S\). Consequently every formal
power series in finitely many commuting or noncommuting variables can be
evaluated strongly at Hahn elements or fixed-size matrices whose entries have
positive valuation. The nonconstant support lies in \(S^*\setminus\{0\}\).
\end{lemma}
\begin{proof}
Use Higman's finite-word lemma for the well-quasi-ordered alphabet \(S\)
\cite{Higman}. In any infinite sequence of words, an earlier word embeds into
a later word with each letter no larger than the corresponding letter.
Such an embedding makes the total sum no larger; unless the words are
identical, either an extra positive letter or an increased letter makes it
strictly larger. Therefore a strictly decreasing sequence of word sums is
impossible, and infinitely many distinct words cannot have one fixed sum.
This proves the two support assertions.

For the evaluation assertion take the finite union of the positive supports of
all input entries. Each coefficient of a fixed total exponent involves only
finitely many such words. The finite choices of variables and intermediate
matrix indices add only finitely many contributions per word. Thus the
formal evaluation is strongly summable, including all word lengths at once.
\end{proof}

The same argument applies to a set-supported surreal calculation after
localizing the supports. It justifies rearrangements in what follows; it does
not assert that \(n\gamma\) is cofinal in an arbitrary ordered value group.

\begin{theorem}[The infinitesimal exponential and logarithm]\label{e3:so:thm:explog}
On a divisible Hahn workspace, and on the full surreal class by strong
normal-form evaluation, the maps
\begin{align}
 \Exp_{\mm}(X)&=\sum_{n\ge0}\frac{X^n}{n!},
       &&X\in\so(3,\mm),\label{e3:so:eq:infexp}\\
 \Log_{\mm}(U)&=\sum_{n\ge1}\frac{(-1)^{n+1}}n(U-I)^n,
       &&U\in K,\label{e3:so:eq:inflog}
\end{align}
are inverse bijections \(\so(3,\mm)\leftrightarrow K\).
They commute with conjugation by any rotation. Their transported multiplication
is the strongly evaluated Baker--Campbell--Hausdorff law, not vector addition.
\end{theorem}
\begin{proof}
The support lemma makes both expressions meaningful. Formal one-variable
identities show that exponential and logarithm are inverse on the ambient
matrix ideals and groups \(M_3(\mm)\) and \(I+M_3(\mm)\). Substitution is
legitimate because the nonconstant arguments have positive support.

For skew \(X\), transposition gives \(\Exp_{\mm}(X)^T=\Exp_{\mm}(-X)\).
The product is \(I\), and the result has standard part \(I\), so its determinant
is \(+1\), not \(-1\). Conversely, if \(U\in K\), then
\[
 \Log_{\mm}(U)^T=\Log_{\mm}(U^T)
 =\Log_{\mm}(U^{-1})=-\Log_{\mm}(U).
\]
The last identity is a formal identity in the one matrix \(U-I\).
Conjugation commutes with all finite products and with the supported sums.
Finally, substituting \(X,Y\) in the formal noncommutative identity
\(\log(e^Xe^Y)=\operatorname{BCH}(X,Y)\) is justified by the same word-support
argument.
\end{proof}

In vector coordinates, if
\(\Exp_{\mm}(\widehat u)\Exp_{\mm}(\widehat v)
 =\Exp_{\mm}(\widehat w)\), then
\begin{equation}\label{e3:so:eq:bchvector}
\begin{split}
 w={}&u+v+\tfrac12u\times v\\
 &+\tfrac1{12}\bigl(u\times(u\times v)+v\times(v\times u)\bigr)
 +\text{terms of total degree at least four}.
\end{split}
\end{equation}
The omitted terms are the uniquely specified formal BCH series. If every
coordinate of \(u,v\) has valuation at least \(\gamma>0\), their sum has
valuation at least \(4\gamma\). This is a finite-order error certificate, not
an assertion of convergence through all positive surreal tolerances.

\begin{proposition}[Exact bridge to Cayley coordinates]\label{e3:so:prop:logcay}
For \(t\in\mm^3\),
\begin{equation}\label{e3:so:eq:logcay}
 \Log_{\mm}(\cay(t))
 =2\sum_{j\ge0}\frac{(-1)^j}{2j+1}(t\cdot t)^j\widehat t.
\end{equation}
Its leading term is \(2\widehat t\). In particular, exponential coordinates
and Cayley coordinates have the same valuation.
\end{proposition}
\begin{proof}
The commuting factors \(I+\widehat t\) and \(I-\widehat t\) give
\[
 \log\bigl((I+\widehat t)(I-\widehat t)^{-1}\bigr)
 =\log(I+\widehat t)-\log(I-\widehat t)
 =2\sum_{j\ge0}\frac{\widehat t^{\,2j+1}}{2j+1}.
\]
Now use \(\widehat t^{\,2j+1}=(-t\cdot t)^j\widehat t\).
Every later term has strictly higher valuation than the first.
\end{proof}

\section{Valuation layers and noncommutativity at every scale}\label{e3:so:sec:filtration}

Use the valuation of \cref{e3:sec:scalars}:
\(v(x)=-\max\operatorname{supp}_{\omega}(x)\), \(v(0)=+\infty\),
so \(v(\omega^{-\gamma})=\gamma\) and infinitesimals have positive valuation.
For a vector or matrix, \(v\) is the minimum of the valuations of its entries.
On a Hahn workspace this is the ordinary Hahn valuation with value group
\(\Gamma\). For all surreals the value group is a class; statements involving
one layer still have finite formulas, and direct-sum notation is localized to
set-sized workspaces.

Because all rotation entries have nonnegative valuation, and so do all entries
of the inverse, \(v(Ax)=v(x)\). The equality follows by applying the two
inequalities from \(A\) and \(A^{-1}\). Thus a rotation preserves not only the
norm, but also the leading scale of every vector.

For \(\gamma>0\), define
\begin{equation}\label{e3:so:eq:filtration}
 K_{\ge\gamma}=\{U\in K:v(U-I)\ge\gamma\},\qquad
 K_{>\gamma}=\{U\in K:v(U-I)>\gamma\}.
\end{equation}
These are normal subgroups of \(G\). If \(t\) is the Cayley coordinate of
\(U\in K\), then
\begin{equation}\label{e3:so:eq:valcay}
 v(U-I)=v(t)=v(\Log_{\mm}U).
\end{equation}
For the first equality use the rational chart and its inverse, whose
denominators have nonzero real standard parts. The second follows from
\cref{e3:so:eq:logcay}.

\begin{theorem}[Graded rotation algebra]\label{e3:so:thm:graded}
For \(\alpha,\beta,\gamma>0\),
\begin{equation}\label{e3:so:eq:commbound}
 \comm{K_{\ge\alpha}}{K_{\ge\beta}}\subseteq K_{\ge\alpha+\beta}.
\end{equation}
For surreals, the map
\begin{equation}\label{e3:so:eq:layer}
 \ell_\gamma(U)=2\st\bigl(\omega^\gamma t(U)\bigr)
\end{equation}
induces an isomorphism
\(K_{\ge\gamma}/K_{>\gamma}\cong(\R^3,+)\). On a Hahn workspace replace
\(\omega^{-\gamma}\) by \(t^\gamma\). The commutator between layers is the
ordinary cross product:
\begin{equation}\label{e3:so:eq:gradedbracket}
 \ell_{\alpha+\beta}(\comm{U}{V})
 =\ell_\alpha(U)\times\ell_\beta(V).
\end{equation}
Conjugation on every layer factors through standard part and is the standard
three-dimensional action of \(\SO(3,\R)\).
\end{theorem}
\begin{proof}
The Cayley law is addition modulo terms of degree at least two. For inputs of
valuation at least \(\gamma\), scaling by \(\omega^\gamma\) kills those terms
under standard part. Thus \(\ell_\gamma\) is a homomorphism with kernel
\(K_{>\gamma}\). It is onto, since
\(U=\cay(\tfrac12\omega^{-\gamma}a)\) maps to \(a\in\R^3\).

For the commutator it is convenient to use \(U=e^X,V=e^Y\) with positive
valuation. Formal multiplication gives
\[
 \Log_{\mm}(\comm{e^X}{e^Y})=[X,Y]
  +\text{Lie terms containing at least three letters}.
\]
Every nonzero term contains both \(X\) and \(Y\). Terms after the bracket have
valuation strictly larger than \(\alpha+\beta\), whereas the bracket has at
least that valuation. Since \(X\) has leading vector twice the Cayley vector,
\cref{e3:so:eq:hatbracket} proves \cref{e3:so:eq:commbound,e3:so:eq:gradedbracket} with the stated
factor two. Lastly \cref{e3:so:eq:conjcay} gives
\(\ell_\gamma(AUA^{-1})=\st(A)\ell_\gamma(U)\).
\end{proof}

On a set-sized workspace the associated graded Lie algebra is therefore
\[
 \bigoplus_{\gamma\in\Gamma_{>0}}\so(3,\R)\,T^\gamma,
 \qquad [aT^\alpha,bT^\beta]=(a\times b)T^{\alpha+\beta},
\]
with vector notation for the coefficients. A leading bracket can vanish when
the leading vectors are parallel; the theorem deliberately gives an inclusion
of filtration levels, not equality of valuations in every case.

The quaternion form of this grading is \texttt{squat:thm:graded} (Theorem~8.6 of
the surquaternion report), where the graded bracket of leading imaginary
coefficients is \(2X\times Y\). The two statements agree: the leading
imaginary coefficient \(X\) of the unit-quaternion lift \((1,t)/\sqrt{1+\norm t^2}\)
of \(U\) is the leading coefficient of its Cayley vector \(t\), so
\(\ell_\gamma(U)=2X\). Hence
\(\ell_{\alpha+\beta}(\comm UV)=2(2X\times Y)=\ell_\alpha(U)\times\ell_\beta(V)\);
the factor two is that of the spin differential (\cref{e3:so:prop:liesimple}).

\subsection{Why finite jets are nilpotent but the kernel is perfect}

For any ordinary integer \(N\ge1\),
\[
 K_{\ge\gamma}/K_{>N\gamma}
\]
is nilpotent of class at most \(N\): an \((N+1)\)-fold iterated commutator has
valuation at least \((N+1)\gamma>N\gamma\). This is a fixed-scale,
finite-depth statement. The full group is the union over \emph{all} positive
scales. Its commutator construction can use parameters at scales such as
\(\gamma/2\), outside a chosen \(K_{\ge\gamma}\). Theorem \ref{e3:so:thm:perfect}
therefore remains fully compatible with the nilpotence of these quotients.

It would be incorrect to conclude merely from the BCH description that the
whole \(K\) is a pronilpotent group in an unspecified topology. In particular,
one fixed sequence of multiples \(n\gamma\) need not exhaust all valuation
scales, and the whole group is perfect rather than residually nilpotent in the
sense of its ordinary lower central series.

\subsection{An imported structural perspective}

For a set-sized proper real closed extension \(F\supsetneq\R\), a theorem of
Bays--Peterzil says that the pure group \(K_F\), with parameters allowed, is
bi-interpretable with the valued field \((F;+,{\cdot},\OO_F)\)
\cite[Theorem 1.1 and Remark 1.2]{BP}. Thus the exact infinitesimal group carries
much more information than an abelian tangent space. This is an imported
model-theoretic theorem, not proved here, and its statement is about set-sized
structures. We do not automatically replace its universe by the full proper
class of surreals.

\section{Finite angles and canonical infinite phases}\label{e3:so:sec:angles}

\subsection{Finite-angle trigonometry from strong sums}

For \(\theta=r+\eps\in\OO\), with \(r\in\R\), \(\eps\in\mm\), define
\begin{equation}\label{e3:so:eq:finitephase}
 \cis(\theta)=e^{ir}\sum_{k\ge0}\frac{(i\eps)^k}{k!}
 =\cos\theta+i\sin\theta.
\end{equation}
Only the infinitesimal part is strongly expanded; \(e^{ir}\) is an ordinary
complex constant. This is the finite-angle construction used in the repository
\cite{RepoTrig} and in the restricted-analytic framework \cite{EK}.
The formal exponential identities prove addition, parity, Pythagoras, and the
usual finite trigonometric identities. In particular \(\cis\) is a
homomorphism from \((\OO,+)\) to the algebraic surcomplex unit circle. It
agrees with the real-centered evaluation of \cref{e3:eq:taylor}.

Source~09's theorem on finite angles has two halves, both printed in Part~I.
The map \(\cis\) is onto the unit circle with kernel exactly \(2\pi\Z\): this is
\cref{e3:thm:phase}, whose second proof is source~09's. Every \(A\in G(\No)\)
has a representation \(A=\Rot(n;\cos\theta,\sin\theta)\) with \(0\le\theta\le\pi\),
unique for \(A\ne I\), with the axis oriented uniquely for \(0<\theta<\pi\): this
is \cref{e3:thm:finiteaxis}, whose proof is source~09's.

Consequently the stabilizer of any chosen oriented unit axis has the group
presentation
\begin{equation}\label{e3:so:eq:finitequotient}
 T_n\cong\OO/(2\pi\Z).
\end{equation}
Its finite angles are surreal, not necessarily ordinary real. A rotation can
have an infinitesimal angle, or differ infinitesimally from an ordinary angle.
There is no need to attach an infinite angle to represent any rotation.

\subsection{The Ehrlich--Kaplan global extension does exist}

It would nevertheless be wrong to say that no canonical global surreal
trigonometry is available. The initial-integer-part construction of
Ehrlich--Kaplan \cite[\S11]{EK} supplies one. In terms of normal forms, write
\[
 \No^\uparrow=\{x:x\text{ has only positive normal-form exponents}\},\qquad
 \Oz=\No^\uparrow+\Z.
\]
The class \(\Oz\) is the initial integer part, often called the omnific
integers. It is not the class of ordinary ordinal numbers. The zero element
is included in \(\No^\uparrow\).

The corresponding phase has a particularly transparent expression. For the
three-part decomposition \(\theta=\theta^\uparrow+r+\eps\), set
\begin{equation}\label{e3:so:eq:globalphase}
 \cis_{\mathrm{EK}}(\theta)=\cis(r+\eps)
 =\Cos\theta+i\Sin\theta.
\end{equation}
This is exactly reduction by \(2\pi\Oz\): a real multiple of a purely infinite
normal form is still purely infinite, so
\begin{equation}\label{e3:so:eq:globalperiod}
 2\pi\Oz=\No^\uparrow+2\pi\Z,
 \qquad \ker\cis_{\mathrm{EK}}=2\pi\Oz.
\end{equation}
The decomposition into purely infinite and finite parts is additive. Thus the
global phase is a surjective additive-to-multiplicative homomorphism. This
construction is canonical for the specified initial-integer-part normalization;
it is not an assertion that every conceivable global phase convention agrees
with it.

For example,
\[
 \Cos\omega=1,\qquad\Sin\omega=0,\qquad
 \cis_{\mathrm{EK}}(\omega+\pi/2)=i.
\]
The infinite part is removed before evaluating the finite phase. No divergent
Taylor series at \(\omega\) is being summed. On a fixed axis one therefore has
\begin{equation}\label{e3:so:eq:globalaxisquotient}
 (\No,+)/(2\pi\Oz)\cong T_n,
 \qquad \theta\longmapsto \Rot(n;\Cos\theta,\Sin\theta).
\end{equation}
The finite and global presentations describe the \emph{same} axis subgroup,
with different parameter domains and period classes.

The associated spin parameter is
\[
 q_n(\theta)=\Cos(\theta/2)+n\Sin(\theta/2),
 \qquad \ker q_n=4\pi\Oz.
\]
An ordinary \(2\pi\) rotation has spin value \(-1\), while a purely infinite
parameter such as \(\omega\) has spin value \(+1\). Indeed, multiplication by
any nonzero real preserves \(\No^\uparrow\), so the purely infinite parts of
\(2\pi\Oz\) and \(4\pi\Oz\) coincide; their ordinary integer parts differ.
This prevents misleading interpretations of an infinite parameter as an
ordinary number of physical revolutions.

\section{The global rotation exponential and its differential}\label{e3:so:sec:globalexp}

\subsection{A radial map on the whole Lie algebra}

For \(w\in\No^3\), put \(r=\norm w\). Define
\begin{equation}\label{e3:so:eq:globalexp}
 \mathcal E(\widehat w)=
 \begin{cases}
 I+\dfrac{\Sin r}{r}\widehat w
   +\dfrac{1-\Cos r}{r^2}\widehat w^{\,2},&r\ne0,\\[6pt]
 I,&r=0.
 \end{cases}
\end{equation}
Equivalently, it is the spin image of
\(\Cos(r/2)+(w/r)\Sin(r/2)\). This is a well-defined global map to
\(G(\No)\), and is equivariant under conjugation. It agrees with the strong
matrix exponential for infinitesimal \(w\); finite \(w\) are evaluated using
finite trigonometry, not by assuming strong summability of the uncentered
matrix power series.

The map is already onto on the finite ball \(\norm w\le\pi\), by
\cref{e3:thm:finiteaxis}. On any fixed line of generators it is a homomorphism:
\[
 \mathcal E(\alpha\widehat n)\mathcal E(\beta\widehat n)=\mathcal E((\alpha+\beta)\widehat n).
\]
It is not a homomorphism from the whole additive Lie algebra. The additive
domain is abelian while its image is not; the BCH correction is nonzero even
for small independent generators.

\begin{theorem}[Complete fibers of the rotation exponential]\label{e3:so:thm:expfibers}
Let \(A\ne I\) have the principal representation
\(A=\Rot(n;\cos\theta,\sin\theta)\), \(0<\theta\le\pi\). Then
\begin{equation}\label{e3:so:eq:nonidentityfiber}
 \mathcal E^{-1}(A)=\{\widehat{(\theta+2\pi z)n}:z\in\Oz\}.
\end{equation}
At \(\theta=\pi\) either choice of axis orientation gives the same set.
The identity fiber is
\begin{equation}\label{e3:so:eq:identityfiber}
 \mathcal E^{-1}(I)=\{0\}\ \cup\
 \{\widehat w:\norm w\in(2\pi\Oz)_{>0}\}.
\end{equation}
Restricting to finite generators replaces \(\Oz\) by \(\Z\).
\end{theorem}
\begin{proof}
If \(\mathcal E(\widehat w)\ne I\), its unique fixed line is \(Fw\), so
\(w=\lambda n\). Oddness of sine and evenness of cosine show that
\(\mathcal E(\lambda\widehat n)=\Rot(n;\Cos\lambda,\Sin\lambda)\) for signed \(\lambda\).
Equality of circle values is exactly congruence modulo \(2\pi\Oz\), proving
\cref{e3:so:eq:nonidentityfiber}. When \(\theta=\pi\), negating \(\lambda\) does not
change that congruence class. For the identity, the radial circle value must be
one, which is precisely the condition in \cref{e3:so:eq:identityfiber}.
The finite part of \(\Oz\) is \(\Z\), giving the last assertion.
\end{proof}

The identity fiber is a union of spheres, not a subgroup of the additive
three-dimensional Lie algebra. It must not be called the kernel of a global
group homomorphism. Only a fixed-axis restriction has a kernel.

\subsection{Fine differentiability}

For a function on a surreal finite-dimensional coordinate space, the fine
Fr\'echet derivative at \(x\) is a \(\No\)-linear map \(L\) such that for every
positive surreal \(\eta\), some positive surreal \(\delta\) satisfies
\[
 0<\norm h<\delta\quad\Longrightarrow\quad
 \norm{f(x+h)-f(x)-Lh}<\eta\norm h.
\]
Equivalent finite-dimensional coordinate norms give the same definition.
The quantifiers range over all positive surreal scales. This derivative is not
defined by ordinary real sequences.

The global phase satisfies the usual local derivative identities
\(D\Sin=\Cos\), \(D\Cos=-\Sin\). Indeed, for infinitesimal \(h\) the global
addition formulas reduce the increment to \(\Sin h\) and \(\Cos h\), whose
strong Taylor expansions have bounded quadratic remainders. For any fixed
positive surreal tolerance, one can choose a still smaller radius. The purely
infinite part of the base point causes no obstruction to this local argument.
For \(w\ne0\), differentiating the positive square root gives
\(Dr_w(h)=w\cdot h/r\); the binomial expansion at one justifies its local
quadratic remainder even when \(r\) is infinite or infinitesimal.
Here a quadratic remainder means a bound \(C\norm h^2\) on some neighborhood,
with a fixed nonnegative surreal \(C\) depending on the base point; \(C\) need
not be an ordinary real. Choosing a radius smaller than \(\eta/(C+1)\) then
gives the required fine derivative estimate.

\begin{theorem}[Differential and exact critical radii]\label{e3:so:thm:expdifferential}
Right-translate the differential of \(\mathcal E\) back to the Lie algebra. For
\(w\ne0\), \(r=\norm w\), one obtains
\begin{equation}\label{e3:so:eq:jacobian}
 (d\mathcal E_w(h))\mathcal E(\widehat w)^{-1}=\widehat{\mathcal J(w)h},
 \qquad
 \mathcal J(w)=I+\frac{1-\Cos r}{r^2}\widehat w
          +\frac{r-\Sin r}{r^3}\widehat w^{\,2}.
\end{equation}
Here \(\mathcal E\) is viewed as a function of the vector \(w\), and \(\mathcal J(0)=I\).
For \(r\ne0\),
\begin{equation}\label{e3:so:eq:jacdet}
 \det \mathcal J(w)=\frac{4\Sin^2(r/2)}{r^2}.
\end{equation}
The nonzero critical radii are exactly \(r\in(2\pi\Oz)_{>0}\), and the
fine differential has rank one there. Otherwise it has rank three.
\end{theorem}
\begin{proof}
Write \(w=rn\) and split \(h=h_{\parallel}+h_{\perp}\) along \(n\).
A radial increment changes only the angle and gives
\(\mathcal J(w)h_{\parallel}=h_{\parallel}\). For a perpendicular increment,
\(Dr(h_{\perp})=0\) and \(Dn(h_{\perp})=h_{\perp}/r\).
Differentiate
\(A=cI+(1-c)nn^T+s\widehat n\), with \(c=\Cos r,s=\Sin r\), and multiply by
\(A^T\). The skew matrix that results has vector
\begin{equation}\label{e3:so:eq:jacplane}
 \mathcal J(w)h_{\perp}=\frac{\Sin r}{r}h_{\perp}
                 +\frac{1-\Cos r}{r}\,n\times h_{\perp}.
\end{equation}
This calculation can also be made with \(n=e_3\), then transported by
\cref{e3:so:eq:liecross}. Since \(\widehat w^{\,2}\) is zero on the axis and
\(-r^2I\) on the plane, \cref{e3:so:eq:jacplane} is precisely \cref{e3:so:eq:jacobian}.
The derivative identities preceding the theorem provide quadratic remainders;
finite sums, products, and inverses with nonzero denominators preserve that
property locally. At zero the small-generator expansion is
\(\mathcal E(\widehat h)=I+\widehat h+\smallO(\norm h^2)\), giving \(\mathcal J(0)=I\).

On the perpendicular plane, the determinant of \cref{e3:so:eq:jacplane} is
\[
 \frac{\Sin^2r+(1-\Cos r)^2}{r^2}
 =\frac{2(1-\Cos r)}{r^2}
 =\frac{4\Sin^2(r/2)}{r^2}.
\]
The axis contributes one. Over an ordered field the planar numerator vanishes
exactly when \(\Sin r=0\) and \(\Cos r=1\), equivalently
\(r\in2\pi\Oz\). At such a point the whole perpendicular plane is killed and
the axis survives, proving the rank statement.
\end{proof}

The singular values of the right-trivialized derivative are
\begin{equation}\label{e3:so:eq:singularvalues}
 1,\qquad \frac{2\abs{\Sin(r/2)}}r,\qquad
          \frac{2\abs{\Sin(r/2)}}r.
\end{equation}
This follows directly from \cref{e3:so:eq:jacplane}, whose symmetric square is the
scalar in \cref{e3:so:eq:jacdet} on the perpendicular plane. Thus a noncritical
infinite generator can still be very poorly conditioned transversely. For
example, at \(r=\omega+\pi\) the transverse singular values are
\(2/(\omega+\pi)\), positive but infinitesimal.

At \(w=\omega e_1\), the image is \(I\) and the differential has rank one.
At the ordinary half-turn radius \(r=\pi\), by contrast, the differential is
invertible. The ambiguity of a principal logarithm at half-turns is a branch
issue, not a critical point of the exponential there. Also note the factor of
two relative to the radial unit-quaternion exponential: its imaginary parameter
has half the spatial rotation angle. The four-dimensional quaternion radial map
in \cite{RepoQuaternion} has a different domain and different rank counts.

No inverse function theorem in an unspecified surreal analytic category is
inferred solely from nonsingularity of \cref{e3:so:eq:jacobian}. The formulas and
ranks are direct differential statements. Nor are they chain rules for a
field derivation acting on the coefficients.

\section{Three topologies, and what the spin cover means}\label{e3:so:sec:topology}

This section separates the most important analytical ambiguity. Algebraic
identities do not depend on a topology, but connectedness, convergence, compactness,
and fundamental groups do.

\subsection{The full fine topology}

On \(\No^d\) use neighborhoods
\(\norm{x-a}_{\infty}<\rho\), with every positive surreal \(\rho\) allowed.
On matrices and \(G(\No)\) use the corresponding finite-coordinate topology.
For the full class these are neighborhood schemes, or ordinary topologies in a
suitably larger-universe formulation with the size restrictions retained.

\begin{theorem}[Set-sized discreteness in the full surreal topology]
\label{e3:so:thm:finediscrete}
Every set-sized subset of \(\No^d\), and hence of \(G(\No)\), is closed and
discrete. Every convergent set-indexed net is eventually equal to its limit.
Every set-indexed Cauchy net is eventually constant. Nevertheless \(G(\No)\)
has no isolated points.
\end{theorem}
\begin{proof}
For a point \(a\) and a set-sized subset \(S\), form the set of positive
distances \(\norm{s-a}_{\infty}\) for \(s\in S\setminus\{a\}\).
The separated surreal cut with left set \(\{0\}\) and that right set has a
positive separator \(\rho\) smaller than every listed distance. If the list
is empty, any positive \(\rho\) works. Its ball meets \(S\) in at most
\(\{a\}\). This proves both closedness and discreteness.

A set-indexed net has a set-sized range. Applying the preceding argument at
its limit proves eventual equality. For a Cauchy net, apply the cut argument
to all positive pairwise distances in its range; these form a set. A Cauchy
tail of diameter smaller than the resulting positive separator has only one
value. Finally, \(\cay(te_1)\to I\) as a neighborhood assertion with \(t\)
arbitrarily small in \(\No\), and \(\cay(te_1)\ne I\) for \(t\ne0\).
Translations give no isolated points anywhere in the group.
\end{proof}

The general size phenomenon is already emphasized in the repository's
foundations report \cite{RepoFoundations}; the proof here specializes it to
rotation coordinates. In a universe-relative version, ``set-sized'' means
small relative to the universe in which the whole surreal carrier is large.
It cannot be dropped from the statement.

In particular, an ordinary real interval admits no nonconstant fine-continuous
map into \(G(\No)\): its image is a set, is discrete, and must be connected.
An ordinary irrational-angle cyclic subgroup is a countable closed discrete
subset of \(G(\No)\), even though its standard part can be dense in an ordinary
real circle. Strongly summed exponential partial sums cannot converge in this
fine topology unless they are eventually equal to their limit. This is why the
support theorem, rather than a real convergence argument, was required.

\subsection{The intrinsic topology of a set-sized workspace}

Let \(F\supsetneq\R\) be a set-sized real closed field. Its \emph{intrinsic}
order topology permits only radii from \(F\). If \(F\subseteq\No\), this is
not its topology induced from the full surreal class: the latter is discrete
by \cref{e3:so:thm:finediscrete}, while the intrinsic order topology has no isolated
points. The distinction applies to \(G(F)\) as well.

Even intrinsically, \(F\) is totally separated. Its infinitesimal ideal
\(\mm_F\) is clopen: a positive infinitesimal radius proves openness at zero,
and a point with noninfinitesimal magnitude has a ball disjoint from
\(\mm_F\). Translates and nonzero scalar multiples of \(\mm_F\) are clopen.
For \(a\ne b\), the clopen set \(a+\abs{b-a}\mm_F\) contains \(a\) but not
\(b\). Coordinate projections therefore separate any two points of \(G(F)\)
by a clopen set. Ordinary real-parameter paths into \(G(F)\) are constant, and
its ordinary fundamental group at every point is trivial.

This topology is not locally compact. To verify it rather than merely infer it
from non-Archimedeanness, note that \(\R\subset F\) is closed and discrete:
near a real point choose an infinitesimal radius, near a finite nonreal point
choose less than its distance to its standard part, and near an infinite point
choose a ball avoiding all finite elements. Hence every subset of \(\R\) is
closed in \(F\). Every neighborhood of the identity in \(G(F)\) contains
\(\{\cay(ae_1/n):n\ge1\}\) for a sufficiently small \(a>0\) in \(F\).
This is an infinite closed discrete subset. It is closed even at the boundary
of the Cayley chart, since its trace stays a positive \(F\)-distance from
\(-1\), by \cref{e3:so:eq:caytrace}. A compact neighborhood would contain it as a
closed subset, which is impossible.

On a Hahn workspace, intrinsic valuation convergence is still meaningful and
can be nontrivial. In \(\R((t^\Q))\), for instance, \(t^n\to0\), since the
ordinary integers are cofinal in \(\Q\). In a higher-rank group that conclusion
can fail for a given positive exponent. These are statements about the
workspace, not convergence in the full surreal fine topology.

\subsection{The standard-part quotient topology is discrete}

The subgroup \(K_F\) is open and closed in the intrinsic topology, since it is
the intersection of finitely many conditions \(A_{ij}-I_{ij}\in\mm_F\).
Thus all fibers of standard part are open. It follows that the quotient topology
on
\begin{equation}\label{e3:so:eq:discretequotient}
 G(F)/K_F\cong\SO(3,\R)
\end{equation}
is \emph{discrete}, not the usual real Lie-group topology. Standard part into
the usual \(\SO(3,\R)\) is continuous, but it is not a quotient map onto that
usual topology. The constant section from the usual real Lie group is not
continuous. The same neighborhood conclusions hold for the full surreal class.

Accordingly, a topological reading of the split extension must use the discrete
topology on its real factor. Alternatively, one can deliberately choose the
coarser topology pulled back from the usual real group by standard part. That
topology has the usual real group as its quotient, but does not separate distinct
members of one infinitesimal fiber. These are different choices, not competing
proofs about one topology.

\subsection{Semialgebraic paths, compactness, and the spin fundamental group}

Over a real closed field \(F\), semialgebraic topology restricts maps and
homotopies to semialgebraic ones with parameter intervals \([0,1]_F\).
Closed and bounded semialgebraic sets are definably compact, and the matrix
equations show that \(G(F)\) and \(\Sp(1,F)\) have that property. This is not
ordinary compactness of all open covers.

There is an explicit semialgebraic path from the identity to every rotation.
Choose a unit quaternion lift \(q\) with nonnegative scalar part and put
\begin{equation}\label{e3:so:eq:sapath}
 q(u)=\frac{(1-u)+uq}{\sqrt{N((1-u)+uq)}},\qquad 0\le u\le1\text{ in }F.
\end{equation}
The numerator never vanishes: for \(0\le u<1\) its scalar part is strictly
positive, and at \(u=1\) it is \(q\). Its spin image is the
required path. Thus \(G(F)\) is semialgebraically path connected even when its
ordinary topological connected components are singletons. The parameter
interval and the allowed notion of homotopy have changed.

The spin map is a semialgebraic two-sheeted covering. Around the identity use
the non-half-turn open set and choose the lift with positive scalar part; the
other lift has negative scalar part. Both branches are continuous and
semialgebraic by the explicit square-root formulas, and translation supplies
such neighborhoods at every point. This is also a local two-sheeted covering
for the intrinsic topology of \(F\)-points.

\begin{theorem}[The category in which the familiar fundamental group survives]
\label{e3:so:thm:sahomotopy}
For every set-sized real closed field \(F\),
\[
 \pi_1^{\mathrm{sa}}(G(F),I)\cong\Z/2\Z,
 \qquad \pi_1^{\mathrm{sa}}(\Sp(1,F),1)=0.
\]
In the semialgebraic category the spin map is the simply connected double
cover. For a proper real closed extension of \(\R\), this is not the ordinary
fundamental group of the intrinsic topological space.
\end{theorem}
\begin{proof}
The equations of \(G\), the unit three-sphere, and the spin map are defined
over the rationals. The semialgebraic homotopy comparison and real-closed-field
extension theorem of Baro--Otero \cite[Corollaries 4.3--4.4]{BO} identifies their
semialgebraic homotopy groups with the ordinary homotopy groups of their real
realizations. Over \(\R\) the map is the antipodal double cover
\(S^3\to\mathbb{RP}^3\), giving the displayed groups. The local branches just
constructed identify the transferred covering. For a proper extension,
total separation has already shown that all ordinary loops are constant.
\end{proof}

Likewise, \(\pi_2^{\mathrm{sa}}G(F)=0\) and
\(\pi_3^{\mathrm{sa}}G(F)\cong\Z\), by the same comparison. A claim about a
``surreal universal cover'' is meaningful only after specifying this category
and, for the full class, the compatible universe or workspace interpretation.
The fine topology does not make \(G(\No)\) into an ordinary real Lie group.

\section{Worked examples at surreal scales}\label{e3:so:sec:examples}

\subsection{A tiny angle with a macroscopic displacement}

Take \(\eps=\omega^{-1}\) and
\begin{equation}\label{e3:so:eq:epsrotation}
 A=\cay(\eps e_3)=
 \begin{pmatrix}
 \dfrac{1-\eps^2}{1+\eps^2}&-\dfrac{2\eps}{1+\eps^2}&0\\[5pt]
 \dfrac{2\eps}{1+\eps^2}&\dfrac{1-\eps^2}{1+\eps^2}&0\\[5pt]
 0&0&1
 \end{pmatrix}.
\end{equation}
It belongs to \(K\), is not the identity, and has angle
\(2\atan\eps=2\eps-\tfrac23\eps^3+\cdots\), where the arctangent is the
strong local series. Acting on \(Le_1\) gives
\[
 A(Le_1)-Le_1=
 \left(-\frac{2L\eps^2}{1+\eps^2},\,
             \frac{2L\eps}{1+\eps^2},\,0\right)^T.
\]
For \(L=\omega\), the transverse displacement has standard part two. For
\(L=\omega^2\), it is of order \(\omega\), although the matrix still has
standard part \(I\). An infinitesimal relative error need not be a negligible
absolute error at an infinite spatial scale.

\begin{proposition}[Exact displacement formula]\label{e3:so:prop:displacement}
If \(A=\Rot(n;c,s)\) and \(x_{\perp}=x-(n\cdot x)n\), then
\begin{equation}\label{e3:so:eq:displacement}
 \norm{Ax-x}^2=2(1-c)\norm{x_{\perp}}^2.
\end{equation}
For a finite angle \(\theta\), this is
\(\norm{Ax-x}=2\abs{\sin(\theta/2)}\norm{x_{\perp}}\).
\end{proposition}
\begin{proof}
The axis component cancels. On the plane,
\((A-I)x_{\perp}=(c-1)x_{\perp}+s(n\times x_{\perp})\).
The two vectors are perpendicular with equal norms, so the squared norm is
\(((c-1)^2+s^2)\norm{x_{\perp}}^2=2(1-c)\norm{x_{\perp}}^2\).
\end{proof}

\subsection{A near-half-turn with an infinite Cayley coordinate}

Let \(B=\cay(\omega e_3)\). Then
\[
 c=-1+\frac2{1+\omega^2},\qquad s=\frac{2\omega}{1+\omega^2},
 \qquad \theta=\pi-2\atan(\omega^{-1}).
\]
This is not a half-turn, although its standard part is one. The matrix entries
are bounded; its Cayley coordinate alone is infinite. Its quaternion lift has
scalar part \(1/\sqrt{1+\omega^2}\) and third imaginary part
\(\omega/\sqrt{1+\omega^2}\). The largest-coordinate extraction algorithm
selects the latter rather than dividing by a tiny scalar coordinate.

\subsection{Infinitesimally different axes}

The unit vector
\[
 n=\frac{(1,\eps,\eps^2)^T}{\sqrt{1+\eps^2+\eps^4}}
\]
defines a half-turn \(H_n=2nn^T-I\) with the same standard part as \(H_{e_1}\).
The half-turns are distinct because their fixed lines differ. Their difference
is visible already in the off-diagonal entry
\((H_n)_{12}=2\eps/(1+\eps^2+\eps^4)\). This is axis information that a real
standard-part calculation discards.

\subsection{Noncommutativity first appears at a finer scale}

For positive infinitesimals \(\eps,\delta\), let
\(U=\cay(\eps e_1)\), \(V=\cay(\delta e_2)\). Then
\begin{equation}\label{e3:so:eq:twoscale}
 \comm{U}{V}=I+4\eps\delta\widehat{e_3}
       +\text{terms divisible by }\eps\delta
          \text{ of total degree at least three}.
\end{equation}
In valuation language, the remainder has valuation at least
\(\min(2v(\eps)+v(\delta),v(\eps)+2v(\delta))\), strictly above the displayed
term. This follows either by rational multiplication or by the BCH bracket.
The two rotations commute to first order but not exactly.

\subsection{Infinite generators and local sensitivity}

Put \(w=\omega e_1\) and \(z=\pi e_2\). Equation \cref{e3:so:eq:globalexp} gives
\(\mathcal E(\widehat w)=I\) and \(\mathcal E(\widehat z)=H_{e_2}\). Yet
\[
 \norm{w+z}=\sqrt{\omega^2+\pi^2}
 =\omega+\frac{\pi^2}{2\omega}+\smallO(\omega^{-3}),
\]
so the global phase removes \(\omega\) and leaves an infinitesimal angle.
Consequently
\[
 \mathcal E(\widehat{w+z})=I+\frac{\pi^2}{2\omega}\widehat{e_1}
                        +\smallO(\omega^{-2}),
\]
not \(H_{e_2}\). Here \(\smallO(\omega^{-k})\) means an entrywise bound by a
finite real constant times that scale; it is not an unspecified real limit.
This illustrates why global radial exponentiation is not additive across
noncommuting directions, even when one generator has trivial image.

\subsection{A nonreal copy of the octahedral group}

Let \(H_0\) be the 24 signed permutation matrices of determinant one, and set
\(C=\cay(\eps e_1+\eps^2e_2)\). Then
\[
 H=CH_0C^{-1}
\]
is an exact 24-element subgroup with generally nonreal entries. Its standard
part is \(H_0\), and conjugation by \(C^{-1}\in K\) restores the ordinary
real group. The multiplication table is unchanged; the coordinate realization
has been deformed at several infinitesimal scales.

\part{Computation, limits, and provenance}\label{e3:part:limits}

\section{Exact computation and formalization routes}\label{e3:sec:implementation}

The two sources give complementary implementation plans, one for the geometry
and one for the rotation group. Both are printed; neither describes existing
software.

\subsection{A coordinate-to-invariants algorithm (source 04)}\label{e3:sub:impl04}

For three proposed vertices $P_1,P_2,P_3$ on a sphere of center $O$
and radius $R>0$, an exact workflow is as follows.
\begin{enumerate}[label=\arabic*.,leftmargin=2em]
\item Verify $\norm{P_i-O}^2=R^2$ and set $u_i=(P_i-O)/R$.
Compute $x=u_2\cdot u_3$, $y=u_3\cdot u_1$, $z=u_1\cdot u_2$,
and $\tau=u_1\cdot(u_2\times u_3)$.
\item Check $\tau\ne0$ for a proper triangle. Retain its sign for
orientation, and put $d=|\tau|$. Verify the redundant identity
$d^2=1+2xyz-x^2-y^2-z^2$ as a useful certificate.
\item Compute the side angles with
$a=\atantwo(\norm{u_2\times u_3},x)$ and cyclic variants;
compute $\alpha=\atantwo(d,x-yz)$ and cyclic variants.
All these calls have strictly positive first argument.
\item Compute $M=1+x+y+z$ and $E=2\atantwo(d,M)$.
Return physical side lengths $Ra,Rb,Rc$ and area $R^2E$.
Centers, polar vertices, and transport matrices are obtained from the
explicit formulas above, with each denominator checked before division.
\end{enumerate}

This is a mathematical algorithm over a field with exact operations,
order decisions, positive square roots, and finite-angle evaluation.
It is not a claim that every surreal input has a finite machine
representation. A general Conway normal form has set-sized support,
and a generic element need not come with effective coefficient or
support oracles. An actual computer algebra implementation must name
its representable subclass and its supported comparison procedure.

For exact finite descriptions, it is useful to separate an algebraic
certificate from a displayed numerical approximation. For example,
$\tau\ne0$ can be certified by a leading coefficient even when
$\st(\tau)=0$. Near tangency, an exact discriminant sign should not
be replaced by an ordinary floating-point tolerance. Truncated Hahn
expansions must retain enough support to decide the leading nonzero
term of the predicate being tested; an unresolved cancellation is an
unknown sign, not a zero.

\subsection{An implementation hierarchy for rotations (source 09)}\label{e3:sub:impl09}

There is no finite data structure representing every surreal number with a
decidable exact equality test. A practical rotation package must name its scalar
backend. The finite algebra in this report works over any exact ordered field
with the required square roots; analytic operations require more.

For a backend of rational functions in one formal positive infinitesimal,
Cayley coordinates are especially convenient: composition, inverse, and the
action on vectors are rational formulas. Unit quaternions can be kept
projectively, avoiding square roots until normalization is needed. A bounded
Puiseux or Hahn-series backend can additionally track leading exponents,
standard parts, and finite-order BCH truncations. General supports need an
explicit effective representation; a mathematical support theorem is not an
algorithm for deciding equality of arbitrary input names.

A robust division of responsibilities is the following.
\begin{enumerate}[label=\textup{\arabic*.}]
\item The scalar layer supplies exact arithmetic, order, positive square roots,
      and, when supported, valuation and standard part.
\item The finite rotation layer supplies matrix, quaternion, algebraic
      axis-angle, and Cayley representations, with all exceptional branches.
\item The supported-series layer supplies infinitesimal exp/log and BCH with
      declared error valuations and support certificates.
\item The phase layer distinguishes finite angles from the chosen global
      Ehrlich--Kaplan phase; it never treats an infinite argument as an
      uncentered convergent Taylor input.
\end{enumerate}
The exact rational law on \(\mm^3\) can often replace BCH entirely. BCH is
useful for logarithmic generators and graded errors, not necessary merely to
multiply two rotations.

\subsection{A layered Lean development for the geometry (source 04)}\label{e3:sub:lean04}

At source~04's pin, the repository's \texttt{Surreal/Algebra/Geometry.lean}
works with two-coordinate complexifications and field-valued moduli
\cite{RepoGeometry}. Its approach suggests a clean separation for the
three-dimensional extension. The following are proposed layers, not
assertions that these modules already exist or have been compiled. At the
merge pin \texttt{50cb709} no such module exists.

\begin{center}\small
\begin{tabular}{L{0.22\textwidth}L{0.30\textwidth}L{0.39\textwidth}}
\toprule
Layer & Assumptions & Principal obligations\\
\midrule
Vector algebra & Commutative ring & Dot/cross identities; determinants; Gram polynomials\\
Ordered geometry & Ordered field; positive square roots & Norm inequalities; projections; affine incidence; intersections\\
Orthogonal algebra & Real closed field when needed & Principal axes; axis--angle cosine/sine data; quaternion action\\
Spherical algebra & Unit vectors; nonzero determinant & Tangent and normal formulas; cosine/sine laws; reconstruction\\
Finite angles & Compatible restricted-analytic structure & Phase bijection; inverse branches; angle inequalities\\
Area and transport & Proper minor-arc polygons & Excess branches; finite additivity; edge transport; holonomy\\
Scale analysis & Valuation and standard part & Noncancellation; asymptotic bounds; degeneration certificates\\
\bottomrule
\end{tabular}
\end{center}

An appropriate vector type is $\operatorname{Fin}(3)\to F$ or a
three-coordinate structure. Length should initially be an $F$-valued
function, not a real-valued \texttt{NormedSpace} norm. Theorems about
matrices and finite sums can be proved before introducing any surreal
construction. Angles should not enter the algebraic spherical-law layer
until the circle-phase bridge is available. Proper-class statements are
then obtained externally from compatible set-sized workspaces, rather
than by pretending that all of $\No$ is a small type in one universe.

The exceptional SSA family in \cref{e3:subsec:SSA} is a useful specification
test: a solver returning only a list of at most two triangles is
mathematically incomplete unless it explicitly excludes the case
$\cos b=\sin b\cos\alpha=0$. Its output type should distinguish no
solution, finitely many solutions, and a parametrized family.

\subsection{A Lean development plan for the rotation group (source 09)}\label{e3:sub:lean09}

The repository's formalization ledger is the authority for which existing
statements have actually been matched to compiled Lean declarations
\cite{RepoLedger}. This report does not add a compiled Lean development, and
its proofs must not be marked proved in that ledger solely because they appear
here or pass symbolic tests.

The natural finite-algebra modules would begin with \(F^3\), the hat map, and
\cref{e3:so:eq:hatpowers,e3:so:eq:hatbracket}, then prove the distance-isometry theorem and
axis theorem. Next come the Hamilton action, spin kernel and explicit
surjectivity, the matrix Cayley equivalence, and the centralizer cases. A
standard-part interface then gives the split extension and the kernel chart.
Finite subgroup rigidity can be formalized using a finite sum and the positive
matrix-root interface, without invoking Hahn series.

The torsion-free and unique-divisibility proofs are also finite algebra over a
real closed extension of \(\R\). The single-commutator construction reduces
to one checked quaternion identity, positive fourth roots, and the already
proved transitivity theorem. These are particularly attractive next targets:
they deliver significant group structure without depending on a general
surreal analytic inverse function theorem.

For the analytic modules, formal noncommutative series must be evaluated only
after establishing joint strong summability. The derivative theorem requires a
fine Fr\'echet interface in three coordinates and local quadratic remainders;
it cannot be obtained merely by substituting surreal symbols into a theorem
whose hypotheses require a real Banach space. Semialgebraic fundamental groups
are a further, independent formalization layer.

\subsection{The companion checks}\label{e3:sub:checks}

Each source came with an exact SymPy script. Both are kept byte-identical, under
the prefixed names given here.

\texttt{code/04-three-space-verify\_identities.py} (source~04) checks selected
identities and expansions: the vector and Gram identities, the
spherical-excess polynomial, the quaternion holonomy numerator, the chart
metrics, the local curvature coefficient, and the worked examples. Its recorded
run, \texttt{data/04-three-space-verification.txt}, passed all 19 test groups.
These checks can detect a sign, coefficient, or transcription error in the
displayed formulas. They do not prove the analytic transfer theorem, establish
Hahn summability, discharge every branch condition, or formally verify the
geometry. Those distinctions are part of the mathematical specification, not
merely software caveats.

\texttt{code/09-so3-verify.py} (source~09) checks the hat identities, quaternion
multiplication and rotation formulas, orthogonality and determinant of the
Cayley chart, its composition law, the scalar commutator identity used in
\cref{e3:so:thm:perfect}, selected leading commutators, and the Jacobian
calculation in an axis-adapted frame. It also checks all 576 products of the
24-element octahedral group and an exact rational averaging intertwiner. Its
recorded run, \texttt{data/09-so3-verification.txt}, passed all 24 checks.
Rational specialization tests supplement, but do not replace, symbolic
identities. These checks can catch sign errors, factor-of-two errors, and
incorrect matrix branches. They do not prove real closedness of $\No$, the
Neumann support lemma, a global topology theorem, or a Lean kernel judgment.
Those claims rest on the mathematical proofs and explicitly imported results
stated in the text.

Both recorded runs used Python~3.13.5 and SymPy~1.14.0. For this merge both
scripts were rerun with Python~3.14.4 and SymPy~1.14.0. The rerun reproduced
19/19 groups and 24/24 checks; only the version line and the timing line
differ. Neither script writes a file.

\section{Non-claims and limits, by source}\label{e3:sec:nonclaims}

A merged report must not silently upgrade either source. This section records
every limitation, disclaimer and priority caveat stated by either manuscript.
Most also appear at the point of use.

\subsection{Source 04 (Part I)}\label{e3:sub:nonclaims04}
\begin{enumerate}[label=\textup{G\arabic*.},leftmargin=2.6em]
\item No priority claim is made, and no named open problem is claimed to be
 resolved. This covers in particular the uniform gnomonic estimate
 (\cref{e3:thm:local-area}), the exceptional-family classification
 (\cref{e3:prop:SSA}), and their combination into a surreal development.
 Repeated classical identities are not presented as new results merely
 because their arguments are surreal.
\item Established surreal analytic machinery and classical identities are
 cited: \cite{vdDE,vdDEErr,EK,Neumann,DLMF,Oosterom}. The proofs supply the
 precise extensions used.
\item The symbolic checks test selected polynomial identities and formal
 coefficients only. They do not verify the whole development or its
 foundations (\cref{e3:sub:checks}).
\item No Lean formalization is claimed. The layers of \cref{e3:sub:lean04} are
 proposals, not assertions that the modules exist or have been compiled.
\item Real-valued metric completeness, ordinary Riemann integration, and a
 global oscillatory differential equation are not imported without additional
 hypotheses.
\item The finite theory does not assert that no distinguished global sine or
 cosine exists (\cref{e3:rem:infinite-input}). It needs no phase for infinite
 inputs. The zero-centred sine series at $\omega$ is not strongly summable,
 whatever global definition is chosen.
\item The distance is field-valued, not an ordinary real-valued metric, and
 $F^3$ is not made a Banach space over $\R$.
\item The ``double cover'' of \cref{e3:sec:rotations} is an algebraic two-to-one
 homomorphism, not a topological covering assertion in the fine topology.
\item The spherical-coordinate metric and Jacobian are finite identities in
 restricted-analytic coordinate derivatives. They are not an integral
 transformation theorem for every function on $F^3$.
\item The minimizing property of minor arcs concerns finite chains of minor
 great-circle arcs. No supremum over arbitrary partitions, no real-valued length
 metric, and no unrestricted class of nondefinable curves is introduced.
\item Area is a finitely additive functional on geodesic polygons. The cap
 formula extends one specific analytic formula. It asserts neither that
 arbitrary subsets have a surreal measure nor that Darboux sums define it, and no
 countably additive area measure is inferred.
\item Derivatives in \cref{e3:sec:differential} are taken with respect to
 external chart or curve parameters. They are not an intrinsic derivation of
 $\No$, and no existence theorem for arbitrary differential equations is
 invoked.
\item Parallel transport is canonical by explicit analytic continuation. It is
 unique within the definable restricted-analytic class, not among arbitrary
 intrinsically differentiable fields.
\item The normal-coordinate estimate \cref{e3:eq:normal-limit} is absolute. It
 claims no small relative error when $p-q$ is exceptionally small.
\item The coordinate-to-invariants algorithm does not claim that every surreal
 input has a finite machine representation. An implementation must name its
 representable subclass and comparison procedure.
\item Completeness, ordinary topological compactness, existence for arbitrary
 smooth ODEs, and a supremum-based length for every curve are not claimed. Finite
 dimension does not permit unrestricted summation of arbitrary families of
 surreals. Ordinary-size angles do not imply ordinary-size arc lengths, and a
 global phase at infinite arguments would be an additional convention.
\item The repository comparison is not an exhaustive audit of every file, or
 of the published work on non-Archimedean geometry. The repository is a
 research source, not a substitute for published foundations.
\item No completed repository integration, general-purpose surreal computer
 algebra, or formal verification is asserted.
\end{enumerate}

\subsection{Source 09 (Part II)}\label{e3:sub:nonclaims09}
\begin{enumerate}[label=\textup{S\arabic*.},leftmargin=2.6em]
\item The manuscript is an AI-assisted research exposition, subject to
 mathematical review. It makes no claim of new priority and no claim of
 complete machine-checked formalization.
\item No new priority is claimed for classical quaternion theory or for the
 perfectness of the infinitesimal rotation group. Bays--Peterzil record that
 every element is a commutator \cite[\S3.1]{BP}, and \cref{e3:so:thm:perfect}
 gives an elementary exact construction of that fact.
\item The quaternion spin description, standard-part reduction of unit
 quaternions, infinitesimal exponential and logarithm, and valuation filtration
 in the surquaternion report, finite angles in the trigonometry report, and the
 foundations report's distinction of strong summation from fine convergence are
 starting points, not discoveries of the source.
\item The localization result, \cref{e3:so:prop:localization}, does not
 assert that the subfield is closed under strong sums, or that its intrinsic
 topology is the one induced from $\No$.
\item The homogeneous-space statement \cref{e3:so:eq:torus} is algebraic and
 semialgebraic, not a statement about ordinary real manifolds.
\item $G(F)\cong\mathbb P^3(F)$ is a statement about $F$-points, not an
 isomorphism of algebraic varieties.
\item Connectedness and absolute simplicity of the algebraic group imply
 neither abstract simplicity of $G(F)$ nor connectedness of its fine
 topology.
\item Unique divisibility of $K_F$ does not make it a $\Q$-vector space.
\item \Cref{e3:so:thm:finiterigidity} gives conjugate realizations. It does not
 say that finite subgroups are literally constant matrices, and it says nothing
 about infinite or set-generated subgroups.
\item The dual number $\kappa$ is formal, not a surreal infinitesimal. For an
 angular-velocity equation, existence and uniqueness of solutions is a separate
 question. The Berarducci--Mantova derivation is not the variable derivative of
 \cref{e3:so:sec:globalexp}.
\item The BCH truncation valuations are finite-order error certificates, not
 convergence through all positive surreal tolerances. The support lemma does
 not assert that $n\gamma$ is cofinal.
\item \Cref{e3:so:thm:graded} gives an inclusion of filtration levels, not
 equality of valuations in every case. The BCH description does not make $K$
 pronilpotent in an unspecified topology.
\item The Bays--Peterzil bi-interpretability theorem is imported, not proved,
 and it concerns set-sized structures. It is not automatically transferred to
 the full class.
\item The Ehrlich--Kaplan phase is canonical for its specified normalization.
 This is not an assertion that every conceivable global phase convention agrees
 with it.
\item The identity fiber of the global rotation exponential is not a kernel,
 and that exponential is not a homomorphism on the additive Lie algebra.
\item No inverse function theorem is inferred from nonsingularity of
 \cref{e3:so:eq:jacobian}, and the derivative formulas are not chain rules for
 a field derivation.
\item The size restriction in \cref{e3:so:thm:finediscrete} cannot be dropped.
 The spin fundamental group $\Z/2\Z$ is semialgebraic, not the ordinary
 fundamental group of the intrinsic or fine space. Definable compactness is not
 ordinary compactness. The fine topology does not make $G(\No)$ an ordinary real
 Lie group, and a ``surreal universal cover'' is meaningful only after a
 category is specified.
\item No finite data structure represents every surreal number with decidable
 equality, and a support theorem is not an equality algorithm.
\item No compiled Lean development is added. The proofs must not be marked
 proved in the formalization ledger solely because they appear here or pass
 symbolic tests.
\item The checks do not prove real closedness of $\No$, the Neumann support
 lemma, a topology theorem, or a Lean kernel judgment (\cref{e3:sub:checks}).
\item There is no claim to classify all abstract subgroups or automorphisms,
 to construct a Haar measure on the full class, or to obtain a physical model
 merely from a change of scalars.
\item The repository was not rebuilt for the source, and the source changed no
 repository file.
\end{enumerate}

\subsection{The merge}\label{e3:sub:nonclaimsmerge}
\begin{enumerate}[label=\textup{M\arabic*.},leftmargin=2.6em]
\item The merge adds no mathematical result. Each statement comes from one of
 the sources, and each statement proved by both is marked in
 \cref{e3:sub:provenance}. Merge-written remarks only relate statements to
 each other or to other reports (for example the factor-two comparison after
 \cref{e3:so:thm:graded}). \Cref{e3:thm:finiteaxis} and \cref{e3:thm:isometries}
 combine clauses that the two sources stated separately; the proofs printed
 are the sources' proofs.
\item The renamings of \cref{e3:sub:notation} change symbols, not statements.
\item \Cref{e3:thm:spin} is cited from the surquaternion report and not reproved.
 Its proof there is outside the scope of this report's review.
\item The report has not been refereed. None of its labels is mapped to a Lean
 declaration.
\end{enumerate}

\section{Conclusion and directions}\label{e3:sec:conclusion}

The first two paragraphs are source~04's conclusion, and the rest is source~09's.

Three-dimensional surreal geometry has a large, rigorous core that is
both familiar and genuinely more informative than its real standard-part
picture. Affine incidence, orthogonality, quadrics, rotations, and finite
volumes are ordered-field algebra. Canonical finite angles supply the
usual trigonometric language for all directions, even when coordinates
and physical lengths are infinite. Spherical geometry then admits exact
sine and cosine laws, reconstruction criteria, centers, a finite-polygon
area functional, and a curvature--holonomy identity.

The indispensable extra care concerns branch choices, degeneracy, and
the meaning of analysis. A zero standard part is not a zero determinant;
a small determinant need not mean small spherical area; a locally
constant, intrinsically differentiable function need not be globally constant; and a
formal integral symbol is not a measure. With those distinctions made
explicit, infinitesimal geometry becomes a calculus of exact scales
rather than a collection of heuristic limiting pictures.

The correct analogue of the ordinary spatial rotation group is not a new
multiplication invented for surreal coordinates: it is the positive-definite
special orthogonal class group \(\SO(3,\No)\). Its axes, algebraic angles,
quaternion double cover, and rational charts have exact finite formulas.
Every matrix entry is bounded by one, but the entries can retain arbitrarily
fine surreal information.

The strongest structural contrast with the ordinary group is the extension
\(K\rtimes\SO(3,\R)\). The real quotient does not exhaust the geometry.
The kernel has a full hierarchy of cross-product layers, yet the entire kernel
is perfect and uniquely divisible. It contains no finite torsion and no
nontrivial finite quotient. Finite subgroups have no genuinely new forms:
finite averaging and polar normalization conjugate them to real subgroups.

Finite angles suffice to represent every rotation. The canonical global phase
also exists, but its period class is \(2\pi\Oz\), not just ordinary multiples
of \(2\pi\). The resulting radial rotation exponential has explicitly described
fibers and an exact Jacobian. Infinite generators can be critical or severely
ill-conditioned without producing unbounded rotation matrices.

Finally, the phrase ``like \(\SO(3)\)'' has a category-dependent answer. The
algebraic group is three-dimensional and connected. Over a set-sized real closed
field it is semialgebraically connected and definably compact, with spin
fundamental group \(\Z/2\Z\). In the non-Archimedean intrinsic topology it is
totally separated and not locally compact. In the full surreal fine topology,
every set-sized subset is closed and discrete. None of these assertions
contradicts the others.

Further work should therefore specify a category before posing a problem.
Substantial directions include normal-subgroup classification compatible with
convex valuations; effective multi-scale rotation algorithms with support and
conditioning certificates; integration of the finite proofs into the
repository's Lean development; and comparison of semialgebraic, valued-field,
and fine analytic representation theories. We do not claim to classify all
abstract subgroups or automorphisms, construct a Haar measure on the full class,
or obtain a physical model merely from a change of scalars.

\appendix
\section{Formula and convention reference for Part II}\label{e3:so:app:formulas}

\renewcommand{\arraystretch}{1.25}
\begin{longtable}{L{0.27\textwidth}L{0.67\textwidth}}
\toprule
Object & Convention or formula\\
\midrule
\endfirsthead
\toprule Object & Convention or formula\\\midrule
\endhead
Rotation group & \(G(F)=\{A:A^TA=I,\det A=1\}\), acting on column vectors.\\
Hat map & \(\widehat u\,v=u\times v\), with the standard right-handed cross product.\\
Rodrigues & \(\Rot(n;c,s)=cI+(1-c)nn^T+s\widehat n\).\\
Half-turn & \(H_n=2nn^T-I\); its centralizer is a line stabilizer, not just an axis circle.\\
Spin map & \(p(a,u)=(a^2-\norm u^2)I+2uu^T+2a\widehat u\), for \(a^2+\norm u^2=1\).\\
Cayley chart & \(\cay(t)=((1-t\cdot t)I+2tt^T+2\widehat t)/(1+t\cdot t)\).\\
Cayley product & \(t\star u=(t+u+t\times u)/(1-t\cdot u)\), where the denominator is nonzero.\\
Infinitesimal kernel & \(K=\ker\st\), identified with all \(\mm^3\) by the Cayley chart.\\
Logarithmic leading term & \(\Log_{\mm}\cay(t)=2\widehat t+\text{higher terms}\).\\
Valuation & \(v(\omega^{-\gamma})=\gamma\); infinitesimals have \(v>0\).\\
Graded coordinate & \(\ell_\gamma(U)=2\st(\omega^\gamma t(U))\).\\
Group commutator & \(\comm{U}{V}=UVU^{-1}V^{-1}\).\\
Lie bracket & \([\widehat u,\widehat v]=\widehat{u\times v}\).\\
Finite phase period & \(\ker\cis=2\pi\Z\), on \(\OO\).\\
Global phase period & \(\ker\cis_{\mathrm{EK}}=2\pi\Oz\), on \(\No\).\\
Global exponential & Radial map \(\mathcal E\) of \cref{e3:so:eq:globalexp}; its identity fiber is not an additive kernel.\\
Fine Jacobian & Right-trivialized as \((d\mathcal E_w h)\mathcal E_w^{-1}=\widehat{\mathcal J(w)h}\).\\
Critical radii & \(r\in(2\pi\Oz)_{>0}\), rank one; \(r=0\) is regular.\\
\bottomrule
\end{longtable}
\renewcommand{\arraystretch}{1}

\section{Source, dependency, and verification ledger}\label{e3:app:ledger}

\subsection{Pins and files}

\begin{center}\small
\begin{tabular}{L{0.2\textwidth}L{0.72\textwidth}}
\toprule
Item & Value\\
\midrule
Source 04 pin & \texttt{d22a5b35d5b3040e870c3cfd6d1c8f7259e094b0}\\
Source 09 pin & \texttt{dcf86662b574988e75d08d515d3a2a5ba7bbc0d7}\\
Placement & \texttt{5fe7f8d} (both manuscripts placed, source~04's text as the
 starting file)\\
Merge written at & \texttt{50cb709}\\
Source 04 files & \texttt{code/04-three-space-verify\_identities.py},
 \texttt{data/04-three-space-verification.txt},
 \texttt{data/04-three-space-requirements.txt}\\
Source 09 files & \texttt{code/09-so3-verify.py},
 \texttt{data/09-so3-verification.txt},
 \texttt{data/09-so3-requirements.txt}\\
\bottomrule
\end{tabular}
\end{center}
The manuscripts themselves, their PDFs, READMEs and source~09's checksum list
are not shipped. The code and data files are byte-identical to the delivered
ones.

\subsection{Source 04: inspected sources and dependencies}

At its pin, source~04 specifically inspected the trigonometry report and the
Lean module \texttt{Surreal/Algebra/Geometry.lean}. The repository is a
research source, not a substitute for published foundations. The source
comparison is not an exhaustive audit of every file or of all published
work on non-Archimedean geometry.

The main dependencies of Part~I are these. Real closedness and the
restricted-analytic Hahn-field model theorem are imported from
\cite{vdDE}, with support-combinatorial cautions informed by its erratum
\cite{vdDEErr}. The canonical global extension available from an initial
integer part is attributed to \cite{EK}; Part~I needs only its finite
restriction. The classical spherical identities are recorded in
\cite{DLMF}; their algebraic proofs and sign conditions are provided
here. The solid-angle expression is classical and attributed to
\cite{Oosterom}. No priority claim is made for the remaining formulations,
including the uniform gnomonic estimate, the exceptional-family
classification, or their combination into a surreal development.
The proofs are mathematical proofs, not Lean proof terms.

\subsection{Source 09: inspected sources and dependency ledger}

At its pin, source~09 reviewed the documentation index, the introductory
status of the formalization ledger, the surquaternion README and the relevant
surquaternion source sections. Its bibliography separates the report
directories identified through that index from the source sections actually
inspected. It did not rebuild the repository.

\begin{longtable}{L{0.29\textwidth}L{0.36\textwidth}L{0.26\textwidth}}
\toprule
Result family & Essential assumptions & Evidence in this report\\
\midrule
\endfirsthead
\toprule Result family & Essential assumptions & Evidence in this report\\\midrule
\endhead
Isometries, axes, spin, Cayley, centralizers & Real closed ordered field; finite matrices & Direct finite proofs and selected symbolic identities; spin cited from \texttt{squat:thm:rotations}.\\
Standard-part splitting & Real closed \(F\supseteq\R\) & Direct ring-homomorphism proof.\\
Torsion-free and unique divisibility & Real closed \(F\supseteq\R\) & Algebraic circle and root proofs.\\
Perfect infinitesimal group & Proper real closed extension; fourth roots & Explicit commutator construction; known precedent \cite{BP}.\\
Finite-subgroup rigidity & Finite group, real closed \(F\supseteq\R\) & Averaging and positive polar factor, fully proved.\\
Strong exp/log and BCH & Set-supported normal forms or a divisible Hahn workspace & Support proof using Higman's lemma; formal identities.\\
Global phase & Ehrlich--Kaplan initial-integer-part normalization & Imported normalization, explicitly expressed and analyzed.\\
Rotation exponential differential & Fine variable derivatives and local Taylor remainder bounds & Direct radial/perpendicular calculation.\\
Full fine discreteness & Full surreal cut property; small index sets & Direct separating-radius proof.\\
Semialgebraic homotopy & Set-sized real closed field & Imported comparison theorem \cite{BO} and explicit spin branches.\\
Machine formalization & A compiled proof assistant development & Not supplied; no such verification is claimed.\\
\bottomrule
\end{longtable}

\begin{thebibliography}{99}
\raggedright

\bibitem{vdDE}
L. van den Dries and P. Ehrlich,
\emph{Fields of surreal numbers and exponentiation},
Fundamenta Mathematicae \textbf{167} (2001), no.~2, 173--188.
\href{https://doi.org/10.4064/fm167-2-3}{doi:10.4064/fm167-2-3}.

\bibitem{vdDEErr}
L. van den Dries and P. Ehrlich,
\emph{Erratum to ``Fields of surreal numbers and exponentiation''},
Fundamenta Mathematicae \textbf{168} (2001), no.~3, 295--297.
\href{https://doi.org/10.4064/fm168-3-5}{doi:10.4064/fm168-3-5}.

\bibitem{EK}
P. Ehrlich and E. Kaplan,
\emph{Surreal ordered exponential fields},
Journal of Symbolic Logic \textbf{86} (2021), no.~3, 1066--1115,
especially \S11.
\href{https://doi.org/10.1017/jsl.2021.59}{doi:10.1017/jsl.2021.59}.
See also the citation-correction erratum, \textbf{87} (2022), no.~2, 871,
\href{https://doi.org/10.1017/jsl.2022.5}{doi:10.1017/jsl.2022.5}.
Source~09 cites the preprint version,
\url{https://arxiv.org/abs/2002.07739v3}.

\bibitem{Neumann}
B. H. Neumann,
\emph{On ordered division rings},
Transactions of the American Mathematical Society \textbf{66} (1949),
202--252.
\href{https://www.jstor.org/stable/1990552}{Primary journal archive}.

\bibitem{Higman}
G. Higman,
\emph{Ordering by divisibility in abstract algebras},
Proceedings of the London Mathematical Society (3) \textbf{2} (1952), 326--336.
\href{https://doi.org/10.1112/plms/s3-2.1.326}{doi:10.1112/plms/s3-2.1.326}.

\bibitem{DLMF}
F. W. J. Olver et al., eds.,
\emph{NIST Digital Library of Mathematical Functions},
\S4.42, ``Solution of Triangles,'' especially \S4.42(iii),
``Spherical Triangles.''
\href{https://dlmf.nist.gov/4.42}{Official online section},
consulted September 22, 2026.

\bibitem{Oosterom}
A. van Oosterom and J. Strackee,
\emph{The solid angle of a plane triangle},
IEEE Transactions on Biomedical Engineering \textbf{BME-30} (1983),
no.~2, 125--126.
\href{https://doi.org/10.1109/TBME.1983.325207}{doi:10.1109/TBME.1983.325207}.

\bibitem{Voight}
J. Voight,
\emph{Quaternion Algebras},
Graduate Texts in Mathematics \textbf{288}, Springer, 2021.
Author's official and corrected editions:
\url{https://jvoight.github.io/quat.html}.

\bibitem{BP}
M. Bays and Y. Peterzil,
\emph{Definability in the group of infinitesimals of a compact Lie group},
Confluentes Mathematici \textbf{11}, no.~2 (2019), 3--23.
\href{https://doi.org/10.5802/cml.58}{doi:10.5802/cml.58}.
\url{https://arxiv.org/abs/1901.10831}.

\bibitem{BO}
E. Baro and M. Otero,
\emph{On o-minimal homotopy groups},
Quarterly Journal of Mathematics \textbf{61}, no.~3 (2010), 275--289.
\href{https://doi.org/10.1093/qmath/hap011}{doi:10.1093/qmath/hap011}.
\url{https://arxiv.org/abs/0810.0365}.

\bibitem{RepoTrig}
\repo,
\emph{Trigonometry on the Surcomplex Plane: Canonical Phases,
Arbitrary-Scale Geometry, and Infinitesimal Degeneration},
\path{docs/surcomplex/trigonometry/article.tex}.
Source~04 inspected it at
\href{https://github.com/VladimirReshetnikov/Surreal/blob/d22a5b35d5b3040e870c3cfd6d1c8f7259e094b0/docs/surcomplex/trigonometry/article.tex}{\texttt{d22a5b35d5b3}};
source~09 identified it through the documentation index at
\href{https://github.com/VladimirReshetnikov/Surreal/tree/dcf86662b574988e75d08d515d3a2a5ba7bbc0d7/docs/surcomplex/trigonometry}{\texttt{dcf86662b574}}.
Labels cited here are those present at \texttt{50cb709}.

\bibitem{RepoGeometry}
\repo,
\emph{Finite geometry over an ordered base field},
\path{Surreal/Algebra/Geometry.lean}, inspected by source~04 at
\href{https://github.com/VladimirReshetnikov/Surreal/blob/d22a5b35d5b3040e870c3cfd6d1c8f7259e094b0/Surreal/Algebra/Geometry.lean}{\texttt{d22a5b35d5b3}}.

\bibitem{RepoQuaternion}
\repo,
\emph{Surquaternions: Algebra, Infinitesimal Geometry, and Support-Controlled Analysis},
\path{docs/surquaternions/surquaternions/article.tex}.
Source~09 inspected the README and the relevant algebraic and phase sections at
\href{https://github.com/VladimirReshetnikov/Surreal/blob/dcf86662b574988e75d08d515d3a2a5ba7bbc0d7/docs/surquaternions/surquaternions/article.tex}{\texttt{dcf86662b574}}.
Theorem numbers cited here are those at \texttt{50cb709}.

\bibitem{RepoFoundations}
\repo,
\emph{Foundations},
\path{docs/foundations-and-computation/foundations/}; topology and size
conventions, identified by source~09 through the documentation index at
\href{https://github.com/VladimirReshetnikov/Surreal/tree/dcf86662b574988e75d08d515d3a2a5ba7bbc0d7/docs/foundations-and-computation/foundations}{\texttt{dcf86662b574}}.

\bibitem{RepoLedger}
\repo,
\emph{Formalization coverage and dependency ledger},
\path{docs/FORMALIZATION.md}; see also \path{docs/README.md}.
Source~09 inspected the introductory status definitions at
\href{https://github.com/VladimirReshetnikov/Surreal/blob/dcf86662b574988e75d08d515d3a2a5ba7bbc0d7/docs/FORMALIZATION.md}{\texttt{dcf86662b574}}.

\end{thebibliography}
\end{document}
